\documentclass[11pt]{article}
\usepackage[utf8]{inputenc}
\usepackage{amsmath,amssymb,amsthm,mathtools}
\usepackage[margin=2.5cm]{geometry}
\usepackage[hypertexnames=false]{hyperref}
\usepackage{enumitem}
\usepackage{algorithm,algorithmic}
\usepackage{booktabs}
\usepackage{multirow}
\usepackage{longtable}
\usepackage{float}
\usepackage{multicol}
\usepackage{graphicx}
\graphicspath{{../output/figures/mean_field/}{../output/figures/vmf/}{figure/}}
\usepackage{wrapfig}
\usepackage{pifont}
\usepackage{pgfplots}
\pgfplotsset{compat=1.18}

\newtheorem{theorem}{Theorem}[section]
\newtheorem{lemma}[theorem]{Lemma}
\newtheorem{proposition}[theorem]{Proposition}
\newtheorem{corollary}[theorem]{Corollary}
\newtheorem{definition}[theorem]{Definition}
\newtheorem{remark}[theorem]{Remark}
\newtheorem{assumption}{Assumption}
\renewcommand{\theassumption}{A\arabic{assumption}}

\DeclareMathOperator*{\argmin}{arg\,min}
\DeclareMathOperator{\softmax}{softmax}
\DeclareMathOperator{\diag}{diag}
\newcommand{\N}{\mathbb{N}}
\newcommand{\R}{\mathbb{R}}
\newcommand{\E}{\mathbb{E}}
\newcommand{\Prob}{\mathbb{P}}
\newcommand{\cL}{\mathcal{L}}
\newcommand{\cW}{\mathcal{W}}
\newcommand{\cS}{\mathcal{S}}
\newcommand{\cN}{\mathcal{N}}
\newcommand{\cB}{\mathcal{B}}
\newcommand{\cD}{\mathcal{D}}  % output-space domain ball (distinct from barrier \cB)
\newcommand{\cE}{\mathcal{E}}
\newcommand{\cY}{\mathcal{Y}}
\DeclareMathOperator{\vmf}{vMF}
\newcommand{\nperm}{\mathfrak{p}}  % neuron permutation (renamed from \pi to avoid conflict with mixing weights \pi_k)
% Named label for table rows (assumption references)
% \ref{#1} will display #2 (e.g. "A1") rather than the section number
\makeatletter
\newcommand{\namedlabel}[2]{%
  \protected@edef\@currentlabel{#2}%
  \phantomsection\label{#1}#2}
\makeatother

\title{Linear Mode Connectivity Beyond Convex Losses}
\author{}
\date{}

\begin{document}
\maketitle

\begin{abstract}
All existing theoretical guarantees for linear mode connectivity (LMC)
of neural networks require convex, Lipschitz losses.
We remove this convexity requirement: for any
\emph{$\rho$-semi-convex} loss, we prove that the interpolation
barrier of wide two-layer networks decomposes into a transfer term
(controlled by input-weight alignment) and a semi-convexity penalty
proportional to $\rho$, both vanishing as width grows.
We instantiate this general framework on two non-convex mixture-model
families.
For \emph{Mixture Density Networks} (MDNs), a log-sum-exp Hessian
decomposition yields $\rho = O(\Lambda^4/s_{\min}^4)$, giving a
barrier of $O(\sqrt{d/N} + K/N)$.
For \emph{von Mises-Fisher (vMF) mixture networks}, the
exponential-family structure of each component makes the per-component
Hessian negative semi-definite; the resulting semi-convexity constant
is a \emph{universal} $\rho_{\vmf} \le 6$---independent of all model
parameters---with no architectural constraints required.
Under a conditional head-regularity hypothesis, both rates improve
to $O(d/N + K/N)$.
We complement the theory with two-level alignment algorithms and
empirical validation on synthetic benchmarks.
\end{abstract}

%======================================================================
\section{Introduction}\label{sec:intro}
%======================================================================

Two independently trained neural networks often lie in the same loss
basin: after permuting one network's neurons, the linear interpolation
path between them stays at low loss throughout
(Figure~\ref{fig:schematics}).
This \emph{linear mode connectivity} (LMC) phenomenon---originally
demonstrated via nonlinear paths~\cite{garipov2018,draxler2018} and
strengthened to linear connectivity by
Frankle et al.~\cite{frankle2020}---has become a practical foundation
for model merging~\cite{wortsman2022soups,ilharco2023task,yadav2023ties},
federated averaging, and loss-landscape
diagnostics~\cite{entezari2022,ainsworth2023}.

\begin{wrapfigure}{r}{0.52\textwidth}
  \centering
  \vspace{-12pt}
  \includegraphics[width=0.50\textwidth]{schematics.pdf}
  \caption{Schematic of mixture-model LMC.
    After neuron permutation (red) and component permutation (blue),
    the \textcolor{orange}{origin} is transformed close to the
    \textcolor{blue!70!black}{target} in the loss landscape,
    enabling low-barrier linear interpolation.}
  \label{fig:schematics}
  \vspace{-8pt}
\end{wrapfigure}

On the theoretical side,
Ferbach et al.~\cite{ferbach2024} established the first quantitative
guarantee: for wide two-layer networks with i.i.d.\ mean-field sampled
weights and a \emph{convex, Lipschitz} loss, the interpolation barrier
vanishes as width grows.
Their proof rests on two pillars:
(i)~an optimal-transport (OT) alignment of neurons, with an
activation-matching property specific to scalar output heads, and
(ii)~convexity of the loss, which converts the alignment error directly
into a barrier bound.
Several works extend LMC to other
settings~\cite{zhan2025permutation,theus2025,kanoh2025dte,tran2025moe},
but all assume convex output losses.

Many important model families train with non-convex losses.
Mixture Density Networks
(MDNs)~\cite{bishop1994,graves2013,ha2018worldmodels,makansi2019,papamakarios2016}
model conditional densities as weighted sums of Gaussians;
von Mises--Fisher (vMF) mixture
networks~\cite{banerjee2005,gopal2014} model directional data on the
unit sphere~$S^{d_y-1}$.
Both fall outside the convex-loss framework:
their NLLs $-\log\sum_k\pi_k f_k$ are non-convex, breaking
pillar~(ii), while each neuron's output head is a vector (in
$\R^{3K}$ for MDNs, $\R^{K(1+d_y)}$ for vMF), breaking
pillar~(i).
Like Mixture-of-Experts (MoE)~\cite{tran2025moe}, these outputs are
symmetric sums over components, so establishing LMC requires a
\emph{two-level alignment}: a neuron permutation in the hidden layer
together with a component permutation in the output layer
(Figure~\ref{fig:schematics}).

Our key insight is that, for non-convex losses, the degree of
\emph{semi-convexity}~$\rho$ determines the quality of the barrier
bound.
A loss $\ell$ is $\rho$-semi-convex if $\ell + \tfrac{\rho}{2}\|\cdot\|^2$
is convex; the constant $\rho \ge 0$ captures how far the loss departs
from convexity.
We prove a general LMC theorem (Theorem~\ref{thm:general}) in which the
interpolation barrier decomposes into a transfer term (controlled by
input-weight alignment) and a semi-convexity penalty proportional
to~$\rho$, both vanishing as width grows.
Crucially, $\rho$ is \emph{loss-specific}: it depends on the analytic
structure of the negative log-likelihood, not on the network
architecture, and its magnitude dictates the tightness of the bound.

Two instantiations reveal dramatically different behaviour.
For MDNs, a log-sum-exp Hessian decomposition yields a
parameter-dependent $\rho = O(\Lambda^4/s_{\min}^4)$, where
$\Lambda$ is an output-magnitude bound and $s_{\min}$ is the scale
floor (Section~\ref{sec:output_geometry}).
For vMF mixture networks, the exponential-family structure of each
component makes the per-component Hessian negative semi-definite,
yielding a \emph{universal} $\rho_{\vmf} \le 6$---independent of all
model parameters (Section~\ref{sec:vmf}).
This contrast demonstrates that \emph{loss structure governs linear
mode connectivity}: it is the geometry of the output-space loss, not
the network width alone, that controls the barrier.

We resolve the two technical obstacles outlined above by replacing
both of Ferbach et al.'s pillars.
In place of~(i), Proposition~\ref{prop:transfer} gives a
\emph{direct output-coupling} bound for multi-dimensional heads.
In place of~(ii), the Hessian decomposition reduces the
semi-convexity constant to a manageable, loss-specific quantity.
Our contributions are as follows.
\begin{itemize}[leftmargin=*,itemsep=2pt]
  \item \textbf{General semi-convex LMC theorem}
    (Theorem~\ref{thm:general}):
    barrier $= $ transfer $+$ $\rho$-penalty, both $\to 0$ as
    $N \to \infty$; strictly generalises Ferbach et
    al.~\cite{ferbach2024} ($\rho = 0$ special case).
  \item \textbf{MDN instantiation}
    (Corollary~\ref{cor:asymptotic}):
    $\rho = O(\Lambda^4/s_{\min}^4)$, barrier
    $O(\kappa_W\sqrt{d/N} + K/N)$.
  \item \textbf{vMF instantiation}
    (Corollary~\ref{cor:vmf_barrier}):
    universal $\rho_{\vmf} \le 6$, barrier
    $O(\sqrt{d_y}\,\kappa_W\sqrt{d/N} + K/N)$.
  \item \textbf{Refined rate under head regularity}
    (Corollary~\ref{cor:struct_rate}):
    $O(\kappa_W^2 d/N + K/N)$ for both families;
    two-level alignment algorithms
    (Section~\ref{sec:algorithm}).
\end{itemize}

\noindent
Table~\ref{tab:refinement} summarises the progression of results across
the two loss families.

\begin{table}[t]
\centering\small
\caption{Barrier rates for MDN and vMF mixture applications.
  Each row within a block adds assumptions to the previous one.}
\label{tab:refinement}
\begin{tabular}{@{}lllll@{}}
\toprule
\textbf{Loss} & \textbf{Result} & \textbf{Additional assumptions}
  & \textbf{Alignment} & \textbf{Rate} \\
\midrule
\multirow{2}{*}{MDN}
  & Cor.~\ref{cor:asymptotic}
  & sub-Gaussian $W$
  & $W$-OT
  & $O(\kappa_W\sqrt{d/N} + K/N)$ \\
  & Cor.~\ref{cor:struct_rate}
  & $+$ head regularity, $C^2$ act.
  & $W$-OT
  & $O(\kappa_W^2 d/N + K/N)$ \\
\midrule
vMF
  & Cor.~\ref{cor:vmf_barrier}
  & sub-Gaussian $W$
  & $W$-OT
  & $O(\sqrt{d_y}\,\kappa_W\sqrt{d/N} + K/N)$ \\
\bottomrule
\end{tabular}
\end{table}

%======================================================================
\section{Related Work}\label{sec:related}
%======================================================================

\paragraph{Linear mode connectivity theory.}
Garipov et al.~\cite{garipov2018} and Draxler et al.~\cite{draxler2018}
demonstrated nonlinear low-loss paths between minima;
Frankle et al.~\cite{frankle2020} strengthened this to linear paths
after neuron permutation.
Entezari et al.~\cite{entezari2022} conjectured that
permutation-invariance alone explains LMC, and
Ainsworth et al.~\cite{ainsworth2023} provided practical matching
algorithms.
Ferbach et al.~\cite{ferbach2024} gave the first quantitative barrier
bound for wide mean-field networks under a convex, Lipschitz loss.
Subsequent extensions treat multi-layer
networks~\cite{zhan2025permutation}, convolutional
architectures~\cite{theus2025}, and data-dependent
transport~\cite{kanoh2025dte}, yet all retain the convexity assumption.
Among these, Zhan et al.~\cite{zhan2025permutation} achieve a
dimension-free $O(m^{-1/2})$ rate but in a teacher--student setting
with known ground-truth structure;
Theus et al.~\cite{theus2025} extend the empirical scope to
Transformers but provide no quantitative barrier bounds.
Ferbach et al.'s result therefore remains the only quantitative
barrier guarantee in the general mean-field two-layer setting
with i.i.d.\ sampled weights, making it the natural baseline
for theoretical extensions.
Our work takes this baseline and removes the convexity
assumption via semi-convexity, strictly generalising their
result (Theorem~\ref{thm:general} subsumes their bound as
the $\rho=0$ special case;
see Remark~\ref{rem:ferbach_recovery}).

\paragraph{Mixture-of-Experts LMC.}
Tran et al.~\cite{tran2025moe} study LMC for Mixture-of-Experts models,
introducing a two-level alignment (expert permutation $+$ neuron
permutation).
We share the two-level alignment structure but address a distinct
challenge: the non-convexity of the mixture-model NLL, which does not
arise in the MoE setting where each expert uses a convex loss.

\paragraph{Model merging.}
LMC underpins practical model merging methods including model
soups~\cite{wortsman2022soups}, task
arithmetic~\cite{ilharco2023task}, TIES-Merging~\cite{yadav2023ties},
and weight matching~\cite{ainsworth2023}.
Our barrier bounds provide theoretical grounding for applying these
techniques to mixture density networks.

\paragraph{Mean-field theory.}
The mean-field regime of two-layer networks---where weights are sampled
i.i.d.\ from a common law that evolves under training---has been
studied extensively~\cite{mei2018meanfield,rotskoff2022}.
De~Bortoli et al.~\cite{debortoli2020} establish quantitative
propagation-of-chaos bounds that justify Assumption~\ref{ass:mf_law}
in our setting (see Appendix~\ref{app:A4_mf}).

\paragraph{Mixture density networks.}
Bishop~\cite{bishop1994} introduced MDNs for modelling conditional
densities as Gaussian mixtures.
Graves~\cite{graves2013} applied MDNs to handwriting synthesis, and
Ha \& Schmidhuber~\cite{ha2018worldmodels} used them in world models.
vMF mixture models for directional data were studied
by Banerjee et al.~\cite{banerjee2005} and Gopal \&
Yang~\cite{gopal2014}.

%======================================================================
\section{Preliminaries}\label{sec:prelim}
%======================================================================

\subsection{Two-Layer Mixture Networks}\label{sec:general_arch}

We consider two-layer networks with $N$ hidden neurons, input
$x \in \R^d$, and $1/N$ mean-field scaling.
The backbone is
$\phi(x) = \bigl(\sigma(w_i^\top x)\bigr)_{i=1}^N$,
followed by a linear head producing raw outputs
\begin{equation}\label{eq:network_output}
  z(x) = \frac{1}{N}\sum_{i=1}^N v_i\,\sigma(w_i^\top x)
  \;\in\; \R^p,
\end{equation}
where $v_i \in \R^p$ is the head weight and
$w_i \in \R^d$ the input weight of neuron~$i$.
The output dimension $p$ and the mapping from $z$ to a conditional
density $p(y \mid x)$ are model-specific; see
Sections~\ref{sec:mdn} and~\ref{sec:vmf}.

\paragraph{Interpolation barrier.}
Let $P$ denote the data distribution over $(x,y)$.
Given two parameter vectors $\theta_A,\theta_B$ and population loss
$\cE(\theta)=\E_P[\ell(z(x;\theta),y)]$, the
\emph{interpolation barrier} is
\[
  \cB(\theta_A,\theta_B)
  = \sup_{t\in[0,1]}\Bigl\{
    \cE((1-t)\theta_A + t\theta_B)
    - (1-t)\cE(\theta_A) - t\cE(\theta_B)
  \Bigr\}.
\]
The networks are \emph{linearly mode connected} if
$\cB \le \varepsilon$ for a small $\varepsilon > 0$ after suitable
neuron permutation.

\paragraph{Component permutation symmetry.}
When the loss involves a mixture of $K$ components,
$\ell(P_\tau z, y) = \ell(z, y)$ for any component
permutation $\tau \in \cS_K$
($\cS_K$: symmetric group on $\{1,\dots,K\}$),
where $P_\tau$ is the corresponding block permutation matrix.
This gives rise to a \emph{two-level alignment}: optimising over
both a neuron permutation $\nperm \in \cS_N$ and a component
permutation $\tau \in \cS_K$
(Figure~\ref{fig:schematics}), yielding permuted head weights
$\widetilde{v}_{B,i} = P_\tau v_{B,\nperm(i)}$.

\paragraph{Neuron permutation via $W$-optimal transport.}
Throughout this paper the neuron permutation $\nperm \in \cS_N$ is
defined as the \textbf{$W$-optimal-transport ($W$-OT) coupling}---the
linear assignment problem on \emph{input weights} $w_i \in \R^d$ only:
\begin{equation}\label{eq:OT_perm}
  \nperm \;\coloneqq\; \argmin_{q \in \cS_N}
  \frac{1}{N}\sum_{i=1}^N \|w_{A,i} - w_{B,q(i)}\|^2.
\end{equation}
This $\nperm$ minimises the input-weight mismatch
$B_N = \frac{1}{N}\sum_i\|w_{A,i}-w_{B,\nperm(i)}\|^2$.
As Proposition~\ref{prop:transfer} shows, only $B_N$---not the
head-weight mismatch $A_N$---appears in the interpolation barrier,
so aligning on $W$ alone is sufficient for the baseline barrier bound.
Component correspondence is handled separately by $\tau$ at
$O(K^3)$ cost (Section~\ref{sec:algorithm}).

\subsection{MDN Architecture and Loss}\label{sec:mdn}

A Mixture Density Network (MDN) with $K$ Gaussian components
models $y \in \R$ conditioned on $x \in \R^d$.
The output dimension is $p = 3K$, with
\begin{equation}\label{eq:mdn_output}
  z(x) = \frac{1}{N}\sum_{i=1}^N v_i \sigma(w_i^\top x)
  \in \R^{3K},
\end{equation}
partitioned as $z = (z_\pi, z_\mu, z_s) \in \R^K \times \R^K \times \R^K$.
The conditional density is
\[
  p(y \mid x) = \sum_{k=1}^K \pi_k(x)\,
  \cN(y;\,\mu_k(x),\,s_k(x)^2),
\]
where $\pi_k = \softmax(z_\pi)_k$ (mixture weights), $\mu_k = z_{\mu,k}$,
and $s_k = e^{z_{s,k}}$ (log-standard-deviation parametrisation).
We distinguish $\pi_k$ (network-output mixture weights) from the
EM responsibilities $\gamma_k = \pi_k\varphi_k / \sum_j \pi_j\varphi_j$
(posterior component probabilities given $(x,y)$), which appear in
the Hessian analysis of Section~\ref{sec:output_geometry}.
All gradient and Hessian computations are with respect to the raw
outputs~$z$. The MDN loss is
\begin{equation}\label{eq:mdn_loss}
  \cL_{\mathrm{MDN}}(z, y)
  = -\log\!\Biggl(\sum_{k=1}^K \pi_k\,
    \frac{1}{\sqrt{2\pi}\,s_k}\,
    \exp\!\Bigl(-\frac{(y-\mu_k)^2}{2s_k^2}\Bigr)\Biggr).
\end{equation}

The completeness of the component permutation symmetry for MDNs
is analysed in Section~\ref{sec:symmetry};
see also Tran et al.~\cite{tran2025moe} for the analogous
symmetry in MoE architectures.

\subsection{Standing Assumptions}\label{sec:assumptions}

The following assumptions extend the mean-field framework
of Ferbach et al.~\cite{ferbach2024}, replacing
their convex loss assumption with our semi-convexity analysis.
Assumptions~\ref{ass:arch}--\ref{ass:bdd_resp} are loss-specific;
\ref{ass:activ}--\ref{ass:subgauss} apply to all losses.
Each theorem statement lists the assumptions it requires.

\begin{assumption}[Architecture]\label{ass:arch}
The two networks $A$ and $B$ are two-layer networks parameterised
as in~\eqref{eq:network_output} with the $1/N$ mean-field scaling.
The per-neuron parameter space is $\R^{p+d}$, where $v_i \in \R^p$
is the head weight vector and $w_i \in \R^d$ is the input weight.
\emph{MDN-specific:}
the log-standard-deviation outputs are clamped from below:
$z_{s,k} \mapsto \max(z_{s,k},\,\underline{z}_s)$ for a fixed
$\underline{z}_s \in \R$, ensuring
$s_k = e^{z_{s,k}} \ge s_{\min} \coloneqq e^{\underline{z}_s} > 0$ for
\emph{all} network outputs, including those at interpolated parameters.
This is standard practice in MDN
implementations~(e.g.\ \texttt{torch.clamp}).
For vMF mixtures, no such clamp is needed.
\end{assumption}

\begin{assumption}[Bounded response]\label{ass:bdd_resp}
The response variable is bounded: $|y| \le Y_{\max}$ $P$-a.s.\
for some finite $Y_{\max} > 0$ (MDN), or
$y \in S^{d_y-1}$ so that $\|y\| = 1$ (vMF).
\end{assumption}

\begin{assumption}[Activation regularity]\label{ass:activ}
The activation $\sigma:\R\to\R$ satisfies:
\begin{enumerate}[label=(\roman*),leftmargin=*,itemsep=0pt]
  \item $|\sigma(u)| \le S_\sigma$ for all $u\in\R$, for some $S_\sigma \in \R_+$
    (boundedness);
  \item $|\sigma(u) - \sigma(v)| \le L_\sigma|u-v|$ for all $u,v\in\R$,
    for some $L_\sigma \in \R_+$ (Lipschitz continuity).
\end{enumerate}
When stated explicitly, we additionally require:
\begin{enumerate}[label=(\roman*),leftmargin=*,itemsep=0pt,start=3]
  \item $\sigma \in C^2(\R)$ with
    $\|\sigma''\|_\infty \le L_{\sigma'} < \infty$
    (second-order smoothness).
\end{enumerate}
For the standard sigmoid $\sigma(u) = 1/(1+e^{-u})$, conditions (i)--(iii) hold with
$S_\sigma = 1$, $L_\sigma = 1/4$, and
$L_{\sigma'} = 1/(6\sqrt{3}) \approx 0.096$.
\end{assumption}

\begin{assumption}[Finite input moments]\label{ass:input_mom}
The input distribution satisfies $\E_P[\|x\|^2] \le M_x^2$
for some finite constant $M_x > 0$.
Write $\Sigma_x \coloneqq \E_P[xx^\top]$ for the input covariance
(so $\mathrm{tr}(\Sigma_x) = M_x^2$).
When stated explicitly, we additionally require a finite fourth moment:
$\E_P[\|x\|^4] \le M_{x,4}^4$.
\end{assumption}

\begin{assumption}[Mean-field law from finite-time training]\label{ass:mf_law}
Both networks are trained by (mean-field) gradient flow for time
$T < \infty$, starting from i.i.d.\ initialisation
$\theta_{X,i}(0) \sim \mathfrak{m}_0$ with
$\E_{\mathfrak{m}_0}[\|\theta\|^4] < \infty$
(e.g.\ Gaussian, He, or Xavier initialisation).
After training, the neuron parameters $\theta_{X,i} = (v_{X,i}, w_{X,i})$,
$i=1,\ldots,N$, of each network $X \in \{A, B\}$ are i.i.d.\ draws
from the time-$T$ law $\mathfrak{m}_* \coloneqq \mathfrak{m}_T$ on
$\R^{p+d}$.
\end{assumption}

The i.i.d.\ structure is justified via propagation of
chaos~\cite{mei2018meanfield,rotskoff2022,debortoli2020} at
$O(1/\sqrt{N})$ error; see Appendix~\ref{app:A4_mf}.
Empirical validation (neuron density overlays and predictive
convergence across seeds) is provided in
Appendix~\ref{app:a5_validation}.

The following moment quantities of $\mathfrak{m}_*$ appear in the
barrier bounds:
\begin{equation}\label{eq:moment_defs}
  m_{2,*} \coloneqq \int\!\|V\|^2\,d\mathfrak{m}_*,
  \qquad
  m_{4,*} \coloneqq \max_{j=1,\ldots,p}
    \E_{\mathfrak{m}_*}[(V)_j^4].
\end{equation}
In the finite-time mean-field setting, Proposition~\ref{prop:a6_auto}
(Appendix~\ref{app:a6_proposition}) gives a sufficient condition
ensuring that both are finite.
The per-coordinate fourth moment $m_{4,*}$ enables sharper
concentration of the head moment (Theorem~\ref{thm:hp}).

\begin{assumption}[Sub-Gaussian input weights]\label{ass:subgauss}
The $W$-marginal of $\mathfrak{m}_*$ is
$\kappa_W^2$-sub-Gaussian:
for all $\lambda \in \R^d$,
\[
  \E_{\mathfrak{m}_*}\!\bigl[\exp\bigl(\langle\lambda, W - \E[W]\rangle\bigr)\bigr]
  \;\le\; \exp\!\Bigl(\frac{\kappa_W^2\|\lambda\|^2}{2}\Bigr).
\]
\end{assumption}

This holds at Gaussian initialisation; preservation under training
is not guaranteed in general (the flow map is locally, not globally,
Lipschitz; see Appendix~\ref{app:a7_validation}).
The $W$-marginal sub-Gaussian condition (Assumption~\ref{ass:subgauss})
is empirically supported across all mean-field experiments
(Appendix~\ref{app:a7_validation}).
Sub-Gaussianity yields the parametric $\cW_2$ convergence rate
$\eta_N = O(\kappa_W\sqrt{d/N})$, with $d$ entering only linearly,
via a covering-number argument~\cite{lei2020,weed2019}.

%======================================================================
\section{Theoretical Results}\label{sec:loss_geometry}
%======================================================================

We first state a general barrier theorem that applies to any
$\rho$-semi-convex loss, then specialise to MDNs and vMF mixtures.
All proofs are in Appendix~\ref{app:proofs}.

\subsection{General Framework for Semi-Convex Losses}
\label{sec:general_framework}

A loss $\ell:\R^p \times \cY \to \R$ is \emph{$\rho$-semi-convex}
on a domain $\cD \subset \R^p$ if
$\nabla^2_z \ell(z,y) \succeq -\rho I$ for all $z \in \cD$,
$y \in \cY$; it is \emph{$C_L^\infty$-Lipschitz} on $\cD$ if
$|\ell(z_1,y)-\ell(z_2,y)| \le C_L^\infty \|z_1-z_2\|_\infty$.
Consider a two-layer network
$z(x)=\frac{1}{N}\sum_{i=1}^N v_i\,\sigma(w_i^\top x)\in\R^p$
with head weights $v_i\in\R^p$, input weights $w_i\in\R^d$, and
$1/N$ mean-field scaling,
under Assumptions~\ref{ass:activ} and~\ref{ass:input_mom}.
Define $M_V^\infty$ as in~\eqref{eq:head_moment_inf}
and set $R=S_\sigma M_V^\infty$.

\begin{theorem}[General LMC barrier for semi-convex losses]
\label{thm:general}
Let $\ell$ be $\rho$-semi-convex and $C_L^\infty$-Lipschitz on
$\cD_R = \{z:\|z\|_\infty \le R\}$.
Fix any neuron permutation $\nperm \in \cS_N$.
Then for all $t\in[0,1]$:
\begin{equation}\label{eq:general_bound}
  \boxed{
  \cB \;\le\;
  \underbrace{\tfrac{1}{2}\,C_L^\infty\,L_\sigma\,M_V^\infty
    M_x\sqrt{B_N}}_{\text{transfer}}
  \;+\;
  \underbrace{\tfrac{\rho}{8}\,\Delta_N}_{\text{semi-convexity}}
  }
\end{equation}
where $B_N = \frac{1}{N}\sum_i\|\Delta w_i\|^2$ and
$\Delta_N = \E_P[\|z_A(x)-z_B'(x)\|^2]$
(with $z_B'$ the permuted network-$B$ output).
\end{theorem}

The proof (Appendix~\ref{sec:proof_main}) combines two
loss-independent steps: the transfer bound
(Proposition~\ref{prop:transfer}) controls the
parameter-to-output gap, and $\rho$-semi-convexity replaces
convexity in the output-interpolation step.
The high-probability bound (Theorem~\ref{thm:hp}) and asymptotic
rate (Corollary~\ref{cor:asymptotic}) carry over with the
loss-specific $\rho$ and $C_L^\infty$.

\begin{remark}[Subsumption of Ferbach et al.]\label{rem:ferbach_recovery}
Setting $\rho=0$ (convex loss) eliminates the semi-convexity
penalty, reducing~\eqref{eq:general_bound} to
$\cB \le \tfrac{1}{2}\,C_L^\infty L_\sigma M_V^\infty M_x\sqrt{B_N}$.
This strictly generalises the barrier guarantee
of Ferbach et al.~\cite{ferbach2024} (Corollary~3.2)
in two respects:
(i)~it provides an explicit rate
$\cB=O(\kappa_W\sqrt{d/N})$
(Corollary~\ref{cor:asymptotic}),
whereas Ferbach et al.\ establish qualitative convergence
($\forall\,\mathrm{err}>0,\;\exists\,N_{\min}$), and
(ii)~Proposition~\ref{prop:transfer} replaces their
activation-matching property (which requires scalar output heads)
with a direct output-coupling bound valid for
multi-dimensional heads ($p\ge 1$).
For $\rho > 0$, our framework further extends
to non-convex losses that lie entirely outside the scope of
prior work.
\end{remark}
To apply this framework, one must compute $\rho$ and $C_L^\infty$
for the loss at hand.
For mixture-model NLLs of the form $\ell = -\log\sum_k e^{u_k}$,
the log-sum-exp Hessian decomposition~\eqref{eq:hess_decomp}
reduces this to bounding per-component quantities
$\|\nabla u_k\|$ and $\|\nabla^2 u_k\|_{\mathrm{op}}$.

\subsection{MDN Output-Space Geometry}\label{sec:output_geometry}

We now compute $\rho$ and $C_L^\infty$ for the MDN loss to
instantiate Theorem~\ref{thm:general}.
Let $\cD_R = \{z \in \R^{3K} : \|z\|_\infty \le R\}$ denote the
output-space domain; our goal is to obtain $\rho$ with a manageable
dependence on~$R$.
A direct Hessian bound via $\nabla^2(-\log p)$ requires a lower bound
on $p$, yielding a doubly-exponential $\rho$ (vacuous in practice).
We avoid $p_{\min}$ entirely by exploiting the log-sum-exp structure.
Writing $u_k = \log(\pi_k\varphi_k)$ with
$\varphi_k \coloneqq \cN(y;\mu_k,s_k^2)$,
we have $\cL_{\mathrm{MDN}} = -\log\sum_k e^{u_k}$, and the
standard log-sum-exp Hessian identity~\cite{boyd2004} gives
\begin{equation}\label{eq:hess_decomp}
  \nabla^2 \cL_{\mathrm{MDN}}
  = -\sum_k \gamma_k \nabla^2 u_k - \mathrm{Cov}_\gamma(\nabla u_k),
\end{equation}
where $\gamma_k = e^{u_k}/\sum_j e^{u_j}$ are the EM responsibilities.
Since $\gamma_k \in [0,1]$ absorb the $1/p$ denominator, both the
Lipschitz ($C_L$) and semi-convexity ($\rho$) bounds reduce to
controlling $\|\nabla u_k\|$ and $\|\nabla^2 u_k\|_{\mathrm{op}}$ on
$\cD_R$---finite quantities with no $p_{\min}$ factor.

The \emph{scale floor} $s_{\min} = e^{\underline{z}_s}$ is the
architectural parameter from Assumption~\ref{ass:arch}, guaranteeing
$s_k \ge s_{\min} > 0$ for all network outputs.
Because the clamp is applied coordinate-wise \emph{after} the linear
head, the clamped output of any network---including one with
interpolated parameters---satisfies $s_k \ge s_{\min}$.
Set $\Lambda = Y_{\max} + R$.

\begin{proposition}[Loss regularity]\label{prop:loss_regularity}
On $\cD_R$ with $|y| \le Y_{\max}$ and $\min_k e^{z_{s,k}} \ge s_{\min}$:
\begin{equation}\label{eq:refined_bounds}
  C_L^\infty \;\le\; \hat c_0\,\frac{\Lambda^2}{s_{\min}^2},
  \qquad
  C_L \;\le\; \hat c_1\,\frac{\Lambda^2}{s_{\min}^2},
  \qquad
  \rho \;\le\; \frac{1}{s_{\min}^2}
  + 2\hat c_2\,\frac{\Lambda^4}{s_{\min}^4},
\end{equation}
where $\hat c_0, \hat c_1, \hat c_2 > 0$ are universal constants.
The simplified forms above assume $\Lambda \ge 1$;
for general $\Lambda$ the full expressions
from the proof~\eqref{eq:grad_sq_bound}, \eqref{eq:rho_bound} hold
with additive universal constants
(e.g.\ $C_L^\infty \le 3 + \hat c_0' \Lambda^2/s_{\min}^2$).
In all downstream applications $\Lambda = Y_{\max} + R \ge Y_{\max}$,
which is $\ge 1$ in typical regression settings.
\end{proposition}

The $\ell_\infty$-Lipschitz constant $C_L^\infty$ is the one used in
the transfer bound (Proposition~\ref{prop:transfer}).

While the bound on $\rho$ is stated uniformly in $K$,
the Hessian decomposition~\eqref{eq:hess_decomp} reveals a
qualitative gap between $K=1$ and $K \ge 2$.
For $K=1$ the EM weight $\gamma_1 \equiv 1$, so the
covariance term $\mathrm{Cov}_\gamma(\nabla u_k)$ vanishes
identically; the semi-convexity is then controlled by
$-\nabla^2 u_1$ alone, giving $\rho_{K=1} \le 1/s_{\min}^2$
independently of $\Lambda$.
For $K \ge 2$ the covariance term is non-zero---the competition
factor $\gamma_k(1-\gamma_k)$ reaches $1/4$ at
$\gamma_k = 1/2$---and drives the $O(\Lambda^4/s_{\min}^4)$
scaling in~\eqref{eq:refined_bounds}.

\subsection{Main Theorems}\label{sec:main_theorem}

The output-space bounds above apply to the linear interpolation of
outputs $\hat z_t = (1-t)z_A + tP_\tau z_B$, but the barrier involves
the interpolated \emph{network} $z_{M_t}$.
Define the \emph{per-coordinate head moment}
\begin{equation}\label{eq:head_moment_inf}
  M_V^\infty \;\coloneqq\;
  \max\!\Bigl\{
    \max_{j=1,\ldots,p}\!\Bigl(\frac{1}{N}\sum_{i=1}^N (v_{A,i})_j^2\Bigr)^{\!1/2},\;
    \max_{j=1,\ldots,p}\!\Bigl(\frac{1}{N}\sum_{i=1}^N (P_\tau v_{B,\nperm(i)})_j^2\Bigr)^{\!1/2}
  \Bigr\}
\end{equation}
and set $R = S_\sigma M_V^\infty$.
A Cauchy--Schwarz argument shows that all network outputs
$z_A$, $P_\tau z_B$, $z_{M_t}$, $\hat z_t$ lie in $\cD_R$
(Lemma~\ref{lem:auto_path_domain} in the appendix).

The following proposition quantifies the transfer gap
$z_{M_t} - \hat z_t$ via the input-weight mismatch
$B_N = \frac{1}{N}\sum_{i} \|\Delta w_i\|^2$,
where $\Delta w_i = w_{A,i} - w_{B,\nperm(i)}$.

\begin{proposition}[Direct output coupling]\label{prop:transfer}
Under Assumptions~\ref{ass:activ} and~\ref{ass:input_mom},
the gap between the interpolated network output $z_{M_t}$ and the
output interpolation $\hat z_t$ satisfies
\begin{equation}\label{eq:transfer}
  \sup_{t\in[0,1]}
  \bigl(\E_P[\|z_{M_t}(x) - \hat z_t(x)\|_\infty^2]\bigr)^{1/2}
  \;\le\; \tfrac{1}{2}\,L_\sigma M_V^\infty M_x \sqrt{B_N},
\end{equation}
where $M_V^\infty$ is the per-coordinate head moment~\eqref{eq:head_moment_inf}
and $B_N = \frac{1}{N}\sum_i \|\Delta w_i\|^2$ is the input-weight mismatch.
\end{proposition}

The bound depends only on the input-weight mismatch $B_N$,
not on the head-weight mismatch~$A_N$:
the output-coupling gap decomposes into activation differences
that depend on $\Delta w_i$ but not on head-weight differences;
the head weights enter only through their amplitudes $M_V^\infty$.
Combined with the loss-specific $C_L^\infty$,
the transfer term is independent of $K$
(the semi-convexity term carries the $K$-dependence
via $\Delta_N$; see Theorem~\ref{thm:general}).

The endpoint output mismatch
$\Delta_N = \E_P[\|z_A(x) - P_\tau z_B(x)\|^2]$ is bounded via
a parameter-space decomposition
(Lemma~\ref{lem:delta_control} in the appendix).

\noindent
The \emph{aligned barrier} $\cB(\theta_A,\theta_B;\nperm,\tau)$ is
defined by permuting network~$B$'s weights to
$(P_\tau v_{B,\nperm(i)}, w_{B,\nperm(i)})$ before interpolation.

\begin{theorem}[MDN barrier---deterministic]\label{thm:main}
Theorem~\ref{thm:general} applied to $\ell = \cL_{\mathrm{MDN}}$
under Assumptions~\ref{ass:arch}--\ref{ass:input_mom}, with
$p = 3K$,
$C_L^\infty$ and $\rho$ from
Proposition~\ref{prop:loss_regularity}
($C_L^\infty = O(\Lambda^2/s_{\min}^2)$,
$\rho = O(\Lambda^4/s_{\min}^4)$,
$\Lambda = Y_{\max} + S_\sigma M_V^\infty$),
yields: for any $(\nperm,\tau) \in \cS_N \times \cS_K$
and all $t \in [0,1]$,
\begin{equation}\label{eq:main_bound}
  \begin{aligned}
  \E_P[\cL_{\mathrm{MDN}}(z_{M_t}, y)]
  &\le (1-t)\,\E_P[\cL_{\mathrm{MDN}}(z_A, y)]
  + t\,\E_P[\cL_{\mathrm{MDN}}(z_B, y)]\\
  &\quad + \tfrac{1}{2}\,C_L^\infty L_\sigma M_V^\infty M_x \sqrt{B_N}
  + \tfrac{\rho}{2}\,t(1-t)\,\Delta_N.
  \end{aligned}
\end{equation}
\end{theorem}

This is deterministic and loss-independent in structure: the
MDN-specific content enters only through $C_L^\infty$ and $\rho$.
The transfer term is $K$-independent;
$K$ enters only through the semi-convexity term via~$\Delta_N$.

Taking $\sup_{t\in[0,1]}$ and using $\sup_t t(1-t) = \tfrac{1}{4}$, the barrier form is
\begin{equation}\label{eq:barrier_form}
  \cB \;\le\; \tfrac{1}{2}\,C_L^\infty L_\sigma M_V^\infty M_x \sqrt{B_N}
  \;+\; \frac{\rho}{8}\,\Delta_N.
\end{equation}

To obtain a high-probability guarantee, we invoke the mean-field model:

\begin{theorem}[LMC barrier---high-probability existence]\label{thm:hp}
Let $\ell$ be $\rho$-semi-convex and $C_L^\infty$-Lipschitz on
$\cD_R$ with $R = S_\sigma M_V^\infty$ (so that
Theorem~\ref{thm:general} applies).
Under Assumptions~\ref{ass:arch}, \ref{ass:activ},
and~\ref{ass:input_mom}, additionally suppose
Assumption~\ref{ass:mf_law}: both networks $A$ and $B$ are
drawn i.i.d.\ from $\mathfrak{m}_* = \mathfrak{m}_T$.
Let $m_{2,*}, m_{4,*}$ be the moment quantities
defined in~\eqref{eq:moment_defs}, and assume they are finite.
(Proposition~\ref{prop:a6_auto} gives a sufficient criterion
for this in the finite-time mean-field setting.)
Let $\delta \in (0,1)$ and let $\eta_N$ be any quantity satisfying
$\Prob\!\bigl(\cW_2(\hat\nu_W^X, \mathfrak{m}_{*,W}) > \eta_N\bigr)
\le \delta/6$ for each network $X \in \{A,B\}$.
Then with probability at least $1-\delta$
over the draws of $A$ and $B$, the $W$-OT neuron permutation $\nperm$
from~\eqref{eq:OT_perm} together with $\tau = \mathrm{id}$
yield a barrier bound~\eqref{eq:general_bound} with
\begin{equation}\label{eq:head_moment}
  (M_V^\infty)^2 \le m_{\infty,*}^2 + \varepsilon_N(\delta), \quad
  B_N \le 4\eta_N^2, \quad
  \Delta_N \le \frac{6S_\sigma^2 m_{2,*}}{N\delta},
\end{equation}
where
$\varepsilon_N(\delta) \coloneqq \sqrt{6p\,m_{4,*}/(N\delta)}$.
\end{theorem}

The Chebyshev correction $\varepsilon_N(\delta)$ vanishes as $N\to\infty$.
In particular, as $N \to \infty$ under the mean-field model,
the barrier $\cB$ converges to zero (Corollary~\ref{cor:asymptotic}).

We now derive the asymptotic rate.\label{sec:asymptotic_rate}
By the SLLN, as $N\to\infty$ the per-coordinate head moment
converges: for each coordinate~$j$,
$(1/N)\sum_i (v_i)_j^2 \to \E_{\mathfrak{m}_*}[(V)_j^2]$.
Define the \emph{per-coordinate limiting moment}
$m_{\infty,*}^2 \coloneqq \max_j \E_{\mathfrak{m}_*}[(V)_j^2]$.
Then $M_V^\infty \to m_{\infty,*}$ almost surely,
$R \to S_\sigma m_{\infty,*}$, and the loss constants
$C_L^\infty$, $\rho$ stabilise at finite values.
Note that $m_{\infty,*}^2 = O(1)$ under natural scaling
(each coordinate has bounded second moment),
without requiring any head-normalisation assumption.
Since $m_{2,*} = \sum_j \E[(V)_j^2] \le p \cdot m_{\infty,*}^2$,
the semi-convexity term contributes
$\rho \cdot \Delta_N = O(m_{2,*}/N) = O(p/N)$.

Substituting the probabilistic bounds from Theorem~\ref{thm:hp} into the barrier form~\eqref{eq:barrier_form} yields the asymptotic rate:

\begin{corollary}[Asymptotic barrier under $W$-OT alignment]
\label{cor:asymptotic}
Under the hypotheses of Theorem~\ref{thm:hp},
additionally suppose Assumption~\ref{ass:subgauss}.
Let $\delta \in (0,1)$.  Then the $W$-OT alignment $\nperm$
of~\eqref{eq:OT_perm} yields, with probability at least $1-\delta$:
\begin{equation}\label{eq:barrier_bound}
  \boxed{
  \cB \;\le\;
  \underbrace{C_{\mathrm{tr}}(\delta,N)\,\kappa_W\sqrt{\frac{d}{N}}}_{\substack{\text{transfer}\\[2pt]
    \scriptscriptstyle O(\sqrt{d/N})}}
  \;+\;
  \underbrace{C_{\mathrm{sc}}(\delta,N)\,\frac{m_{2,*}}{N}}_{\substack{\text{semi-convexity}\\[2pt]
    \scriptscriptstyle O(K/N)}}
  \;\;\xrightarrow{\;N\to\infty\;}\;\; 0
  }
\end{equation}
where
$\hat M_V(\delta,N) \coloneqq (m_{\infty,*}^2 + \varepsilon_N(\delta))^{1/2}$,
$R(\delta,N) = S_\sigma \hat M_V(\delta,N)$,
and $C_L^\infty$, $\rho$ are the loss-specific regularity constants
evaluated at this~$R$
(e.g.\ Proposition~\ref{prop:loss_regularity} for MDN,
Proposition~\ref{prop:vmf_regularity} for vMF).
The prefactors are
\begin{equation}\label{eq:prefactors}
  C_{\mathrm{tr}}(\delta,N)
  \;=\; C_L^\infty \cdot L_\sigma\,\hat M_V(\delta,N)
        \cdot M_x,
  \qquad
  C_{\mathrm{sc}}(\delta,N)
  \;=\; \frac{3\,\rho\,S_\sigma^2}{4\,\delta}.
\end{equation}
Since $\varepsilon_N(\delta) = O(\sqrt{p/(N\delta)})$,
$\hat M_V(\delta,N) \to m_{\infty,*}$, and both prefactors converge
to finite SLLN limits $C_{\mathrm{tr}}^*$, $C_{\mathrm{sc}}^*$.
The barrier simplifies to
\begin{equation}\label{eq:barrier_rate_sg}
  \cB \;=\; O\!\left(\kappa_W\sqrt{\frac{d}{N}} + \frac{p}{N}\right),
\end{equation}
where $p/N = O(K/N)$ for both MDN ($p=3K$) and vMF ($p=K(1\!+\!d_y)$).
To achieve a target barrier $\cB \le \xi$ ($\xi > 0$),
each term must be at most $\xi/2$, giving
\begin{equation}\label{eq:width_req}
  N \;\ge\;
  \max\!\left\{
    \frac{4\,(C_{\mathrm{tr}}^*)^2 \kappa_W^2\,d}{\xi^2},\;
    \frac{2\,C_{\mathrm{sc}}^*\,m_{2,*}}{\xi}
  \right\}.
\end{equation}
The first (transfer) entry scales as $\xi^{-2}$ and dominates
for small $\xi$ or large $d$;
the second (semi-convexity) entry scales as $\xi^{-1}$ but carries
the loss-specific constant $C_{\mathrm{sc}}^*$ and can dominate at moderate~$\xi$.
In either regime the required width is polynomial in~$d$
with no exponential dependence.
\end{corollary}

\paragraph{MDN prefactor convergence.}
For MDN, the loss-specific constants depend on $R$ through
$\Lambda = Y_{\max} + R$ as
$C_L^\infty = O(\Lambda^2/s_{\min}^2)$ and
$\rho = O(\Lambda^4/s_{\min}^4)$
(Proposition~\ref{prop:loss_regularity}).
Taylor-expanding around $\hat M_V = m_{\infty,*}$:
\begin{equation}\label{eq:Ctr_rate}
  C_{\mathrm{tr}}(\delta,N)
  \;=\; C_{\mathrm{tr}}^*
  + O\!\left(\sqrt{\frac{p}{N\delta}}\right),
  \qquad
  C_{\mathrm{sc}}(\delta,N)
  \;=\; C_{\mathrm{sc}}^*
  + O\!\left(\sqrt{\frac{p}{N\delta}}\right),
\end{equation}
where $\Lambda^* = Y_{\max} + S_\sigma m_{\infty,*}$ and
the SLLN limits are
\begin{equation}\label{eq:prefactor_limits}
  C_{\mathrm{tr}}^*
  \;=\; \hat c_0\,\frac{(\Lambda^*)^2}{s_{\min}^2}
        \cdot L_\sigma\,m_{\infty,*}
        \cdot M_x,
  \qquad
  C_{\mathrm{sc}}^*
  \;=\; \frac{3\,S_\sigma^2}{4\,\delta}
        \!\left(\frac{1}{s_{\min}^2}
        + 2\hat c_2\,\frac{(\Lambda^*)^4}{s_{\min}^4}\right).
\end{equation}
Both corrections are lower-order, leaving the asymptotic
rate~\eqref{eq:barrier_rate_sg} unchanged.
For vMF, $C_L^\infty \le 2 + 2\sqrt{d_y}$ and $\rho \le 6$
are independent of~$R$ (Proposition~\ref{prop:vmf_regularity}),
so the prefactors stabilise immediately with no Taylor correction
needed.

\subsection{Improved Rate via Head Regularity}
\label{sec:head_regularity}

The baseline rate $O(\kappa_W\sqrt{d/N})$ in the transfer term
leaves a gap with the empirical decay $O(N^{-1.3})$ observed for
$K=1$ (Appendix~\ref{app:nonvacuous_detail}).
We now show that this gap can be closed under
a structural hypothesis on the trained law:
even with $W$-only alignment, the transfer rate improves to
$O(d/N)$ when the gauge-fixed head $\tilde V$ is approximately
a Lipschitz function of the hidden weight~$W$.

\paragraph{Head regularity hypothesis.}
Define the \emph{centred head weights}
\begin{equation}\label{eq:centred_head}
  \tilde v_{i,\pi_k}
  \;\coloneqq\;
  (v_i)_{\pi_k} - \frac{1}{K}\sum_{k'=1}^K (v_i)_{\pi_{k'}},
  \qquad
  \tilde v_i
  \;\coloneqq\;
  \bigl(\tilde v_{i,\pi},\; (v_i)_{K+1:p}\bigr)
  \;\in\; \R^{p},
\end{equation}
where the $\pi$-block (mixing-weight coordinates) is projected onto the
hyperplane $\{u \in \R^K : \mathbf{1}^\top u = 0\}$ and the remaining
$p-K$ component-parameter coordinates are left unchanged.
The common-mode component is barrier-neutral
(Proposition~\ref{prop:barrier_neutrality}).
Fix canonical gauge
($\frac{1}{K}\sum_k (v_{X,i})_{\pi_k} = 0$ for all $X,i$) and suppose
that for some map $g:\R^d \to \R^{p}$ and residuals $r_{X,i}$,
\begin{equation}\label{eq:head_reg_decomp}
  \tilde v_{X,i} = g(w_{X,i}) + r_{X,i},
  \qquad X \in \{A,B\},
\end{equation}
with $g$ satisfying the $\ell_2 \to \ell_\infty$ Lipschitz condition
\begin{equation}\label{eq:head_reg_lip}
  \|g(w) - g(w')\|_\infty \le L_g \|w - w'\|
  \qquad \forall\, w,w' \in \R^d.
\end{equation}
For a given alignment $(\nperm,\tau)$ define the residual mismatch
\begin{equation}\label{eq:head_residual_mismatch}
  \mathcal R_N(\nperm,\tau)
  \;\coloneqq\;
  \max_{j=1,\ldots,p}
  \Bigl(
    \frac{1}{N}\sum_{i=1}^N
    \bigl(r_{A,i} - P_\tau r_{B,\nperm(i)}\bigr)_j^2
  \Bigr)^{1/2}.
\end{equation}
This is a structural hypothesis on the trained law
($\tilde V$ is approximately a regular function of~$W$),
not an assumption on the optimisation algorithm.
Intuitively, neurons with similar input directions learn similar
output heads under mean-field dynamics; the residual $r_{X,i}$
captures the remaining stochasticity.
We probe this hypothesis empirically in
Section~\ref{sec:experiments}.

\begin{remark}[Structural support for $K\!=\!1$]
\label{rmk:head_reg_k1}
For $K=1$, the mean-field head drift
$b_v(w,m_s) = -\E_P[\sigma(w^\top x)\,\nabla_z\cL(z_{m_s}(x),y)]$
depends on the particle only through~$w$, not~$v$.
The conditional mean $g(w) \coloneqq \E[V_T \!\mid\! W_T = w]$
is thus Lipschitz under mild non-degeneracy conditions on the
flow map, providing structural evidence for the
decomposition~\eqref{eq:head_reg_decomp}.
However, the residual $r = V_T - g(W_T)$ captures the initial
head variance $\mathrm{Var}(V_0)$, which does not vanish under the
mean-field law; the stronger condition
$\mathcal{R}_N = O(\sqrt{B_N})$ therefore remains a hypothesis.
\end{remark}

\begin{proposition}[Fast transfer from conditional head regularity]
\label{prop:transfer_struct}
Under Assumptions~\ref{ass:arch}, \ref{ass:activ}\,(i)--(iii),
and~\ref{ass:input_mom} (including the fourth-moment condition
$\E_P[\|x\|^4] \le M_{x,4}^4$),
suppose additionally that the
gauge-fixed heads admit the decomposition~\eqref{eq:head_reg_decomp}
with~\eqref{eq:head_reg_lip}. Then for any alignment $(\nperm,\tau)$,
the per-coordinate gauge-fixed head mismatch
$M_{\Delta \tilde V}^\infty
\coloneqq
\max_{j} (\frac{1}{N}\sum_{i}
(\tilde v_{A,i} - P_\tau \tilde v_{B,\nperm(i)})_j^2)^{1/2}$
satisfies
\begin{equation}\label{eq:head_mismatch_struct}
  M_{\Delta \tilde V}^\infty
  \;\le\;
  L_g \sqrt{B_N} + \mathcal R_N(\nperm,\tau).
\end{equation}
Hence, if $\mathcal R_N(\nperm,\tau) = O(\sqrt{B_N})$
and $Q_N \coloneqq \frac{1}{N}\sum_{i} \E_P[(\Delta w_i^\top x)^4]
= O(B_N^2)$ along a width sequence, then the transfer
term in Theorem~\ref{thm:general} is $O(B_N) = O(d/N)$
under $W$-only alignment.
\end{proposition}

Proposition~\ref{prop:transfer_struct} identifies a concrete mechanism
by which the empirical barrier can decay faster than the baseline
$O(N^{-1/2})$ transfer law even without explicit head matching:
the conditional head map $W \mapsto \tilde V$ need only be regular
up to a residual whose aligned mismatch scales like $\sqrt{B_N}$.
This provides a target for diagnostics:
measure whether $M_{\Delta \tilde V}^\infty$ already tracks
$\sqrt{B_N}$ under $W$-only alignment
(Section~\ref{sec:experiments}).

Substituting into the barrier bound yields:

\begin{corollary}[Conditional fast barrier under $W$-OT alignment]
\label{cor:struct_rate}
Under the hypotheses of Theorem~\ref{thm:hp},
additionally suppose Assumption~\ref{ass:activ}\,(iii),
the fourth-moment condition in Assumption~\ref{ass:input_mom},
Assumption~\ref{ass:subgauss},
and the structural hypothesis of
Proposition~\ref{prop:transfer_struct}.
Let $\delta \in (0,1)$ and let $(\nperm,\tau)$ be the aligned pair from
Theorem~\ref{thm:hp}.
Assume that on the same high-probability event one further has
\begin{equation}\label{eq:struct_aux}
  \mathcal R_N(\nperm,\tau) \le C_R\,\eta_N,
  \qquad
  Q_N \le C_Q\,\eta_N^4
\end{equation}
for deterministic constants $C_R, C_Q > 0$ independent of~$N$.
Then, with probability at least $1-\delta$,
\begin{equation}\label{eq:struct_rate}
  \boxed{
  \cB
  \;\le\;
  C_{\mathrm{fast}}(\delta,N)\,\eta_N^2
  + C_{\mathrm{sc}}(\delta,N)\,\frac{m_{2,*}}{N}
  \;=\;
  O\!\left(\kappa_W^2\frac{d}{N} + \frac{K}{N}\right).
  }
\end{equation}
\end{corollary}

Corollary~\ref{cor:struct_rate} converts the transfer bottleneck from
$O(\kappa_W\sqrt{d/N})$ to $O(\kappa_W^2 d/N)$ without requiring
full-parameter alignment.

\paragraph{The $K=1$ case.}
For a single Gaussian component ($K=1$), the head regularity
hypothesis is significantly easier to satisfy.
The centring projection~\eqref{eq:centred_head} is trivial
(the $\pi$-block is one-dimensional, so
$\tilde v_{i,\pi} = 0$ for all $i$),
leaving only the $(\mu, s)$ block.
Furthermore, the semi-convexity constant reduces to
$\rho_{K=1} \le 1/s_{\min}^2$---independent of
$\Lambda$---because the covariance term
$\mathrm{Cov}_\gamma(\nabla u_k)$ vanishes when $K=1$
(Proposition~\ref{prop:loss_regularity}).
Together with the empirical observation that the barrier
decays as $O(N^{-1.3})$ at $K=1$
(Appendix~\ref{app:nonvacuous_detail}),
faster than the baseline $O(N^{-1/2})$ but consistent with
the $O(1/N)$ rate of Corollary~\ref{cor:struct_rate},
this provides strong evidence that head regularity holds
for single-component MDNs.

The bounded Lipschitz activation assumption
(Assumption~\ref{ass:activ}\,(i)--(ii)) is shared with
Ferbach et al.~\cite{ferbach2024}, who likewise require bounded Lipschitz
activations in the mean-field two-layer setting; our baseline results
(Theorems~\ref{thm:general}--\ref{thm:hp},
Corollary~\ref{cor:asymptotic}) therefore do not impose any
additional restriction on the activation function
beyond the existing literature.
The head-regularity result (Corollary~\ref{cor:struct_rate})
additionally requires $C^2$ smoothness
(Assumption~\ref{ass:activ}\,(iii)).

\subsection{Application to von Mises--Fisher Mixtures}
\label{sec:vmf}

We now instantiate Theorem~\ref{thm:general} on a second
mixture-model family: von Mises--Fisher (vMF) mixtures on the unit
sphere.

\paragraph{Motivation.}
Normalised embeddings and cosine similarity have become standard in
modern deep learning, making the unit sphere a natural output space.
vMF mixture density networks---where a neural network outputs the
parameters of a vMF mixture conditioned on the input---arise in
several active areas:
\emph{(i)}~probabilistic 3D orientation and pose estimation, where
multimodal angular predictions are modelled by vMF
mixtures~\cite{prokudin2018deep,liu2024probabilistic};
\emph{(ii)}~molecular property prediction with uncertainty
quantification via spherical mixture density
networks~\cite{han2023sphericalmdn};
and \emph{(iii)}~directional prediction in particle physics
detectors~\cite{ghrear2024direction}.
Linear mode connectivity guarantees are directly relevant whenever
such models are independently trained and subsequently merged, e.g.\
in federated or distributed learning scenarios.

The key finding of this section is that the semi-convexity constant
is a \emph{universal constant}, independent of all model parameters.

\paragraph{Architecture and loss.}
A vMF mixture network models conditional densities on
$S^{d_y-1} \subset \R^{d_y}$:
\[
  p(y \mid x) = \sum_{k=1}^K \pi_k(x)\;
  C_{d_y}\!\bigl(\kappa_k(x)\bigr)\,
  \exp\!\bigl(\kappa_k(x)\,\mu_k(x)^\top y\bigr),
\]
where $\mu_k \in S^{d_y-1}$ is the mean direction,
$\kappa_k \ge 0$ the concentration, and
$C_{d_y}(\kappa) =
\kappa^{d_y/2-1}\bigl/\bigl((2\pi)^{d_y/2}\,
I_{d_y/2-1}(\kappa)\bigr)$
the normalising constant
($I_\nu$: modified Bessel function of the first kind).
We adopt the \emph{natural parameterisation}
$\eta_k \coloneqq \kappa_k\mu_k \in \R^{d_y}$
(so that $\kappa_k=\|\eta_k\|$, $\mu_k=\eta_k/\|\eta_k\|$).
The two-layer network output is
$z(x) = \frac{1}{N}\sum_{i=1}^N v_i\,\sigma(w_i^\top x)
\in \R^{K(1+d_y)}$,
partitioned as $z = (z_\pi,\,\eta_1,\ldots,\eta_K)$ with
$z_\pi \in \R^K$ and $\eta_k \in \R^{d_y}$.
The NLL is
\begin{equation}\label{eq:vmf_loss}
  \cL_{\vmf}(z,y)
  = -\log\sum_{k=1}^K \pi_k\,
    C_{d_y}(\|\eta_k\|)\,\exp(\eta_k^\top y),
\end{equation}
with $\pi_k = \softmax(z_\pi)_k$.
Setting $u_k = \log\pi_k + \log C_{d_y}(\|\eta_k\|) + \eta_k^\top y$,
the log-sum-exp Hessian
decomposition~\eqref{eq:hess_decomp} applies verbatim.

\paragraph{Component permutation symmetry.}
As with MDNs, $\cL_{\vmf}(P_\tau z,\,y) = \cL_{\vmf}(z,\,y)$
for any $\tau \in \cS_K$, so two-level alignment
$(\nperm,\tau) \in \cS_N \times \cS_K$ applies.

\begin{proposition}[vMF loss regularity]\label{prop:vmf_regularity}
For $y \in S^{d_y-1}$ and any $z \in \R^{K(1+d_y)}$:
\begin{equation}\label{eq:vmf_bounds}
  C_{L,\vmf}^\infty \;\le\; 2 + 2\sqrt{d_y},
  \qquad
  \rho_{\vmf} \;\le\; 6.
\end{equation}
In particular, $\rho_{\vmf}$ is a universal constant---independent
of $K$, $\kappa_k$, $R$, and all other model parameters.
No architectural constraints (scale floors, concentration ceilings)
are required.
\end{proposition}

\begin{proof}[Proof sketch]
Because vMF belongs to the exponential family with natural parameter
$\eta_k$ and sufficient statistic $y$, the per-component
log-normaliser $\psi(\eta) \coloneqq -\log C_{d_y}(\|\eta\|)$ is
convex, so
$\nabla^2_{\eta_k} u_k = -\nabla^2\psi(\eta_k) \preceq 0$.
Combined with the PSD softmax Hessian $C_\pi \succeq 0$, the
weighted-Hessian term satisfies
$-\sum_k \gamma_k \nabla^2 u_k \succeq 0$
(cf.\ MDN, where this contributes $-1/s_{\min}^2$).
For the covariance term, the gradient
$\nabla_{\eta_k} u_k = y - A_{d_y}(\kappa_k)\,\mu_k$,
where $A_{d_y}(\kappa) = I_{d_y/2}(\kappa)/I_{d_y/2-1}(\kappa)
\in [0,1)$ is the mean resultant length.
Since $\|y\|=1$ and $A_{d_y}(\kappa_k)\|\mu_k\|<1$, we have
$\|\nabla_{\eta_k} u_k\| \le 2$.
Together with the $z_\pi$-block bound $\|e_k-\pi\|\le\sqrt{2}$:
$\max_k\|\nabla u_k\|^2 \le 6$, giving
$\rho_{\vmf} \le 6$ via~\eqref{eq:hess_decomp}.
The Lipschitz bound follows from
$C_{L,\vmf}^\infty
= \sup_z\|\nabla\cL_{\vmf}\|_1 \le 2 + 2\sqrt{d_y}$
(full details in Appendix~\ref{app:vmf_proof}).
\end{proof}

Substituting~\eqref{eq:vmf_bounds} into
Theorem~\ref{thm:general} with $p=K(1+d_y)$ yields:

\begin{corollary}[vMF barrier under $W$-OT alignment]
\label{cor:vmf_barrier}
Under Assumptions~\ref{ass:activ}, \ref{ass:input_mom},
\ref{ass:mf_law}, and~\ref{ass:subgauss}
with $p = K(1+d_y)$ and $y \in S^{d_y-1}$,
the vMF mixture barrier satisfies, with probability $\ge 1-\delta$:
\begin{equation}\label{eq:vmf_rate}
  \boxed{
  \cB_{\vmf}
  \;=\;
  O\!\left((1+\sqrt{d_y})\,\kappa_W\sqrt{\frac{d}{N}}
  + \frac{K}{N}\right)
  }
\end{equation}
with prefactors that are universal constants---no dependence on
$s_{\min}$, $\Lambda$, or domain-radius~$R$.
\end{corollary}

\paragraph{Comparison.}
Table~\ref{tab:loss_comparison} contrasts the loss-geometry
constants.
The vMF case achieves dramatically better constants because
(i)~the observation space is compact ($S^{d_y-1}$), bounding
gradient norms without domain-radius factors, and
(ii)~the exponential-family structure makes each $u_k$ concave
in $\eta_k$, so the weighted Hessian contributes zero rather than
$-1/s_{\min}^2$.

\begin{table}[t]
\centering\small
\caption{Loss-geometry constants for the two mixture-model
  applications of Theorem~\ref{thm:general}.}
\label{tab:loss_comparison}
\begin{tabular}{@{}lcc@{}}
\toprule
 & \textbf{MDN (Gaussian)} & \textbf{vMF Mixture} \\
\midrule
Observation space & $\R$ (bounded by $Y_{\max}$) & $S^{d_y-1}$ (compact) \\
Output dim per component & $3$ & $1+d_y$ \\
$C_L^\infty$ & $O(\Lambda^2/s_{\min}^2)$ & $O(\sqrt{d_y})$ \\
$\rho$ & $O(\Lambda^4/s_{\min}^4)$ & $\le 6$ (universal) \\
Architectural constraint & scale floor $s_{\min}$ & none \\
\bottomrule
\end{tabular}
\end{table}

%======================================================================
\section{Alignment Algorithms}
\label{sec:algorithm}
%======================================================================

Theorem~\ref{thm:general} guarantees that
\emph{some} pair $(\nperm, \tau) \in \cS_N \times \cS_K$ achieves a
vanishing barrier.
We first show that permutations are the \emph{only} barrier-relevant
symmetries for softmax-based mixture models, then develop a practical
alignment algorithm.

\subsection{Symmetry and Barrier Neutrality}\label{sec:symmetry}

Beyond $\cS_N$ and $\cS_K$, any softmax-based mixture model
(including both MDNs and vMF mixtures) has a continuous symmetry:
\emph{softmax translation}
$v_i \mapsto v_i + \delta_i\, e_\pi$
(where $e_\pi = (\mathbf{1}_K, \mathbf{0}_{p-K}) \in \R^{p}$)
adds a common-mode shift to the $\pi$-logits, which softmax absorbs.
The full symmetry group is
$\cS_N \times \cS_K \times \R^N_\pi$
(Proposition~\ref{prop:symmetry}, Appendix~\ref{app:symmetry_full}).
Crucially, softmax translation is barrier-neutral:

\begin{proposition}[Barrier neutrality of softmax translation]
\label{prop:barrier_neutrality}
For any mixture-model parameter sets $\theta^A, \theta^B$ and arbitrary
$\delta \in \R^N$, setting
$\theta^{B'} = \{(w_i^B, v_i^B + \delta_i\,e_\pi)\}$ yields
$\cB(\theta^A, \theta^{B'}) = \cB(\theta^A, \theta^B)$.
\end{proposition}

\noindent
The proof (Appendix~\ref{app:proof_neutrality}) is a direct
interpolation calculation.
The alignment task therefore reduces to finding the permutation
pair $(\nperm, \tau)$.

\subsection{Hellinger-Affinity Matching}\label{sec:hellinger}

The key design choice is \emph{how to measure component correspondence}:
weight matching~\cite{ainsworth2023} compares raw parameter
vectors, but two components with different weight vectors can produce
nearly identical conditional Gaussians.
We instead match components at the \emph{distributional level} via the
Hellinger affinity.

Given calibration inputs $\{x_n\}_{n=1}^{n_{\mathrm{cal}}}$, define the
\emph{Hellinger affinity} between component $j$ of network $A$ and
component $k$ of network $B$ as
\begin{equation}\label{eq:hellinger_cost}
  H_{jk} = \sum_{n=1}^{n_{\mathrm{cal}}}
    \underbrace{\sqrt{\pi_{A,j}(x_n)\,\pi_{B,k}(x_n)}}_{\text{weight similarity}}
    \cdot\;
    \underbrace{\mathrm{BC}\!\bigl(
      \cN(\mu_{A,j}(x_n),\,s_{A,j}(x_n)^2),\;
      \cN(\mu_{B,k}(x_n),\,s_{B,k}(x_n)^2)
    \bigr)}_{\text{shape similarity}},
\end{equation}
where
\[
  \mathrm{BC}\!\bigl(\cN(\mu_1,s_1^2),\cN(\mu_2,s_2^2)\bigr)
  = \sqrt{\frac{2s_1 s_2}{s_1^2+s_2^2}}\,
    \exp\!\Bigl(-\frac{(\mu_1-\mu_2)^2}{4(s_1^2+s_2^2)}\Bigr)
\]
is the Bhattacharyya coefficient between two 1-D Gaussians.
Each term in~\eqref{eq:hellinger_cost} is the geometric mean of the
two components' mixture weights $\sqrt{\pi_{A,j}\pi_{B,k}}$ multiplied
by their distributional similarity $\mathrm{BC}$, both evaluated at
input $x_n$.
(Note: $\pi_{A,j}(x_n)$ here denotes the softmax mixture weight output
by network $A$, distinct from the EM responsibilities
$\gamma_k(x,y)$ used in the Hessian analysis of Section~\ref{sec:output_geometry}.)
$H_{jk}$ is large when components $j$ and $k$ model similar
conditional distributions across the calibration set.
The BC formula above is for 1-D Gaussians (MDN); for other
mixture families (e.g.\ vMF), one substitutes the appropriate
Bhattacharyya coefficient.
The component permutation is then found by solving the LAP:
\[
  \tau \leftarrow \arg\max_{q \in \cS_K}\sum_{k=1}^K H_{k,q(k)},
\]
which finds the bijective matching that maximises total distributional
affinity across components.
Because the output distribution is invariant to hidden-neuron
permutation, $H_{jk}$ is also invariant to $\nperm$, so it can be
computed independently and reused across algorithms.

The theoretical bound (Theorem~\ref{thm:hp}) uses $\tau = \mathrm{id}$,
which already gives $\Delta_N = O(1/N)$ under the common mean-field
law.
In practice, the Hellinger-affinity $\tau$ can further reduce
$\Delta_N$ at finite~$N$; a formal guarantee that it minimises
$\Delta_N$ is left as an open problem.

\paragraph{Full algorithm.}
Combining activation correlation for hidden neurons
(following Ainsworth et al.~\cite{ainsworth2023}) with the Hellinger
affinity for components yields Algorithm~\ref{alg:am}.
Because $H_{jk}$ is invariant to $\nperm$, the two stages are
independent and no iteration is needed.

\begin{algorithm}[t]
\caption{Activation Matching with Hellinger Head (AM)}
\label{alg:am}
\begin{algorithmic}[1]
\REQUIRE Mixture-model parameters $\theta_A, \theta_B$; components $K$;
  calibration set $\{x_n\}_{n=1}^{n_{\mathrm{cal}}}$.
\ENSURE $\nperm \in \cS_N$, $\tau \in \cS_K$.
\STATE \textbf{Hidden permutation} (activation-space):
  form the $N \times N$ inner-product matrix
  $C_{ij}^{\mathrm{AM}} = \sum_{n=1}^{n_{\mathrm{cal}}} \sigma(w_{A,i}^\top x_n)\,\sigma(w_{B,j}^\top x_n)$.
  Solve $\nperm \leftarrow \arg\max_{q \in \cS_N}\sum_i C_{i,q(i)}^{\mathrm{AM}}$
  via the Hungarian algorithm.
  \hfill\COMMENT{$O(n_{\mathrm{cal}}N^2 + N^3)$}
\STATE \textbf{Component permutation} (distribution-space):
  compute $H$ via~\eqref{eq:hellinger_cost}.
  Solve $\tau \leftarrow \arg\max_{q \in \cS_K}\sum_k H_{k,q(k)}$.
  \hfill\COMMENT{$O(n_{\mathrm{cal}}K^2 + K^3)$}
\RETURN $(\nperm, \tau)$.
\end{algorithmic}
\end{algorithm}

%======================================================================
\section{Experiments}\label{sec:experiments}
%======================================================================

We validate the framework empirically on both MDN and vMF mixture
models.

\paragraph{Setup.}
Throughout this section we fix the alignment backend of
Algorithm~\ref{alg:am} ($W$-OT hidden matching $+$ Hellinger
component matching) and compare \emph{initialisation schemes}
for the MDN heads:
\textbf{random}, \textbf{$k$-means} (initial means are set by
$k$-means and then trained), and \textbf{$k$-means frozen}
(the same initial means, kept fixed during training).
The reported runs use the two synthetic benchmarks already employed in
Appendix~\ref{app:assumption_validation},
\texttt{synthetic} and \texttt{synthetic\_transpose},
with $K=5$, widths
$N \in \{512, 1024, 2048, 4096\}$,
full-batch gradient descent, and $10$ random seeds per setting.
For each $(\text{dataset}, N, \text{init})$ combination we evaluate
the $5$ seed pairs
$(0\text{--}1, 2\text{--}3, \ldots, 8\text{--}9)$
already cached in the numerical study.
We record the aligned barrier $\cB$, the unaligned barrier,
the effective scale floor $s_{\min}$, and the transfer-side
diagnostics $M_{\Delta \tilde V}^\infty$ and $Q_N$ from the refined
coupling of Proposition~\ref{prop:transfer_refined}.
These diagnostics \emph{probe} the head-regularity route of
Corollary~\ref{cor:struct_rate}; they do not by themselves verify the
full structural hypothesis, since $\mathcal{R}_N$ is not estimated
directly.

\paragraph{Experimental questions.}
\begin{enumerate}[label=(\arabic*),leftmargin=*,itemsep=2pt]
  \item \emph{Component anchoring and barrier reduction
    (Figure~\ref{fig:kmeans_diag}, Table~\ref{tab:kmeans_rep}):}
    does symmetry breaking at initialisation materially lower the
    aligned barrier, even when the matching backend is held fixed?
  \item \emph{Transfer-side diagnostics
    (Figure~\ref{fig:kmeans_diag}, Table~\ref{tab:kmeans_rep}):}
    do the initialisations that reduce $\cB$ also improve
    $s_{\min}$, $M_{\Delta \tilde V}^\infty$, and $Q_N$,
    moving the models toward the regime suggested by
    Proposition~\ref{prop:transfer_struct}?
  \item \emph{Qualitative component stability
    (Figure~\ref{fig:kmeans_components}):}
    do anchored initialisations produce more persistent component
    identities across random seeds?
\end{enumerate}

\paragraph{Barrier and diagnostic trends.}
Figure~\ref{fig:kmeans_diag} summarises the width dependence of the
aligned barrier, the effective scale floor $s_{\min}$, and the
gauge-fixed head mismatch $M_{\Delta \tilde V}^\infty$.
The main empirical pattern is that $k$-means-based initialisation can
substantially improve the post-training geometry, but the strongest
variant is \emph{dataset-dependent}.
On \texttt{synthetic}, the frozen anchor is best: at $N=1024$ the
aligned barrier drops from $8.92$ (random) to $0.202$, while
$M_{\Delta \tilde V}^\infty$ falls from $41.6$ to $15.4$ and
$Q_N$ from $3.03 \times 10^5$ to $1.18 \times 10^5$.
On \texttt{synthetic\_transpose}, the learnable $k$-means
initialisation is clearly superior: at $N=1024$ the aligned barrier
drops from $5.81$ to $0.229$, $s_{\min}$ rises from
$1.75 \times 10^{-5}$ to $1.19 \times 10^{-3}$, and
$Q_N$ falls by roughly one order of magnitude.
Because the unaligned barriers remain large
(Table~\ref{tab:kmeans_rep}),
the gain is not merely that the trained models become trivially close
before alignment; the initialisation changes the geometry of the final
solutions in a way the alignment can exploit.

\begin{table}[t]
\centering
\small
\renewcommand{\arraystretch}{1.12}
\caption{Representative $K=5$, full-batch results at $N=1024$
  (average over $5$ aligned seed pairs).
  Lower $\cB$, $M_{\Delta \tilde V}^\infty$, and $Q_N$ are better;
  larger $s_{\min}$ indicates a wider effective scale floor.}
\label{tab:kmeans_rep}
\begin{tabular}{llccccc}
\toprule
Dataset & Init & $\cB$ & $\cB_{\mathrm{unal.}}$ & $s_{\min}$ &
$M_{\Delta \tilde V}^\infty$ & $Q_N$ \\
\midrule
\texttt{synthetic}
  & random & 8.92 & 41.8 & $7.85 \times 10^{-5}$ & 41.6 & $3.03 \times 10^5$ \\
\texttt{synthetic}
  & $k$-means & 7.65 & 121 & $\mathbf{1.33 \times 10^{-4}}$ & 39.0 & $1.75 \times 10^5$ \\
\texttt{synthetic}
  & $k$-means frozen & $\mathbf{0.202}$ & 40.6 & $1.08 \times 10^{-4}$ & $\mathbf{15.4}$ & $\mathbf{1.18 \times 10^5}$ \\
\midrule
\texttt{synthetic\_transpose}
  & random & 5.81 & $1.23 \times 10^3$ & $1.75 \times 10^{-5}$ & 94.1 & $5.52 \times 10^6$ \\
\texttt{synthetic\_transpose}
  & $k$-means & $\mathbf{0.229}$ & 87.7 & $\mathbf{1.19 \times 10^{-3}}$ & $\mathbf{41.1}$ & $\mathbf{5.50 \times 10^5}$ \\
\texttt{synthetic\_transpose}
  & $k$-means frozen & 6.45 & 65.7 & $2.62 \times 10^{-5}$ & 58.5 & $5.71 \times 10^6$ \\
\bottomrule
\end{tabular}
\end{table}

\begin{figure}[t]
\centering
\includegraphics[width=\textwidth]{kmeans_init_diagnostics_full_batch.pdf}
\caption{Full-batch $K=5$ diagnostics with fixed
  $W$-OT $+$ Hellinger alignment.
  Columns show the aligned empirical barrier, the effective scale floor
  $s_{\min}$, and the refined head mismatch
  $M_{\Delta \tilde V}^\infty$ as width varies.
  The best anchored variant depends on the dataset, but the variants
  with the smallest barrier also improve the transfer-side
  diagnostics, which is qualitatively consistent with the
  head-regularity route.}
\label{fig:kmeans_diag}
\end{figure}

\paragraph{Qualitative component geometry.}
Figure~\ref{fig:kmeans_components} shows the learned mixture
components across all $10$ seeds at $N=1024$.
For \texttt{synthetic}, random initialisation often yields visibly
redundant or poorly separated heads, whereas the frozen anchor produces
cleaner, more persistent component families.
For \texttt{synthetic\_transpose}, the learnable $k$-means
initialisation produces markedly more stable seed-to-seed component
trajectories than the random baseline.
We therefore interpret the practical benefit not as a single universal
``best'' initialiser, but as \emph{component anchoring}:
better symmetry breaking at initialisation can steer training toward
solutions with lower aligned barrier and better transfer diagnostics.

\begin{figure}[t]
\centering
\includegraphics[width=0.49\textwidth]{synthetic/K5/full_batch/predictive/w1024_components.png}
\hfill
\includegraphics[width=0.49\textwidth]{synthetic/K5_kmeans_frozen/full_batch/predictive/w1024_components.png}

\vspace{0.6em}

\includegraphics[width=0.49\textwidth]{synthetic_transpose/K5/full_batch/predictive/w1024_components.png}
\hfill
\includegraphics[width=0.49\textwidth]{synthetic_transpose/K5_kmeans/full_batch/predictive/w1024_components.png}
\caption{Component plots at $N=1024$ for the two datasets.
  Top row: \texttt{synthetic} with random initialisation (left) and
  $k$-means frozen (right).
  Bottom row: \texttt{synthetic\_transpose} with random initialisation
  (left) and learnable $k$-means (right).
  Each panel aggregates all $10$ seeds.
  The anchored initialisations yield visibly more stable component
  identities across seeds.}
\label{fig:kmeans_components}
\end{figure}

\subsection{vMF Mixture Experiments}\label{sec:vmf_experiments}

We now validate the vMF instantiation
(Corollary~\ref{cor:vmf_barrier}) on synthetic directional data on
$S^1$ ($d_y = 2$), using the same two-level alignment backend.

\paragraph{Setup.}
We generate a \texttt{rotating} dataset: $K$ vMF components on $S^1$
whose mean directions rotate with the input $x \in [-1,1]$, with
varying concentration~$\kappa$.
Training follows the MDN protocol: 10K epochs of full-batch Adam with
cosine annealing, widths $N \in \{512, 1024, 2048, 4096\}$, and 10
random seeds.
We compute barriers for 5~seed pairs and report averages.

\paragraph{Bound quality.}
Table~\ref{tab:mdn_vmf_comparison} contrasts the barrier
\emph{bound quality} (ratio $\cB_{\mathrm{wc}}/\cB_{\mathrm{emp}}$)
across the two loss families.
The settings differ ($K$, dataset, output dimension), so absolute
barrier values are not directly comparable.
The key contrast is in the constants: $\rho_{\vmf} = 6$ versus
the MDN's $\rho \sim 10^{20}$, and $C_L^\infty = 4.83$ versus
$\sim\!6 \times 10^4$.
Consequently, the worst-case bound-to-empirical ratio drops from
$\sim\!10^{21}$ (MDN) to $\sim\!3 \times 10^3$ (vMF).
The data-dependent ratios are more comparable ($\sim\!6 \times 10^3$
vs.\ $\sim\!1.3 \times 10^3$), but the vMF achieves this
\emph{without any architectural constraints} (no scale floor,
no clipping).

\input{table_comparison_draft.tex}

Figure~\ref{fig:vmf_barrier_scaling} shows the barrier-vs-width
scaling for both $K = 1$ and $K = 3$.
For $K = 1$, the aligned barrier decays cleanly from $0.019$ at
$N = 512$ to $0.004$ at $N = 4096$, consistent with the predicted
$O(1/N)$ rate.
For $K = 3$, the aligned barrier is approximately flat around $1.0$
at the tested widths; the output-space mismatch $\Delta_N$ is not yet
converging, indicating that the component alignment quality---rather
than width---is the bottleneck at these scales.
In both cases, alignment reduces the barrier by one to two orders of
magnitude relative to the unaligned baseline.

\begin{figure}[t]
\centering
\includegraphics[width=0.49\textwidth]{rotating/K1/full_batch/barrier_scaling.pdf}
\hfill
\includegraphics[width=0.49\textwidth]{rotating/K3/full_batch/barrier_scaling.pdf}
\caption{vMF barrier scaling on $S^1$ (\texttt{rotating} dataset).
  Left: $K = 1$; the aligned barrier (blue) decays as $\sim 1/N$.
  Right: $K = 3$; the aligned barrier is flat at the tested widths,
  dominated by the component-mismatch $O(K/N)$ term.
  Theoretical worst-case and data-dependent bounds (triangles) both
  track the empirical barrier with a gap of $\sim\!10^3$.}
\label{fig:vmf_barrier_scaling}
\end{figure}

\paragraph{Real-data validation.}
We also run the vMF pipeline on a real directional-data task: predicting
wind direction (on~$S^1$) from temperature, pressure, and humidity
in the Jena Climate dataset~\cite{jena_climate}.
With $K = 1$, 3 seed pairs, and $N \in \{512, 1024\}$, the aligned
barrier decays from $0.29$ to $0.21$, consistent with the synthetic
results.
The data-dependent bound ($\sim\!850$--$990$) and bound-to-empirical
ratio ($\sim\!3{,}400$--$4{,}000$) mirror the synthetic experiment,
confirming that the vMF framework produces meaningful bounds on
real data.

%======================================================================
\section{Conclusion}\label{sec:conclusion}
%======================================================================

We have established a general linear mode connectivity framework for
two-layer networks with semi-convex losses
(Theorem~\ref{thm:general}), removing the convexity assumption that
underpins all prior theoretical guarantees.
The framework decomposes the barrier into a loss-independent transfer
term and a semi-convexity penalty, both vanishing with width.
Instantiating the framework on MDNs yields a barrier of
$O(\kappa_W\sqrt{d/N}+K/N)$ via a log-sum-exp Hessian decomposition
that avoids the doubly-exponential dependence on $p_{\min}$.
A second instantiation on vMF mixtures demonstrates the framework's
generality: the exponential-family structure yields a \emph{universal}
$\rho \le 6$ with no architectural constraints, producing clean,
parameter-free barrier bounds.
Experiments on synthetic $S^1$ data illustrate the framework:
for $K = 1$, the aligned barrier decays as $\sim 1/N$
(Figure~\ref{fig:vmf_barrier_scaling}), and the vMF worst-case bound
is finite and width-dependent, in contrast to the MDN's
vacuous $\sim\!10^{21}$
(Table~\ref{tab:mdn_vmf_comparison}).
For $K = 3$, the barrier does not yet decay at the tested widths,
suggesting that the component-alignment quality is the practical
bottleneck.
A conditional head-regularity hypothesis improves the rate
to $O(d/N + K/N)$ for both families.

\paragraph{Scope of the mean-field framework.}
Under the mean-field parameterisation, both the transfer term
and the semi-convexity penalty vanish as $N\to\infty$
regardless of loss convexity, so the asymptotic conclusion
$\cB\to 0$ holds for any $\rho$-semi-convex loss with finite~$\rho$.
The contribution of the present work is therefore not the
asymptotic conclusion itself, but the \emph{finite-width,
loss-specific} barrier decomposition.
The semi-convexity constant~$\rho$ enters the bound through
the output-interpolation inequality---a property of the loss
function, independent of the weight distribution---and
determines the barrier at practical widths: the semi-convexity
penalty $\rho\cdot O(p/N)$ can dominate the transfer term
$O(\sqrt{d/N})$ when $\rho$ is large.
This is visible in the MDN vs.\ vMF contrast:
$\rho_{\mathrm{MDN}} = O(\Lambda^4/s_{\min}^4)$ yields
a vacuous bound ($\sim\!10^{21}$, Table~\ref{tab:mdn_vmf_comparison})
at moderate widths, whereas the universal
$\rho_{\vmf} \le 6$ gives a finite, width-dependent bound,
a prediction confirmed experimentally
(Figure~\ref{fig:vmf_barrier_scaling}).
The gap between mean-field theory and practical (finite-width,
trained) networks is shared by all works in this line, including
ours; closing this gap remains an important open problem.

\paragraph{Limitations.}
Our theory is restricted to two-layer networks under the $1/N$
mean-field parameterisation; extending to deeper architectures
requires layer-wise propagation of the semi-convexity bound.
For MDNs, the loss regularity constants
(Proposition~\ref{prop:loss_regularity}) depend polynomially on
$1/s_{\min}$; the practical tightness depends on the trained
$s_{\min}$.
The mean-field i.i.d.\ assumption (Assumption~\ref{ass:mf_law})
is justified via propagation of chaos but becomes less accurate
for long training horizons or narrow networks.
The theoretical bound uses $\tau = \mathrm{id}$, which suffices
under the common mean-field law; the Hellinger affinity criterion
used in practice lacks a formal guarantee of minimising $\Delta_N$
at finite~$N$.

\paragraph{Future work.}
Natural extensions include:
(i)~multi-layer networks using iterated output-coupling,
(ii)~further mixture-model instantiations (Poisson, Beta, Gamma
mixtures) that share the log-sum-exp structure,
(iii)~extension to multivariate MDNs ($y \in \R^p$, $p > 1$)
and vMF mixtures on products of spheres,
(iv)~a formal link between Hellinger affinity and $\Delta_N$
minimisation, and
(v)~validation on real directional-data benchmarks beyond synthetic $S^1$.

%======================================================================
% APPENDICES
%======================================================================
\clearpage
\appendix

%======================================================================
\section{Notation Reference}\label{app:notation}
%======================================================================

Table~\ref{tab:notation} collects all notation.
All norms $\|\cdot\|$ are Euclidean ($\ell_2$) unless noted;
all logarithms are natural.
\emph{Subscript conventions:}
a numeric subscript $k$ indexes mixture components ($\pi_k, \mu_k, s_k$);
a subscript $i$ indexes neurons ($v_i, w_i$);
a subscript $*$ denotes an asymptotic ($N\!\to\!\infty$) limiting quantity
($m_{2,*}, m_{\infty,*}, \mathfrak{m}_*$);
a subscript $N$ denotes a finite-sample ($N$-neuron) empirical quantity
($A_N, B_N, \Delta_N, \eta_N$).

\begin{table}[H]
\centering
\caption{Complete notation reference.}\label{tab:notation}
\scriptsize
\setlength{\tabcolsep}{3pt}
\renewcommand{\arraystretch}{1.05}
\begin{multicols}{2}
\noindent
\begin{tabular}[t]{@{}lp{4.6cm}@{}}
\toprule
\textbf{Symbol} & \textbf{Description} \\
\midrule
\addlinespace[2pt]
\multicolumn{2}{@{}l}{\textbf{Architecture \& dimensions}} \\
\addlinespace[1pt]
$N \in \N$     & Network width \\
$d \in \N$     & Input dimension \\
$p \in \N$     & Output dimension ($3K$ for MDN, $K(1\!+\!d_y)$ for vMF) \\
$K \in \N$     & Number of mixture components \\
$\sigma\colon\R\!\to\!\R$ & Activation function (not to be confused with $\kappa_W$) \\
$v_i \in \R^{p}$ & Head weight for neuron $i$ \\
$w_i \in \R^d$    & Input weight for neuron $i$ \\
$\theta_i\!=\!(w_i,v_i) \in \R^{d+p}$ & Full parameter for neuron $i$ \\
$z(x) \in \R^{p}$ & Raw network output \\
$z_\pi \in \R^K$ & Mixing-weight block (shared by all mixture losses) \\
\addlinespace[3pt]
\multicolumn{2}{@{}l}{\textbf{MDN-specific quantities}} \\
\addlinespace[1pt]
$z_\mu, z_s \in \R^K$ & Mean, log-scale blocks (MDN) \\
$[v]_{(k)} \in \R^3$ & $k$-th component block of $v$ (MDN) \\
$\pi_k \in [0,1]$ & Mixture weight; $\softmax(z_\pi)_k$ \\
$\mu_k \in \R$ & Component mean; $z_{\mu,k}$ \\
$s_k \in \R_{>0}$ & Component std.\ dev.; $e^{z_{s,k}}$ \\
$s_{\min} \in \R_{>0}$ & Scale floor; $\min_k s_k$ (not a component) \\
$\varphi_k \in \R_{>0}$ & Gaussian $\cN(y;\mu_k,s_k^2)$ \\
$u_k \in \R$   & $\log(\pi_k\varphi_k)$ \\
$\gamma_k \in [0,1]$ & Responsibility; $e^{u_k}/\!\sum_j e^{u_j}$ \\
\addlinespace[3pt]
\multicolumn{2}{@{}l}{\textbf{Loss \& barrier}} \\
\addlinespace[1pt]
$\ell(z,y) \in \R$ & Loss function (general) \\
$\cL_{\mathrm{MDN}}(z,y)$ & MDN negative log-likelihood \\
$\cL_{\vmf}(z,y)$ & vMF negative log-likelihood \\
$\cE(\theta) \in \R$ & Population loss \\
$\cB \in \R$  & Interpolation barrier \\
$M_t \in \R^{N(p+d)}$ & Interpolated params; $(1\!-\!t)\theta_A\!+\!t\theta_B$ \\
$z_{M_t} \in \R^{p}$ & Output along param.\ path \\
$\hat z_t \in \R^{p}$ & Output interpolation; $(1\!-\!t)z_A\!+\!tz_B$ \\
\addlinespace[3pt]
\multicolumn{2}{@{}l}{\textbf{Loss-geometry constants}} \\
\addlinespace[1pt]
$C_L^\infty \in \R_+$ & $\ell_\infty$-Lipschitz const.\ ($\|\nabla\cL\|_1$) \\
$\bar C_L \in \R_+$ & Data-dep.\ Lipschitz const.~\eqref{eq:CL_dd} (App.~\ref{app:data_dep_barrier}) \\
$C_L \in \R_+$ & $\ell_2$-Lipschitz const.\ on $\cD_R$ \\
$\rho \in \R_+$ & Semi-conv.\ const.; $\nabla^2\cL\!\succeq\!{-}\rho I$ \\
$\rho_{\mathrm{dir}}(x,y) \in \R_+$ & Directional semi-conv.\ cost~\eqref{eq:rho_dir} (App.~\ref{app:data_dep_barrier}) \\
$\bar\chi \in [0,1]$ & Competition mass; $\sum_k\gamma_k(1\!-\!\gamma_k)$ \\
$Q(a) \in \R^{2\times 2}$ & Hessian building block; $\bigl(\begin{smallmatrix}1 & 2a\\2a & 2a^2\end{smallmatrix}\bigr)$ \\
$R \in \R_+$  & Output-domain radius; $S_\sigma M_V^\infty$ \\
$Y_{\max} \in \R_+$ & Response bound; $|y|\!\le\! Y_{\max}$ \\
$\Lambda \in \R_+$ & $Y_{\max}+R$ \\
$S_\sigma \in \R_+$ & Activation bound; $|\sigma|\!\le\! S_\sigma$ \\
$L_\sigma \in \R_+$ & Activation Lipschitz const. \\
$L_{\sigma'} \in \R_+$ & $\|\sigma''\|_\infty$; curvature const. \\
$M_V^\infty \in \R_+$ & Per-coord.\ head moment~\eqref{eq:head_moment_inf} \\
$\bar M_{V,N} \in \R_+$ & RMS head norm; $(\frac{1}{N}\!\sum_i\!\|P_\tau v_{B,\nperm(i)}\|^2)^{1/2}$ \\
$M_x \in \R_+$ & Input moment bound; $\sqrt{\E_P[\|x\|^2]}$ \\
$M_{x,4} \in \R_+$ & 4th input moment; $(\E_P[\|x\|^4])^{1/4}$ \\
$\kappa_W \in \R_+$ & Sub-Gaussian param.\ of $\mathfrak{m}_{*,W}$ \\
$\Sigma_x \in \R^{d\times d}$ & Input covariance; $\E_P[xx^\top]$ \\
\bottomrule
\end{tabular}

\columnbreak

\noindent
\begin{tabular}[t]{@{}lp{4.6cm}@{}}
\toprule
\textbf{Symbol} & \textbf{Description} \\
\midrule
\addlinespace[2pt]
\multicolumn{2}{@{}l}{\textbf{Domains \& norms}} \\
\addlinespace[1pt]
$\cD_R \subset \R^{p}$ & $\{z : \|z\|_\infty\!\le\! R\}$ \\
$\cD_{R,s_{\min}} \subset \cD_R$ & additionally $\min_k e^{z_{s,k}}\!\ge\! s_{\min}$ \\
$\|\cdot\|$ & Euclidean ($\ell_2$) norm \\
$\|\cdot\|_\infty$ & $\ell_\infty$ norm \\
$\|\cdot\|_{\mathrm{op}}$ & Operator (spectral) norm \\
\addlinespace[3pt]
\multicolumn{2}{@{}l}{\textbf{Alignment \& mismatch}} \\
\addlinespace[1pt]
$A, B$    & Network labels (distinct from $A_N, B_N$) \\
$\cS_N, \cS_K$ & Symmetric groups on $N$, $K$ elts \\
$\nperm \in \cS_N$ & Neuron perm.\ ($W$-OT) \\
$\tau \in \cS_K$ & Component permutation \\
$P_\tau \in \R^{p\times p}$ & Block permutation matrix \\
$e_\pi \in \R^{p}$ & $(\mathbf{1}_K,\mathbf{0}_{p-K})$; softmax-translation dir. \\
$\R^N_\pi$ & Softmax-translation group; Sec.~\ref{sec:symmetry} \\
$A_N \in \R_+$ & Head-weight mismatch; $\frac{1}{N}\!\sum_i\!\|\Delta v_i\|^2$ \\
$B_N \in \R_+$ & Input-weight mismatch; $\frac{1}{N}\!\sum_i\!\|\Delta w_i\|^2$ \\
$\tilde B_N \in \R_+$ & $\Sigma_x$-weighted input mismatch (App.~\ref{app:additional}) \\
$\Delta_N \in \R_+$ & Endpoint output mismatch \\
$\Delta w_i, \Delta v_i$ & Weight-difference vectors (prefix $\Delta$) \\
$\tilde v_i \in \R^{p}$ & Centred head weight~\eqref{eq:centred_head} \\
$M_{\Delta \tilde V}^\infty$ & Gauge-fixed head mismatch moment~\eqref{eq:head_mismatch_inf} \\
$\mathcal R_N(\nperm,\tau)$ & Residual head mismatch~\eqref{eq:head_residual_mismatch} \\
$Q_N \in \R_+$ & 4th-moment input mismatch; $\frac{1}{N}\sum_i\E_P[(\Delta w_i^\top x)^4]$ \\
\addlinespace[3pt]
\multicolumn{2}{@{}l}{\textbf{Mean-field \& probability}} \\
\addlinespace[1pt]
$\mathfrak{m}_* \in \mathcal{P}(\R^{p+d})$ & Mean-field law; $\mathfrak{m}_* = \mathfrak{m}_T$ \\
$\mathfrak{m}_{*,W} \in \mathcal{P}(\R^d)$ & $W$-marginal of $\mathfrak{m}_*$ \\
$\mathfrak{m}_{*,U} \in \mathcal{P}(\R^d)$ & Whitened marginal of $U\!=\!\Sigma_x^{1/2}W$ (App.~\ref{app:cov_aware}) \\
$z_*(x) \in \R^{p}$ & Limiting output \\
$g:\R^d \to \R^{p}$ & Conditional head map in Prop.~\ref{prop:transfer_struct} \\
$m_{2,*} \in \R_+$ & Head 2nd moment; $\int\!\|V\|^2 d\mathfrak{m}_*$ \\
$m_{4,*} \in \R_+$ & Per-coord.\ head 4th moment \\
$m_{\infty,*} \in \R_+$ & Per-coord.\ limiting moment; $\max_j\!\sqrt{\E[(V)_j^2]}$ \\
$\hat\nu_W^A, \hat\nu_W^B \in \mathcal{P}(\R^d)$ & Empirical $W$-marginals \\
$\hat\nu_U^A, \hat\nu_U^B \in \mathcal{P}(\R^d)$ & Empirical whitened marginals (App.~\ref{app:cov_aware}) \\
$\eta_N \in \R_+$ & $\cW_2$ convergence rate; $O(\kappa_W\sqrt{d/N})$ under \ref{ass:subgauss} \\
$d_{\Sigma} \in \R_+$ & Intrinsic dimension of the whitened law (Cor.~\ref{cor:sigma_ot}) \\
$\varepsilon_N(\delta) \in \R_+$ & Chebyshev correction; $\sqrt{6p\,m_{4,*}/(N\delta)}$ \\
$\cW_2$  & 2-Wasserstein distance \\
$P$      & Data distribution over $(x,y)$ \\
$\E_P$   & Expectation under $P$ \\
$\E$     & Expectation over network params \\
$\Prob$  & Probability over network params \\
$\delta \in (0,1)$ & Failure probability (also $\delta\!\in\!\R^N$ for softmax translation in Sec.~\ref{sec:symmetry}) \\
\addlinespace[3pt]
\multicolumn{2}{@{}l}{\textbf{Barrier-rate constants}} \\
\addlinespace[1pt]
$C_{\mathrm{tr}}(\delta,N) \in \R_+$ & Transfer-term prefactor \\
$C_{\mathrm{sc}}(\delta,N) \in \R_+$ & Semi-convexity prefactor \\
$\hat M_V(\delta,N) \in \R_+$ & High-prob.\ head moment; $(m_{\infty,*}^2\!+\!\varepsilon_N)^{1/2}$ \\
$\xi \in \R_+$ & Target barrier threshold \\
\addlinespace[3pt]
\multicolumn{2}{@{}l}{\textbf{vMF-specific quantities}} \\
\addlinespace[1pt]
$d_y \in \N$ & Ambient dim.\ of sphere $S^{d_y-1}$ \\
$\eta_k \in \R^{d_y}$ & Natural parameter; $\kappa_k\mu_k$ \\
$\kappa_k \in \R_+$ & Concentration; $\|\eta_k\|$ \\
$\mu_k \in S^{d_y-1}$ & Mean direction \\
$C_{d_y}(\kappa)$ & vMF normalising constant \\
$I_\nu(\kappa)$ & Modified Bessel fn.\ of first kind \\
$A_{d_y}(\kappa)$ & Mean resultant length; $I_{d_y/2}/I_{d_y/2-1}$ \\
\addlinespace[3pt]
\multicolumn{2}{@{}l}{\textbf{Asymptotic notation}} \\
\addlinespace[1pt]
$O(\cdot), \Omega(\cdot)$ & Standard asymptotic bounds \\
$O_T(\cdot)$ & Const.\ depends on $T$, not $N,K$ \\
\addlinespace[3pt]
\multicolumn{2}{@{}l}{\textbf{Appendix-only symbols}} \\
\addlinespace[1pt]
$\alpha \in \R_+$ & Learning rate (App.~\ref{app:mf_sgd}) \\
$b \in \N$ & Mini-batch size (App.~\ref{app:mf_sgd}) \\
$s \in [0,T]$ & Training time (distinct from $s_k$) \\
\bottomrule
\end{tabular}
\end{multicols}
\vspace{-6pt}
\end{table}




%======================================================================
\section{Proofs of Main Results}\label{app:proofs}
%======================================================================

\subsection{Proof of Proposition~\ref{prop:loss_regularity}}\label{sec:proof_loss_regularity}

\begin{proof}
The proof proceeds in two parts: the Lipschitz bound and the
semi-convexity bound. Throughout, we use $s_k \ge s_{\min}$ in place of the
crude bound $s_k \ge e^{-R}$.

\paragraph{Part 1: Lipschitz bound.}
Writing $\gamma_k$ for the responsibilities and
$\psi_k = (y-\mu_k)^2/s_k^2$, the gradient decomposes into three
orthogonal blocks:
\begin{align*}
  \partial\cL/\partial z_{\pi,k} &= -(\gamma_k - \pi_k), \\
  \partial\cL/\partial z_{\mu,k} &= -\gamma_k (y-\mu_k)/s_k^2, \\
  \partial\cL/\partial z_{s,k}   &= -\gamma_k(\psi_k - 1).
\end{align*}
Bounding each block's squared norm:
\begin{itemize}
  \item \textbf{$z_\pi$-block:}
    Since $\gamma_k, \pi_k \ge 0$ for all $k$, the inner product
    $\gamma \cdot \pi = \sum_k \gamma_k \pi_k \ge 0$, so
    $\|\gamma - \pi\|_2^2 = \|\gamma\|_2^2 - 2\gamma\cdot\pi + \|\pi\|_2^2
    \le \|\gamma\|_2^2 + \|\pi\|_2^2 \le \|\gamma\|_1 + \|\pi\|_1 = 2$,
    using $\|\cdot\|_2^2 \le \|\cdot\|_1$ for probability vectors.
    This bound does not involve $K$.
  \item \textbf{$z_\mu$-block:}
    $\textstyle\sum_k \gamma_k^2 (y-\mu_k)^2/s_k^4
    \le \Lambda^2/s_{\min}^4\cdot\textstyle\sum_k \gamma_k^2
    \le \Lambda^2/s_{\min}^4$,
    using $\sum_k \gamma_k^2 \le \sum_k \gamma_k = 1$
    (since $0 \le \gamma_k \le 1$ implies $\gamma_k^2 \le \gamma_k$).
  \item \textbf{$z_s$-block:}
    $\textstyle\sum_k \gamma_k^2(\psi_k-1)^2
    \le (\Lambda^2/s_{\min}^2+1)^2\cdot\textstyle\sum_k\gamma_k^2
    \le (\Lambda^2/s_{\min}^2+1)^2$.
\end{itemize}
Hence $\|\nabla\cL\|^2 = \sum_{k}\bigl[(\gamma_k-\pi_k)^2
+ \gamma_k^2(y-\mu_k)^2/s_k^4
+ \gamma_k^2(\psi_k-1)^2\bigr]$ satisfies
\begin{equation}\label{eq:grad_sq_bound}
  \|\nabla \cL\|^2
  \le 2 + \Lambda^2/s_{\min}^4
  + \bigl(\Lambda^2/s_{\min}^2 + 1\bigr)^2,
\end{equation}
and $C_L = O\!\bigl(\Lambda^2/s_{\min}^2\bigr)$.

\textbf{$\ell_\infty$-Lipschitz constant.}
For $C_L^\infty = \|\nabla\cL\|_1 = \sum_k (|\gamma_k-\pi_k|
+ \gamma_k|y-\mu_k|/s_k^2 + \gamma_k|\psi_k-1|)$, we bound each
block's $\ell_1$ norm.
The $z_\pi$-block: $\sum_k |\gamma_k - \pi_k| \le \sum_k(\gamma_k+\pi_k) = 2$.
The $z_\mu$-block:
$\sum_k \gamma_k|y-\mu_k|/s_k^2
\le \Lambda/s_{\min}^2\sum_k \gamma_k = \Lambda/s_{\min}^2$.
The $z_s$-block:
$\sum_k \gamma_k|\psi_k-1|
\le (\Lambda^2/s_{\min}^2+1)\sum_k \gamma_k
= \Lambda^2/s_{\min}^2+1$.
Summing gives $C_L^\infty \le 3 + \Lambda/s_{\min}^2
+ \Lambda^2/s_{\min}^2 = O\!\bigl(\Lambda^2/s_{\min}^2\bigr)$.

\paragraph{Part 2: Semi-convexity bound.}
We need $\lambda_{\min}(\nabla^2 \cL_{\mathrm{MDN}}) \ge -\rho$.
By~\eqref{eq:hess_decomp} and Weyl's inequality:
\begin{equation}\label{eq:weyl_split}
  \lambda_{\min}(\nabla^2 \cL_{\mathrm{MDN}})
  \ge \lambda_{\min}\!\Bigl(-\sum_k \gamma_k \nabla^2 u_k\Bigr)
  - \|\mathrm{Cov}_\gamma(\nabla u_k)\|_{\mathrm{op}}.
\end{equation}

We bound each term in~\eqref{eq:weyl_split} separately.
For the weighted Hessian $-\sum_k \gamma_k \nabla^2 u_k$,
since $u_k = \log \pi_k + \log \varphi_k$, where $\log\pi_k$ depends
only on $z_\pi$ and $\log\varphi_k$ depends only on $(z_{\mu,k},z_{s,k})$,
all cross-derivatives $\partial^2 u_k/\partial z_{\pi,j}\,\partial z_{\mu,k} = 0$
and $\partial^2 u_k/\partial z_{\pi,j}\,\partial z_{s,k} = 0$.
Hence $\nabla^2 u_k$ is block-diagonal:
\begin{itemize}[leftmargin=*,itemsep=0pt]
\item \emph{$z_\pi$-block}: $(\nabla^2 u_k)_{z_\pi}
  = \nabla^2_{z_\pi}\log\pi_k = -C_\pi$
  where $C_\pi = \mathrm{diag}(\pi)-\pi\pi^\top$ is the softmax
  covariance matrix (PSD since
  $v^\top C_\pi v = \mathrm{Var}_\pi(v) \ge 0$; same for all $k$
  because $\nabla^2_{z_\pi}\log\pi_k$ is independent of~$k$).
  Its eigenvalues lie in $[0,1]$ (each $\pi_k(1-\pi_k) \le 1/4$).
  Therefore $(-\sum_k\gamma_k\nabla^2 u_k)_{z_\pi} = +C_\pi \succeq 0$.
\item \emph{$(z_{\mu,k},z_{s,k})$-block}: Direct computation gives
  \[
    (\nabla^2 u_k)_{(z_{\mu,k},z_{s,k})}
    = -\frac{1}{s_k^2} Q(a_k),
    \quad a_k \coloneqq y - \mu_k,\quad
    Q(a) \coloneqq \begin{pmatrix} 1 & 2a \\ 2a & 2a^2 \end{pmatrix}.
  \]
  For $j \ne k$ the $(z_{\mu,j},z_{s,j})$-block of $\nabla^2 u_k$ is zero.
  Hence the $(z_{\mu,j},z_{s,j})$-block of $-\sum_k\gamma_k\nabla^2 u_k$
  equals $\dfrac{\gamma_j}{s_j^2} Q(a_j)$.
\end{itemize}

The key algebraic fact is that $\lambda_{\min}(Q(a)) \ge -1$ for all $a \in \R$.
The characteristic polynomial gives
$\lambda_{\min}(Q(a)) = \frac{(1+2a^2)-\sqrt{(1+2a^2)^2+8a^2}}{2}$,
and $\lambda_{\min}(Q(a)) \ge -1$ is equivalent to
$(3+2a^2) \ge \sqrt{(1+2a^2)^2+8a^2}$.
Squaring (both sides positive):
$(3+2a^2)^2 = 9+12a^2+4a^4 \ge 1+12a^2+4a^4 = (1+2a^2)^2+8a^2$.
This reduces to $9 \ge 1$.

It follows from $\lambda_{\min}(Q(a)) \ge -1$ that
\[
  \lambda_{\min}\!\Bigl(\frac{\gamma_j}{s_j^2} Q(a_j)\Bigr)
  = \frac{\gamma_j}{s_j^2}\lambda_{\min}(Q(a_j))
  \ge -\frac{\gamma_j}{s_j^2}
  \ge -\frac{1}{s_j^2}
  \ge -1/s_{\min}^2.
\]
Since the $z_\pi$-block contributes non-negatively,
\begin{equation}\label{eq:step1_bound}
  \lambda_{\min}\!\Bigl(-\sum_k \gamma_k \nabla^2 u_k\Bigr) \;\ge\; -1/s_{\min}^2.
\end{equation}

It remains to bound the covariance term
$\|\mathrm{Cov}_\gamma(\nabla u_k)\|_{\mathrm{op}}$.
Writing $g_k = \nabla u_k$, for any unit vector $\omega$:
\begin{align*}
  \omega^\top \mathrm{Cov}_\gamma(g)\, \omega
  &= \mathrm{Var}_\gamma(\omega^\top g_k)
   = \E_\gamma[(\omega^\top g_k)^2] - (\E_\gamma[\omega^\top g_k])^2
   \le \E_\gamma[(\omega^\top g_k)^2] \\
  &\le \E_\gamma[\|g_k\|^2]
   \le \max_k \|g_k\|^2.
\end{align*}
Here the first inequality uses $\mathrm{Var}(X) \le \E[X^2]$, and the
second uses the Cauchy--Schwarz bound $(\omega^\top g_k)^2 \le \|\omega\|^2\|g_k\|^2
= \|g_k\|^2$. Taking the supremum over $\omega$:
\[
  \|\mathrm{Cov}_\gamma(g)\|_{\mathrm{op}}
  \le \max_k \|g_k\|^2.
\]
Each $\|g_k\|^2 = \|\nabla u_k\|^2$ is bounded by the same
per-component argument as in Part~1
(where the same three blocks appear with identical bounds):
the $z_\pi$-component contributes $\|e_k-\pi\|^2 \le 2$;
the $z_{\mu,k}$-component contributes $(y-\mu_k)^2/s_k^4 \le \Lambda^2/s_{\min}^4$;
the $z_{s,k}$-component contributes $(\psi_k-1)^2 \le (\Lambda^2/s_{\min}^2+1)^2$.
Hence $\|\nabla u_k\|^2 \le G^2$ where
$G^2 \coloneqq 2 + \Lambda^2/s_{\min}^4 + (\Lambda^2/s_{\min}^2+1)^2
= O(\Lambda^4/s_{\min}^4)$.

A tighter bound uses the pairwise-variance identity:
\begin{equation}\label{eq:pairwise_identity}
  \|\mathrm{Cov}_\gamma(g)\|_{\mathrm{op}}
  \;\le\; 2\bar\chi\,\max_k\|g_k\|^2,
  \qquad
  \bar\chi \coloneqq \sum_k\gamma_k(1-\gamma_k)\in[0,(K{-}1)/K],
\end{equation}
which follows from the representation
$\mathrm{Cov}_\gamma(g)
= \sum_{j<k}\gamma_j\gamma_k(g_j-g_k)(g_j-g_k)^\top$
and $\sum_{j<k}\gamma_j\gamma_k = \bar\chi/2$.
Taking the universal bound $\bar\chi \le 1$, the covariance
contribution is at most $2G^2$.

Combining~\eqref{eq:step1_bound} and the covariance bound via~\eqref{eq:weyl_split}:
\begin{equation}\label{eq:rho_bound}
  \rho \;\le\; 1/s_{\min}^2 + 2G^2
  \;=\; 1/s_{\min}^2 + 2\bigl(2 + \Lambda^2/s_{\min}^4
  + (\Lambda^2/s_{\min}^2+1)^2\bigr).
\end{equation}
\end{proof}

\subsection{Proof of Proposition~\ref{prop:symmetry}
  (Symmetry characterisation)}\label{app:proof_symmetry}

We first state the genericity condition referenced in the main text.

\begin{definition}[Generic position]\label{def:generic}
Two parameter sets $\{(w_i, v_i)\}_{i=1}^N$ and
$\{(w_i', v_i')\}_{i=1}^{N'}$ are in \emph{generic position}
if (i)~the input weights of the \emph{union}
$\{w_1,\ldots,w_N,w_1',\ldots,w_{N'}'\}$
are pairwise non-proportional:
$w \ne \alpha\, w'$ for every distinct pair $(w,w')$ in the union
and every $\alpha \in \R$
(this includes both within-network and cross-network pairs),
and (ii)~the Gaussian components induced by distinct component indices
are distinct:
$(\mu_k(x), s_k(x)) \ne (\mu_\ell(x), s_\ell(x))$ for $k \ne \ell$
on a set of inputs of positive $P$-measure.
\end{definition}

When the initialisation law $\mathfrak{m}_0$ is absolutely continuous
(e.g.\ Gaussian, He, or Xavier), generic position holds almost
surely for the i.i.d.\ draws from $\mathfrak{m}_*$:
condition~(i) because the input weights are
absolutely continuous, both within each network and across the two
networks drawn independently from the same law,
and condition~(ii) because distinct head-weight
components produce distinct ridge-function combinations
$P$-almost everywhere.

\begin{proof}[Proof of Proposition~\ref{prop:symmetry}]
Suppose $\{(w_i, v_i)\}_{i=1}^N$ and $\{(w_i', v_i')\}_{i=1}^{N'}$
are in generic position and define the same conditional density
$p(y \mid x)$ for $P$-a.e.\ $x$.
We proceed in two stages.

\medskip\noindent\textbf{Stage 1: Component permutation via Gaussian
mixture identifiability.}
For $P$-a.e.\ $x$, the two parameter sets produce two finite Gaussian
mixtures with the same density:
\[
  \sum_{k=1}^K \pi_k(x)\,\cN(y;\,\mu_k(x),s_k(x)^2)
  = \sum_{k=1}^K \pi_k'(x)\,\cN(y;\,\mu_k'(x),s_k'(x)^2)
  \qquad \forall\,y.
\]
By the identifiability theorem for finite Gaussian
mixtures~\cite{yakowitz1968}, the two mixtures have the same number
of components ($K = K$) and there exists a permutation
$\tau_x \in \cS_K$ such that
$\pi_k(x) = \pi_{\tau_x(k)}'(x)$,
$\mu_k(x) = \mu_{\tau_x(k)}'(x)$,
$s_k(x) = s_{\tau_x(k)}'(x)$.
Since the generic-position assumption ensures that the Gaussian
component functions $(\mu_k(\cdot), s_k(\cdot))$ are distinct as
functions of $x$, the permutation $\tau_x$ is independent of $x$;
call it $\tau \in \cS_K$.
Hence, for each component index $k$ and $P$-a.e.\ $x$:
\begin{equation}\label{eq:component_match}
  \mu_k(x) = \mu_{\tau(k)}'(x), \qquad
  s_k(x) = s_{\tau(k)}'(x), \qquad
  \pi_k(x) = \pi_{\tau(k)}'(x).
\end{equation}

\medskip\noindent\textbf{Stage 2: Neuron matching via ridge-function
linear independence.}
The $\mu$ and $\log s$ equalities in~\eqref{eq:component_match}
give, for each component $k$:
\[
  \frac{1}{N}\sum_{i=1}^N (v_i)_{\mu,k}\,\sigma(w_i^\top x)
  = \frac{1}{N'}\sum_{i=1}^{N'} (v_i')_{\mu,\tau(k)}\,\sigma(w_i'{}^\top x)
  \qquad P\text{-a.e.}
\]
(and analogously for the $\log s$ block).
Combining all $2K$ such equalities (over both $\mu$ and $\log s$
blocks) yields a linear relation among the ridge functions
$\{\sigma(w_i^\top \cdot)\}_{i=1}^N \cup
\{\sigma(w_j'{}^\top \cdot)\}_{j=1}^{N'}$.
Since $\sigma$ is continuous and non-polynomial (bounded polynomials
on $\R$ are constant) and $d \ge 2$, the ridge-function
linear-independence theorem~\cite{pinkus1999} implies that functions
$\sigma(w^\top x)$ with pairwise non-proportional weight vectors are
linearly independent over $L^2(P)$.
By the pairwise non-proportionality in the generic-position
assumption, the only way the relation can hold is if $N' = N$ and
there exists a permutation $\nperm \in \cS_N$ with
$w_i' = w_{\nperm(i)}$
and the head-weight coordinates in the $\mu$ and $\log s$ blocks
satisfy
$(v_i')_{\mu,\tau(k)} = (v_{\nperm(i)})_{\mu,k}$,
$(v_i')_{s,\tau(k)} = (v_{\nperm(i)})_{s,k}$
for all $i, k$.

It remains to determine the $\pi$-block.
The equality $\pi_k(x) = \pi_{\tau(k)}'(x)$ means
\[
  \mathrm{softmax}\!\Bigl(
    \frac{1}{N}\sum_i (v_i)_{\pi}\,\sigma(w_i^\top x)
  \Bigr)_k
  = \mathrm{softmax}\!\Bigl(
    \frac{1}{N}\sum_i (v_i')_{\pi}\,\sigma(w_{\nperm(i)}^\top x)
  \Bigr)_{\tau(k)}.
\]
Since softmax is injective up to a common additive
constant~\cite{luce1959}, the $\pi$-logit vectors satisfy
\[
  \frac{1}{N}\sum_i (v_i)_{\pi,k}\,\sigma(w_i^\top x)
  = \frac{1}{N}\sum_i (v_i')_{\pi,\tau(k)}\,\sigma(w_{\nperm(i)}^\top x)
  + c(x)
  \qquad \forall\,k,
\]
where $c(x)$ is the same for all $k$.
Applying ridge-function linear independence once more gives
$(v_i')_{\pi,\tau(k)} = (v_{\nperm(i)})_{\pi,k} + \delta_{\nperm(i)}$
for some constants $\delta_{\nperm(i)} \in \R$ independent of $k$.
Assembling all three blocks, $v_i' = P_\tau v_{\nperm(i)}
+ \delta_{\nperm(i)}\,e_\pi$.
\end{proof}

\subsection{Proof of Proposition~\ref{prop:barrier_neutrality}
  (Barrier neutrality)}\label{app:proof_neutrality}

\begin{proof}
The interpolated head weights are
$v_i^{t'} = (1-t)\,v_i^A + t\,(v_i^B + \delta_i\,e_\pi)
= v_i^t + t\,\delta_i\,e_\pi$,
while $w_i^{t'} = w_i^t$.
Hence the $\pi$-logits satisfy
\[
  z_\pi^{t'}(x)
  = z_\pi^t(x)
    + t\,\underbrace{\frac{1}{N}\sum_{i=1}^N
      \delta_i\,\sigma(w_i^{t\top} x)}_{h(x)}
    \cdot\mathbf{1}_K.
\]
Because $h(x)\cdot\mathbf{1}_K$ is a common-mode shift,
$\mathrm{softmax}(z_\pi^{t'}) = \mathrm{softmax}(z_\pi^t)$,
so the mixing weights coincide at every $t$.
The $\mu$ and $\log s$ components are unchanged.
Therefore the conditional density---and hence the loss---is identical
along the two paths.
\end{proof}

\subsection{Proof of Proposition~\ref{prop:transfer}}\label{sec:transfer}
\begin{proof}
Let $w_i(t) = (1-t)w_{A,i} + tw_{B,\nperm(i)}$ and set
$\Delta w_i = w_{A,i} - w_{B,\nperm(i)}$.  The key decomposition is
\begin{align*}
  z_{M_t}(x) - \hat z_t(x)
  &= \frac{1}{N}\sum_i \Bigl[
    (1-t)v_{A,i}\bigl(\sigma(w_i(t)^\top x) - \sigma(w_{A,i}^{\top}x)\bigr) \\
  &\qquad\qquad
    + tP_\tau v_{B,\nperm(i)}\bigl(\sigma(w_i(t)^\top x)
      - \sigma(w_{B,\nperm(i)}^{\top}x)\bigr)
  \Bigr].
\end{align*}
Since $w_i(t) - w_{A,i} = -t\Delta w_i$ and
$w_i(t) - w_{B,\nperm(i)} = (1-t)\Delta w_i$, the Lipschitz bound on
$\sigma$ (Assumption~\ref{ass:activ}) gives
\[
  |\sigma(w_i(t)^\top x) - \sigma(w_{A,i}^{\top}x)|
  \le L_\sigma t\,|\Delta w_i^\top x|, \quad
  |\sigma(w_i(t)^\top x) - \sigma(w_{B,\nperm(i)}^{\top}x)|
  \le L_\sigma(1-t)\,|\Delta w_i^\top x|.
\]
Taking the $\ell_\infty$ norm coordinate-by-coordinate: for each
coordinate~$j$,
\[
  |(z_{M_t}(x) - \hat z_t(x))_j|
  \le \frac{t(1-t)L_\sigma}{N}
    \sum_i \bigl(|(v_{A,i})_j| + |(P_\tau v_{B,\nperm(i)})_j|\bigr)|\Delta w_i^\top x|.
\]
Applying Cauchy--Schwarz twice ($|\Delta w_i^\top x| \le \|\Delta w_i\|\,\|x\|$,
then over~$i$):
\[
  \frac{1}{N}\sum_i |(v_{A,i})_j|\,|\Delta w_i^\top x|
  \le \|x\|\Bigl(\frac{1}{N}\sum_i (v_{A,i})_j^2\Bigr)^{1/2}
    \Bigl(\frac{1}{N}\sum_i\|\Delta w_i\|^2\Bigr)^{1/2}
  \le M_V^\infty\,\|x\|\sqrt{B_N},
\]
using~\eqref{eq:head_moment_inf}.
The same bound holds for the $P_\tau v_{B,\nperm(i)}$ term, giving
\[
  \|z_{M_t}(x) - \hat z_t(x)\|_\infty
  \le 2t(1-t)\,L_\sigma M_V^\infty\,\|x\|\sqrt{B_N}.
\]
Squaring, taking $\E_P$ (using $\E_P[\|x\|^2] \le M_x^2$),
and $\sup_t 4t^2(1-t)^2 = 1/4$, yields~\eqref{eq:transfer}.
\end{proof}

\subsection{Proof of Proposition~\ref{prop:transfer_refined}}
\label{sec:proof_transfer_refined}

\begin{proof}
We use the following \emph{exact} decomposition of the transfer gap.
Writing $v_i(t) = (1-t)v_{A,i} + tP_\tau v_{B,\nperm(i)}$,
$\Delta v_i = v_{A,i} - P_\tau v_{B,\nperm(i)}$,
and $\Delta\sigma_i(x) = \sigma(w_{B,\nperm(i)}^\top x)
- \sigma(w_{A,i}^\top x)$, one verifies algebraically that
\begin{equation}\label{eq:exact_decomp}
  z_{M_t}(x) - \hat z_t(x)
  = \frac{1}{N}\sum_{i=1}^N \Bigl[
    v_i(t)\,\mathcal{C}_i(x,t)
    + t(1-t)\,\Delta v_i\,\Delta\sigma_i(x)
  \Bigr],
\end{equation}
where
\[
  \mathcal{C}_i(x,t) \;\coloneqq\;
  \sigma(w_i(t)^\top x)
  - \bigl[(1-t)\sigma(w_{A,i}^\top x) + t\sigma(w_{B,\nperm(i)}^\top x)\bigr]
\]
is the \emph{curvature gap} of $\sigma$ at the interpolated point.
For linear $\sigma$, $\mathcal{C}_i \equiv 0$ and only the cross-term remains.

\paragraph{Cross-term.}
By the Lipschitz bound (Assumption~\ref{ass:activ}\,(ii)),
$|\Delta\sigma_i(x)| \le L_\sigma |\Delta w_i^\top x|$.
For each coordinate~$j$:
\[
  \Bigl|\frac{1}{N}\sum_i (\Delta v_i)_j \Delta\sigma_i(x)\Bigr|
  \le \frac{L_\sigma}{N}\sum_i |(\Delta v_i)_j|\,|\Delta w_i^\top x|
  \le L_\sigma\,M_{\Delta \tilde V}^\infty\,\sqrt{\tilde B_N(x)},
\]
by Cauchy--Schwarz with~\eqref{eq:head_mismatch_inf}.
Squaring, taking $\E_P$, and using $\E_P[\tilde B_N(x)] = \tilde B_N$
gives the cross-term contribution
$t(1-t) L_\sigma M_{\Delta \tilde V}^\infty \sqrt{\tilde B_N}$.

\paragraph{Curvature term.}
Since $\sigma \in C^2$ (Assumption~\ref{ass:activ}\,(iii)),
the convexity-gap bound gives
\[
  |\mathcal{C}_i(x,t)|
  \;\le\; \tfrac{1}{2}\,L_{\sigma'}\,t(1-t)\,(\Delta w_i^\top x)^2.
\]
(This follows from writing $\mathcal{C}_i$ as
$t[\sigma(w_i(t)^\top x) - \sigma(w_{B}^\top x)]
+ (1-t)[\sigma(w_i(t)^\top x) - \sigma(w_{A}^\top x)]$
and applying the $C^2$ mean-value remainder to each bracket.)
For each coordinate~$j$:
\[
  \Bigl|\frac{1}{N}\sum_i (v_i(t))_j\,\mathcal{C}_i(x,t)\Bigr|
  \le \frac{L_{\sigma'}}{2}\,t(1-t)\,
  \frac{1}{N}\sum_i |(v_i(t))_j|\,(\Delta w_i^\top x)^2.
\]
Applying Cauchy--Schwarz over~$i$:
\[
  \frac{1}{N}\sum_i |(v_i(t))_j|\,(\Delta w_i^\top x)^2
  \le M_V^\infty\,\sqrt{Q_N(x)},
\]
where
$Q_N(x) = \frac{1}{N}\sum_i (\Delta w_i^\top x)^4$.
Taking $\E_P$ of $Q_N(x)$ gives
$Q_N = \frac{1}{N}\sum_i \E_P[(\Delta w_i^\top x)^4]$,
using the fourth-moment condition $\E_P[\|x\|^4] \le M_{x,4}^4$.

\paragraph{Assembly.}
Combining both terms in~\eqref{eq:exact_decomp},
taking the $\ell_\infty$ norm, squaring, and taking $\E_P$:
\[
  \bigl(\E_P[\|z_{M_t} - \hat z_t\|_\infty^2]\bigr)^{1/2}
  \le t(1-t)\bigl[
    L_\sigma\,M_{\Delta \tilde V}^\infty\sqrt{\tilde B_N}
    + \tfrac{1}{2}\,L_{\sigma'}\,M_V^\infty\sqrt{Q_N}
  \bigr].
\]
Since $\sup_t t(1-t) = 1/4$,
this yields~\eqref{eq:transfer_refined}.
\end{proof}

\subsection{Proof of Corollary~\ref{cor:sigma_ot}}
\begin{proof}
By definition of~\eqref{eq:OT_perm_sigma},
\[
  \tilde B_N
  = \frac{1}{N}\sum_{i=1}^N
    \bigl\|\Sigma_x^{1/2}(w_{A,i} - w_{B,\nperm_{\Sigma}(i)})\bigr\|^2
  = \frac{1}{N}\sum_{i=1}^N
    \|u_{A,i} - u_{B,\nperm_{\Sigma}(i)}\|^2,
\]
where $u_{X,i} = \Sigma_x^{1/2} w_{X,i}$.
This is precisely the empirical $\cW_2^2$ cost between
$\hat\nu_U^A$ and $\hat\nu_U^B$, so $\nperm_{\Sigma}$ minimises
$\tilde B_N$ among all hidden permutations.
The refined barrier bounds in this appendix depend on the hidden
alignment through $\tilde B_N$
(Propositions~\ref{prop:transfer_refined} and~\ref{prop:dir_barrier}),
so any high-probability rate for the whitened empirical measures
substitutes directly for the Euclidean $W$-marginal rate.
(The main-body bounds use $B_N$, for which
$\tilde B_N \le \lambda_{\max}(\Sigma_x)\,B_N$.)
\end{proof}

\subsection{Proof of Proposition~\ref{prop:transfer_struct}}
\begin{proof}
For each coordinate $j$,
\[
  \bigl(\tilde v_{A,i} - P_\tau \tilde v_{B,\nperm(i)}\bigr)_j
  =
  \bigl(g_j(w_{A,i}) - g_j(w_{B,\nperm(i)})\bigr)
  + \bigl(r_{A,i} - P_\tau r_{B,\nperm(i)}\bigr)_j.
\]
Taking the empirical $L^2$ norm over $i$ and applying Minkowski's
inequality gives
\[
  \Bigl(
    \frac{1}{N}\sum_{i=1}^N
    \bigl(\tilde v_{A,i} - P_\tau \tilde v_{B,\nperm(i)}\bigr)_j^2
  \Bigr)^{1/2}
  \le
  \Bigl(
    \frac{1}{N}\sum_{i=1}^N
    \bigl(g_j(w_{A,i}) - g_j(w_{B,\nperm(i)})\bigr)^2
  \Bigr)^{1/2}
  + \mathcal R_N(\nperm,\tau).
\]
Maximising over $j$ and using the Lipschitz condition
~\eqref{eq:head_reg_lip},
\[
  M_{\Delta \tilde V}^\infty
  \le
  L_g \Bigl(\frac{1}{N}\sum_{i=1}^N
    \|w_{A,i} - w_{B,\nperm(i)}\|^2\Bigr)^{1/2}
  + \mathcal R_N(\nperm,\tau)
  = L_g \sqrt{B_N} + \mathcal R_N(\nperm,\tau),
\]
which is~\eqref{eq:head_mismatch_struct}.
Substituting into Proposition~\ref{prop:transfer_refined}
(Appendix~\ref{app:refined_coupling}),
the transfer term becomes
$O(L_\sigma(L_g\sqrt{B_N}+\mathcal R_N)\sqrt{\tilde B_N}
+ L_{\sigma'} M_V^\infty \sqrt{Q_N})$,
with $\tilde B_N \le \lambda_{\max}(\Sigma_x) B_N$.
If $\mathcal R_N(\nperm,\tau) = O(\sqrt{B_N})$ and
$Q_N = O(B_N^2)$, both terms are $O(B_N)$, proving the claim.
\end{proof}

\subsection{Proofs of Theorems~\ref{thm:general} and~\ref{thm:main}}\label{sec:proof_main}

Both proofs use two auxiliary lemmas
(domain containment and endpoint-mismatch control)
together with Proposition~\ref{prop:transfer}.

\begin{lemma}[Path-domain containment]\label{lem:auto_path_domain}
Under Assumptions~\ref{ass:arch} and~\ref{ass:activ},
let $M_V^\infty$ be the per-coordinate head moment~\eqref{eq:head_moment_inf}
and set $R = S_\sigma M_V^\infty$.
Then
$\|z_A(x)\|_\infty,\;\|P_\tau z_B(x)\|_\infty,\;
\|z_{M_t}(x)\|_\infty,\;\|\hat z_t(x)\|_\infty \le R$ for all
$x$, $t \in [0,1]$, and $(\nperm, \tau) \in \cS_N \times \cS_K$.
\end{lemma}
\begin{proof}
For any output $z(x) = \frac{1}{N}\sum_i v_i \sigma(w_i^\top x)$ with
$v_i \in \R^{p}$, each output coordinate $j$ satisfies
\[
  |z_j(x)|
  = \Bigl|\frac{1}{N}\sum_i (v_i)_j\,\sigma(w_i^\top x)\Bigr|
  \le S_\sigma\,\frac{1}{N}\sum_i |(v_i)_j|
  \le S_\sigma\!\Bigl(\frac{1}{N}\sum_i (v_i)_j^2\Bigr)^{\!1/2}
  \le S_\sigma M_V^\infty,
\]
using $|\sigma|\le S_\sigma$, Jensen's inequality
$(\frac{1}{N}\sum_i |a_i|)^2 \le \frac{1}{N}\sum_i a_i^2$
(concavity of $\sqrt{\cdot}$ applied to the empirical mean), and the definition of
$M_V^\infty$~\eqref{eq:head_moment_inf}.
Hence $\|z(x)\|_\infty \le S_\sigma M_V^\infty = R$.

Applying this to $v_i = v_{A,i}$ gives $\|z_A(x)\|_\infty \le R$;
to $v_i = P_\tau v_{B,\nperm(i)}$ gives $\|P_\tau z_B(x)\|_\infty \le R$.
For the interpolated network, the time-dependent weights
$v_i(t) = (1-t)v_{A,i} + tP_\tau v_{B,\nperm(i)}$ satisfy
$\max_j \frac{1}{N}\sum_i (v_i(t))_j^2 \le (M_V^\infty)^2$ by convexity
of $(\cdot)^2$ and~\eqref{eq:head_moment_inf}, giving
$\|z_{M_t}(x)\|_\infty \le R$.
Finally, $\hat z_t(x) = (1-t)z_A(x)+tP_\tau z_B(x)$ satisfies
$\|\hat z_t(x)\|_\infty \le (1-t)\|z_A(x)\|_\infty+t\|P_\tau z_B(x)\|_\infty
\le R$ by convexity of $\cD_R$ (an $\ell^\infty$ ball is convex).
\end{proof}

\begin{lemma}[Endpoint output mismatch]\label{lem:delta_control}
The endpoint mismatch
$\Delta_N = \E_P[\|z_A(x) - P_\tau z_B(x)\|^2]$ satisfies
\[
  \Delta_N \;\le\; 2 S_\sigma^2 A_N
  + 2 L_\sigma^2 M_x^2 \bar M_{V,N}^2 B_N,
\]
where $A_N = \frac{1}{N}\sum_i \|v_{A,i} - P_\tau v_{B,\nperm(i)}\|^2$
and $\bar M_{V,N}^2 = \frac{1}{N}\sum_i \|P_\tau v_{B,\nperm(i)}\|^2$.
\end{lemma}
\begin{proof}
Write
$z_A(x) - P_\tau z_B(x)
= \frac{1}{N}\sum_i [v_{A,i} \sigma(w_{A,i}^\top x)
- P_\tau v_{B,\nperm(i)} \sigma(w_{B,\nperm(i)}^\top x)]$.
Add and subtract $P_\tau v_{B,\nperm(i)} \sigma(w_{A,i}^\top x)$:
\[
  z_A - P_\tau z_B
  = \frac{1}{N}\sum_i \Delta v_i \cdot \sigma(w_{A,i}^\top x)
  + \frac{1}{N}\sum_i P_\tau v_{B,\nperm(i)}
  [\sigma(w_{A,i}^\top x) - \sigma(w_{B,\nperm(i)}^\top x)],
\]
where $\Delta v_i = v_{A,i} - P_\tau v_{B,\nperm(i)}$. Taking the squared
$\ell^2$ norm (using $(a+b)^2 \le 2a^2 + 2b^2$) and applying
Cauchy--Schwarz over $i$: the first term gives
$\le S_\sigma^2 \cdot \frac{1}{N}\sum_i \|\Delta v_i\|^2 = S_\sigma^2 A_N$;
the second term gives (using Lipschitz continuity of $\sigma$)
$\le L_\sigma^2 \|x\|^2 (\frac{1}{N}\sum_i \|P_\tau v_{B,\nperm(i)}\|^2)
(\frac{1}{N}\sum_i \|\Delta w_i\|^2)
\le L_\sigma^2 \|x\|^2 \bar M_{V,N}^2 B_N$.
Taking expectations (using $\E_P\|x\|^2 \le M_x^2$) and the factor of~2:
\begin{equation}\label{eq:delta_bound}
  \Delta_N \le 2 S_\sigma^2 A_N
  + 2 L_\sigma^2 M_x^2 \bar M_{V,N}^2 B_N.
\end{equation}
\end{proof}

\begin{proof}[Proof of Theorem~\ref{thm:general}]
By Lemma~\ref{lem:auto_path_domain}, all outputs
$z_A$, $P_\tau z_B$, $z_{M_t}$, $\hat z_t$ lie in $\cD_R$.
Proposition~\ref{prop:transfer} gives
$(\E_P[\|z_{M_t} - \hat z_t\|_\infty^2])^{1/2}
\le \tfrac{1}{2} L_\sigma M_V^\infty M_x \sqrt{B_N}$.
Since $\ell$ is $C_L^\infty$-Lipschitz on $\cD_R$:
\[
  \E_P[\ell(z_{M_t}, y)]
  \le \E_P[\ell(\hat z_t, y)]
  + \tfrac{1}{2}\,C_L^\infty L_\sigma M_V^\infty M_x \sqrt{B_N}.
\]
Since $\ell$ is $\rho$-semi-convex on $\cD_R$
and $\hat z_t = (1-t)z_A + tP_\tau z_B$ with both endpoints in $\cD_R$:
\[
  \ell(\hat z_t, y)
  \le (1-t)\,\ell(z_A, y) + t\,\ell(P_\tau z_B, y)
  + \tfrac{\rho}{2}\,t(1-t)\,\|z_A - P_\tau z_B\|^2.
\]
Component permutation invariance gives
$\ell(P_\tau z_B, y) = \ell(z_B, y)$.
Combining and taking expectations yields~\eqref{eq:general_bound}
after $\sup_t t(1-t) = 1/4$.
\end{proof}

\begin{proof}[Proof of Theorem~\ref{thm:main}]
Theorem~\ref{thm:main} is the instantiation of
Theorem~\ref{thm:general} with $\ell = \cL_{\mathrm{MDN}}$,
$p = 3K$, and the constants from
Proposition~\ref{prop:loss_regularity}.
The only additional detail is the scale-floor clamp: by
Assumption~\ref{ass:arch}, the log-standard-deviation outputs are
clamped to $z_{s,k} \ge \log s_{\min}$, ensuring
$s_k \ge s_{\min}$ for all outputs including interpolated ones.
Proposition~\ref{prop:transfer} applied with the given $M_V^\infty$
and $B_N$ gives the output approximation
\[
  \sup_t \bigl(\E_P[\|z_{M_t} - \hat z_t\|_\infty^2]\bigr)^{1/2}
  \le \tfrac{1}{2}\,L_\sigma M_V^\infty M_x \sqrt{B_N}.
\]
Next, by Lemma~\ref{lem:auto_path_domain} both $z_{M_t}(x)$ and
$\hat z_t(x)$ lie in $\cD_R$ for all $t, x$.
By the architectural clamp (Assumption~\ref{ass:arch}), both
clamped outputs satisfy $s_k \ge s_{\min}$; since $\hat z_t$ is a convex
combination of clamped endpoint outputs, $(\hat z_t)_{s,k} \ge
\log s_{\min}$ and the clamp does not activate on $\hat z_t$.
For the segment from the clamped $z_{M_t}$ to $\hat z_t$:
all points lie in $\cD_R$ (convexity) with $s_k \ge s_{\min}$
(the set $\{z : z_{s,k} \ge \log s_{\min}\}$ is convex),
so the $\ell_\infty$-Lipschitz bound of
Proposition~\ref{prop:loss_regularity} applies.
Moreover, the clamp is $1$-Lipschitz
($|\!\max(a,c)-\max(b,c)| \le |a-b|$), so the transfer gap
satisfies $\|\tilde z_{M_t} - \hat z_t\|_\infty
\le \|z_{M_t} - \hat z_t\|_\infty$
(where $\tilde z_{M_t}$ denotes the clamped interpolated-network
output).  Hence:
\[
  \E_P[\cL_{\mathrm{MDN}}(z_{M_t}, y)]
  \le \E_P[\cL_{\mathrm{MDN}}(\hat z_t, y)]
  + \tfrac{1}{2}\,C_L^\infty L_\sigma M_V^\infty M_x \sqrt{B_N}.
\]
We now apply the semi-convexity bound (Proposition~\ref{prop:loss_regularity})
to the output interpolation $\hat z_t$.
By Lemma~\ref{lem:auto_path_domain}, both endpoints
$z_A(x), P_\tau z_B(x) \in \cD_R$ for all $x$;
since $\cD_R$ is convex, the interpolant
$\hat z_t(x) = (1-t)z_A(x) + tP_\tau z_B(x) \in \cD_R$
as well, so Proposition~\ref{prop:loss_regularity} applies on $\cD_R$ with
$z_1 = z_A(x)$, $z_2 = P_\tau z_B(x)$:
\[
  \cL_{\mathrm{MDN}}(\hat z_t, y)
  \le (1-t)\,\cL_{\mathrm{MDN}}(z_A, y)
  + t\,\cL_{\mathrm{MDN}}(P_\tau z_B, y)
  + \frac{\rho}{2}\,t(1-t)\,\|z_A - P_\tau z_B\|^2.
\]
Finally, component permutation invariance
$\cL_{\mathrm{MDN}}(P_\tau z_B, y) = \cL_{\mathrm{MDN}}(z_B, y)$
(Section~\ref{sec:mdn})
allows us to combine these bounds;
taking expectations yields~\eqref{eq:main_bound}.
\end{proof}

\begin{proof}[Proof of Proposition~\ref{prop:dir_barrier}]
\textbf{Transfer step (data-dependent).}
By the mean value theorem, for each $x$:
\[
  \cL(z_{M_t}, y) - \cL(\hat z_t, y)
  = \nabla_z \cL(\tilde z_t, y)^\top (z_{M_t} - \hat z_t),
\]
where $\tilde z_t(x)$ lies on the segment $[z_{M_t}(x), \hat z_t(x)]$.
Hence
$|\cL(z_{M_t}, y) - \cL(\hat z_t, y)|
\le \|\nabla_z \cL(\tilde z_t, y)\|_1\,\|z_{M_t} - \hat z_t\|_\infty$.
Taking $\E_P$ and applying Cauchy--Schwarz:
\[
  \E_P[\cL(z_{M_t}, y) - \cL(\hat z_t, y)]
  \le \sqrt{\E_P\!\bigl[\|\nabla_z \cL(\tilde z_t, y)\|_1^2\bigr]}
  \;\cdot\;
  \sqrt{\E_P\!\bigl[\|z_{M_t} - \hat z_t\|_\infty^2\bigr]}.
\]
By the same decomposition as in the proof of
Proposition~\ref{prop:transfer}, retaining the directional bound
$|\Delta w_i^\top x|$ (rather than $\|\Delta w_i\|\|x\|$) and
using $\E_P[(\Delta w_i^\top x)^2] = \Delta w_i^\top \Sigma_x\,\Delta w_i$,
the second factor is bounded by
$\tfrac{1}{2}\,L_\sigma M_V^\infty \sqrt{\tilde B_N}$.
Taking $\sup_t$ and using~\eqref{eq:CL_dd}:
\[
  \sup_t\;\E_P[\cL(z_{M_t}, y) - \cL(\hat z_t, y)]
  \;\le\; \tfrac{1}{2}\,\bar C_L\, L_\sigma M_V^\infty \sqrt{\tilde B_N}.
\]

\textbf{Semi-convexity step (directional).}
For the semi-convexity step,
write $\Delta z = z_A(x) - P_\tau z_B(x)$ and
$H(t,x,y) = \nabla^2_z \ell(\hat z_t(x),y)$.
By a second-order Taylor expansion with integral remainder,
the pointwise convexity deficit is
\[
  \cL(\hat z_t, y)
  - (1-t)\cL(z_A, y) - t\cL(z_B, y)
  = -\frac{t(1-t)}{2}\,
    \Delta z^\top H(\xi_t, x, y)\,\Delta z
\]
for some $\xi_t \in [0,1]$
(here we use the fact that $\hat z_t$ is a linear interpolation and
$\cL$ is $C^2$ on $\cD_R$).
Bounding the right-hand side from above:
\[
  -\Delta z^\top H(\xi_t)\,\Delta z
  \;\le\;
  \Delta z^\top (-H(\xi_t))_+ \,\Delta z
  \;\le\;
  \sup_{t'} \Delta z^\top (-H(t'))_+\, \Delta z
  = \rho_{\mathrm{dir}}(x,y).
\]
Hence
$\cL(\hat z_t, y)
\le (1-t)\cL(z_A, y) + t\cL(z_B, y)
  + \frac{t(1-t)}{2}\,\rho_{\mathrm{dir}}(x,y)$.
Taking $\E_P$, then $\sup_t t(1-t) = 1/4$,
yields~\eqref{eq:dir_bound}.
\end{proof}

\subsection{Proof of Theorem~\ref{thm:hp}}
\begin{proof}
The neuron permutation $\nperm$ is the $W$-OT coupling~\eqref{eq:OT_perm},
defined deterministically.
The component permutation is set to $\tau = \mathrm{id}$
for the theoretical bound; this is justified because both networks
share the same mean-field law~$\mathfrak{m}_*$, so the expected
output $z_*(x)$ is identical for both, and the identity permutation
already makes $\E[\Delta_N] = O(1/N)$ (see Event~3 below).
In practice, a Hellinger-affinity matching
(Section~\ref{sec:algorithm}) can further reduce $\Delta_N$
for finite $N$, but the theoretical rate does not require it.

\textbf{Probability budget.}
We allocate the failure probability $\delta$ across three events,
each allowed failure probability $\delta/3$, and apply a union bound.

\textbf{Event 1 ($\Omega_N$): $W$-marginal convergence.}
Define
\begin{align}\label{eq:event_def}
  \Omega_{N}^{W,A}
    &:= \bigl\{\cW_2(\hat\nu_W^A,\mathfrak{m}_{*,W}) \le \eta_N\bigr\},
    \notag\\
  \Omega_{N}^{W,B}
    &:= \bigl\{\cW_2(\hat\nu_W^B,\mathfrak{m}_{*,W}) \le \eta_N\bigr\},
\end{align}
and $\Omega_N := \Omega_N^{W,A} \cap \Omega_N^{W,B}$.
By hypothesis,
$\Prob(\cW_2(\hat\nu_W^X,\mathfrak{m}_{*,W}) > \eta_N) \le \delta/6$
for each network $X$.
By a union bound over the two networks,
$\Pr(\Omega_N^c) \le \delta/3$.

On $\Omega_N$, by the $W$-OT definition and triangle inequality:
\begin{equation}\label{eq:W_OT_rate}
  B_N = \cW_2^2(\hat\nu_W^A,\hat\nu_W^B)
  \le \bigl(\cW_2(\hat\nu_W^A,\mathfrak{m}_{*,W})
           +\cW_2(\hat\nu_W^B,\mathfrak{m}_{*,W})\bigr)^2
  \le 4\eta_N^2.
\end{equation}

\textbf{Event 2 ($E_V$): per-coordinate head-moment control.}
We bound $(M_V^\infty)^2 = \max_j \frac{1}{N}\sum_i (v_i)_j^2$
directly, coordinate by coordinate.
For each network $X \in \{A,B\}$ and coordinate $j \in \{1,\ldots,p\}$,
the sample mean $\bar{S}_{X,j} \coloneqq \frac{1}{N}\sum_{i=1}^N (v_{X,i})_j^2$
has expectation $\E[\bar{S}_{X,j}] = \E_{\mathfrak{m}_*}[(V)_j^2]
\le m_{\infty,*}^2$ and variance
$\mathrm{Var}(\bar{S}_{X,j}) \le m_{4,*}/N$
(since $\mathrm{Var}(Z^2) \le \E[Z^4]$ for any $Z$).
By Chebyshev's inequality,
$\Pr\!\bigl(\bar{S}_{X,j} > m_{\infty,*}^2 + t\bigr)
\le m_{4,*}/(Nt^2)$.
Setting $t$ so that the right-hand side equals $\delta/(6p)$
(to allow a union bound over $2$ networks $\times$ $p$ coordinates),
we need $t = \sqrt{6p\,m_{4,*}/(N\delta)}$.
Define $\varepsilon_N(\delta) \coloneqq \sqrt{6p\,m_{4,*}/(N\delta)}$ and
$E_V \coloneqq \bigl\{(M_{V,X}^\infty)^2 \le m_{\infty,*}^2 + \varepsilon_N(\delta)
\;\;\forall\, X \in \{A,B\}\bigr\}$.
By the union bound over $2p$ events,
$\Pr(E_V^c) \le \delta/3$.
On $E_V$, $(M_V^\infty)^2 \le m_{\infty,*}^2 + \varepsilon_N(\delta)$,
and as $N \to \infty$, $\varepsilon_N(\delta) \to 0$, so
$(M_V^\infty)^2 \to m_{\infty,*}^2$ at the SLLN rate---without the
crude $\ell_2$-to-$\ell_\infty$ relaxation of the Markov approach.

\textbf{Mean-field CLT bound.}
Under the common mean-field law $\mathfrak{m}_*$
(Assumption~\ref{ass:mf_law}), both $z_A(x)$ and $z_B(x)$
are sample means of $N$ i.i.d.\ draws, converging to the same limit
$z_*(x) \coloneqq \E_{\mathfrak{m}_*}[\sigma(W^\top x)V]$.
Writing
$z_X(x) - z_*(x) = N^{-1}\sum_i \xi_i(x)$ with
$\xi_i(x) = \sigma(W_{X,i}^\top x)V_{X,i} - z_*(x)$ i.i.d.\ mean-zero
under $\mathfrak{m}_*$, the variance of the sample mean gives
\begin{equation}\label{eq:CLT_var}
  \E\!\bigl[\E_P[\|z_X(x)-z_*(x)\|^2]\bigr]
  \;\le\; \frac{S_\sigma^2\,m_{2,*}}{N}
\end{equation}
for each $X\in\{A,B\}$.
Since both networks share $\mathfrak{m}_*$ and are drawn independently,
taking $\tau = \mathrm{id}$ and using the bias-variance decomposition:
\begin{equation}\label{eq:Delta_CLT}
  \E_{A,B}[\Delta_N]
  \;\le\; \E_{A,B}\!\bigl[\E_P[\|z_A(x)-z_B(x)\|^2]\bigr]
  \;=\; \E_P\!\bigl[\mathrm{Var}_A(z_A(x)) + \mathrm{Var}_B(z_B(x))\bigr]
  \;\le\; \frac{2S_\sigma^2\,m_{2,*}}{N},
\end{equation}
where the equality uses the independence of $A$ and $B$
(the cross term $\E_A[z_A]\cdot\E_B[z_B] = z_*\cdot z_*$ cancels),
and each variance term is bounded by~\eqref{eq:CLT_var}.
This improves the na\"{\i}ve triangle-inequality factor from $4$ to $2$.

\textbf{Event 3 ($E_\Delta$): mismatch control.}
By~\eqref{eq:Delta_CLT}, $\E[\Delta_N] \le 2S_\sigma^2 m_{2,*}/N$.
Markov's inequality gives
$\Pr(\Delta_N > 6S_\sigma^2 m_{2,*}/(N\delta)) \le \delta/3$;
call this event $E_\Delta$.

\textbf{Conclusion.}
By the union bound,
$\Pr(\Omega_N \cap E_V \cap E_\Delta) \ge 1 - \delta/3 - \delta/3 - \delta/3 = 1-\delta$.
On this joint event all three bounds hold simultaneously, and
Theorem~\ref{thm:general} yields the barrier inequality~\eqref{eq:general_bound}.
\end{proof}

\subsection{Proof of Corollary~\ref{cor:asymptotic}}
\begin{proof}
All three events $\Omega_N$, $E_V$, $E_\Delta$ of Theorem~\ref{thm:hp}
hold simultaneously with probability $\ge 1-\delta$.

On $\Omega_N$, the $W$-OT bound~\eqref{eq:W_OT_rate} gives $B_N \le 4\eta_N^2$,
so $M_x\sqrt{B_N} \le 2M_x\eta_N$.
Under Assumption~\ref{ass:subgauss}, the sub-Gaussian moment condition
gives, via covering-number bounds~\cite{lei2020}
(see also~\cite[Theorem~2]{weed2019} under $T_2$),
$\E[\cW_2(\hat\nu_W^X,\mathfrak{m}_{*,W})]
\le C\kappa_W\sqrt{d/N}$;
combined with sub-Gaussian concentration of $\cW_2$
around its mean~\cite[Theorem~3.1]{weed2019},
we may set
$\eta_N = C'\kappa_W\sqrt{d/N}
+ C''\kappa_W\sqrt{\log(6/\delta)/N}$
(which is $O(\kappa_W\sqrt{d/N})$ for fixed $\delta$)
to satisfy the hypothesis of Theorem~\ref{thm:hp},
yielding the transfer term $C_{\mathrm{tr}}(\delta,N)\,\kappa_W\sqrt{d/N}$.

On $E_V$,
$(M_V^\infty)^2 \le m_{\infty,*}^2 + \varepsilon_N(\delta)$
with $\varepsilon_N(\delta) = \sqrt{6p\,m_{4,*}/(N\delta)}$,
which controls
$R(\delta,N) = S_\sigma \hat M_V(\delta,N)$
and hence $C_L^\infty$ and $\rho$ via the loss-specific regularity
proposition (Proposition~\ref{prop:loss_regularity} for MDN,
Proposition~\ref{prop:vmf_regularity} for vMF).
As $N \to \infty$, $\hat M_V \to m_{\infty,*}$ and the prefactors
converge to the limits $C_{\mathrm{tr}}^*$, $C_{\mathrm{sc}}^*$.

On $E_\Delta$, $\Delta_N \le 6S_\sigma^2 m_{2,*}/(N\delta)$, yielding
the semi-convexity term $C_{\mathrm{sc}}(\delta,N) \cdot m_{2,*}/N$.

For fixed $K$, $m_{2,*} = O(1)$, so $C_{\mathrm{tr}}(\delta,N) \to C_{\mathrm{tr}}^*$
and $C_{\mathrm{sc}}(\delta,N) \to C_{\mathrm{sc}}^*$ stabilise and
$\cB = O(\kappa_W\sqrt{d/N} + N^{-1})$.
For growing $K$ with $m_{\infty,*}^2 = O(1)$,
$m_{2,*} = O(K)$ (Section~\ref{sec:asymptotic_rate}),
$\varepsilon_N(\delta) = O(\sqrt{K/N})$, and
$\cB = O(\kappa_W\sqrt{d/N} + K/N)$.
\end{proof}

\begin{proof}[Proof of Corollary~\ref{cor:struct_rate}]
On the same high-probability event as Theorem~\ref{thm:hp},
\[
  B_N \le 4\eta_N^2,
  \qquad
  M_x\sqrt{B_N} \le 2M_x\eta_N,
  \qquad
  M_V^\infty \le \hat M_V(\delta,N),
  \qquad
  \Delta_N \le \frac{6S_\sigma^2 m_{2,*}}{N\delta}.
\]
By the auxiliary bounds~\eqref{eq:struct_aux},
\[
  \mathcal R_N(\nperm,\tau) \le C_R\,\eta_N,
  \qquad
  \sqrt{Q_N} \le \sqrt{C_Q}\,\eta_N^2.
\]
Repeating the proof of Theorem~\ref{thm:general} with
Proposition~\ref{prop:transfer_struct} in place of
Proposition~\ref{prop:transfer} gives
\begin{align*}
  \cB
  &\le
  C_L^\infty \!\left[
    \tfrac{1}{4}L_\sigma
    \bigl(2L_g\eta_N + C_R\eta_N\bigr)
    2M_x\,\eta_N
    + \tfrac{1}{8}L_{\sigma'} \hat M_V(\delta,N)\sqrt{C_Q}\,\eta_N^2
  \right]
  + C_{\mathrm{sc}}(\delta,N)\,\frac{m_{2,*}}{N} \\
  &=
  C_{\mathrm{fast}}(\delta,N)\,\eta_N^2
  + C_{\mathrm{sc}}(\delta,N)\,\frac{m_{2,*}}{N},
\end{align*}
which is~\eqref{eq:struct_rate}.
Under Assumption~\ref{ass:subgauss},
$\eta_N = O(\kappa_W\sqrt{d/N})$
by Corollary~\ref{cor:asymptotic}, so
$\eta_N^2 = O(\kappa_W^2 d/N)$ and
the asymptotic rate in~\eqref{eq:struct_rate} follows.
\end{proof}

\subsection{Proof of Proposition~\ref{prop:vmf_regularity}
  (vMF loss regularity)}\label{app:vmf_proof}

\begin{proof}
We follow the same Hessian-decomposition strategy as
Proposition~\ref{prop:loss_regularity}, using
$u_k = \log\pi_k + \log C_{d_y}(\|\eta_k\|) + \eta_k^\top y$
so that
$\cL_{\vmf} = -\log\sum_k e^{u_k}$
and~\eqref{eq:hess_decomp} gives
$\nabla^2\cL_{\vmf}
= -\sum_k\gamma_k\nabla^2 u_k - \mathrm{Cov}_\gamma(\nabla u_k)$.

\paragraph{Gradient bounds.}
The gradient of $u_k$ splits into a $z_\pi$-block and an $\eta_k$-block.

\emph{$z_\pi$-block.}
$\nabla_{z_\pi} u_k = e_k - \pi$ (identical to MDN).
$\|e_k - \pi\|^2 \le 2$.

\emph{$\eta_k$-block.}
Since the vMF distribution belongs to the exponential family with
natural parameter $\eta_k$ and sufficient statistic $y$,
the log-partition function is
$\psi(\eta) = -\log C_{d_y}(\|\eta\|)$
(convex, with gradient
$\nabla_\eta\psi(\eta) = A_{d_y}(\|\eta\|)\,\eta/\|\eta\|$
where $A_{d_y}(\kappa)=I_{d_y/2}(\kappa)/I_{d_y/2-1}(\kappa)$
is the mean resultant length~\cite{mardia2000}).
Therefore
\[
  \nabla_{\eta_k} u_k
  = y - \nabla_\eta\psi(\eta_k)
  = y - A_{d_y}(\kappa_k)\,\mu_k.
\]
Since $\|y\|=1$, $\|\mu_k\|=1$, and $0 \le A_{d_y}(\kappa)<1$
for all $\kappa \ge 0$:
\begin{equation}\label{eq:vmf_grad_bound}
  \|\nabla_{\eta_k} u_k\| \le \|y\| + A_{d_y}(\kappa_k) \le 2.
\end{equation}
At $\kappa_k=0$, $\nabla_{\eta_k} u_k = y$ (the uniform limit),
which is well-defined with $\|y\|=1$.

Combining both blocks:
$\|\nabla u_k\|^2 \le 2 + 4 = 6$ for all $k$.

\paragraph{Lipschitz bound.}
$C_{L,\vmf}^\infty = \sup_z \|\nabla\cL_{\vmf}\|_1$.
The $z_\pi$-block contributes
$\|\gamma-\pi\|_1 \le 2$ (as in MDN).
The $\eta$-blocks contribute
$\sum_k \gamma_k\|y-A_{d_y}(\kappa_k)\mu_k\|_1
\le \sum_k \gamma_k \cdot \sqrt{d_y}\cdot 2 = 2\sqrt{d_y}$,
where $\|v\|_1 \le \sqrt{d_y}\|v\|_2$ for $v \in \R^{d_y}$.
Hence $C_{L,\vmf}^\infty \le 2 + 2\sqrt{d_y}$.

\paragraph{Semi-convexity bound.}
\emph{Weighted Hessian.}
The $z_\pi$-block:
$(-\sum_k\gamma_k\nabla^2 u_k)_{z_\pi} = C_\pi \succeq 0$
(identical to MDN).

The $\eta_k$-block:
$\nabla^2_{\eta_k} u_k = -\nabla^2\psi(\eta_k)$.
Since $\psi$ is the log-partition function of an exponential family,
$\nabla^2\psi(\eta_k) = \mathrm{Cov}_{\vmf(\mu_k,\kappa_k)}[y]
\succeq 0$
(Fisher information matrix).
Therefore
\[
  -\gamma_k\nabla^2_{\eta_k} u_k
  = \gamma_k\nabla^2\psi(\eta_k) \succeq 0
  \qquad \text{for each } k.
\]
Cross-blocks ($\eta_j$ vs $\eta_k$, $j \ne k$) are zero since
$u_k$ depends only on $\eta_k$.
Hence
\begin{equation}\label{eq:vmf_weighted_hess}
  \lambda_{\min}\!\Bigl(-\sum_k\gamma_k\nabla^2 u_k\Bigr) \ge 0.
\end{equation}
This is the key contrast with MDN, where the same quantity is
$\ge -1/s_{\min}^2$.

\emph{Covariance term.}
By the same argument as in the MDN proof
(using $\mathrm{Var}_\gamma(X) \le \E_\gamma[X^2]$):
\[
  \|\mathrm{Cov}_\gamma(\nabla u_k)\|_{\mathrm{op}}
  \le \max_k \|\nabla u_k\|^2 \le 6.
\]

\emph{Assembly.}
By~\eqref{eq:weyl_split}:
\[
  \rho_{\vmf}
  \le 0 + 6 = 6. \qedhere
\]
\end{proof}

%======================================================================
\section{Additional Theoretical Results}\label{app:additional}
%======================================================================

This appendix collects results that complement the main theorems
but are not needed for the core barrier bounds.
Throughout this appendix, we use the
\emph{$\Sigma_x$-weighted input mismatch}
$\tilde B_N \coloneqq \frac{1}{N}\sum_{i=1}^N
\Delta w_i^\top \Sigma_x\,\Delta w_i$
(where $\Sigma_x = \E_P[xx^\top]$),
which satisfies $\tilde B_N \le M_x^2 B_N$
and yields sharper bounds when the data covariance is anisotropic.

\subsection{Full Symmetry Characterisation}\label{app:symmetry_full}

\begin{proposition}[Symmetry characterisation]\label{prop:symmetry}
Let $\sigma$ be bounded and non-polynomial (e.g.\ any activation
satisfying Assumption~\ref{ass:activ}), and let $d \ge 2$.
Assume the two parameter sets are in generic position
(Definition~\ref{def:generic}).
If they define the same
conditional density $p(y \mid x)$, then they are related by
a neuron permutation $\nperm \in \cS_N$, a component permutation
$\tau \in \cS_K$, and a softmax translation $\delta \in \R^N$.
The symmetry group of the mixture-model parameterisation is
$\cS_N \times \cS_K \times \R^N_\pi$.
\end{proposition}

The proof (Appendix~\ref{app:proof_symmetry}) uses a two-level
identification: Gaussian mixture
identifiability~\cite{yakowitz1968} recovers $\tau$, then
ridge-function linear independence~\cite{pinkus1999} recovers
$\nperm$ and $\delta$.

\subsection{Covariance-Aware Hidden Alignment}\label{app:cov_aware}

Define the \textbf{$\Sigma_x$-aware coupling}
\begin{equation}\label{eq:OT_perm_sigma}
  \nperm_{\Sigma}
  \;\coloneqq\;
  \argmin_{q \in \cS_N}
  \frac{1}{N}\sum_{i=1}^N
  \bigl\|\Sigma_x^{1/2}(w_{A,i} - w_{B,q(i)})\bigr\|^2.
\end{equation}

\begin{corollary}[Covariance-aware hidden alignment]\label{cor:sigma_ot}
Let $U = \Sigma_x^{1/2}W$ and choose the hidden permutation
$\nperm_{\Sigma}$ by~\eqref{eq:OT_perm_sigma}.
Then $\nperm_{\Sigma}$ minimises the transfer quantity
$\tilde B_N$ over all hidden-layer permutations, and
$\tilde B_N = \cW_2^2(\hat\nu_U^A,\hat\nu_U^B)$.
In particular, whenever $U$ admits a sub-Gaussian $\cW_2$ rate in an
intrinsic dimension $d_{\Sigma}$, the transfer term scales as
$O(\sqrt{d_{\Sigma}/N})$.
\end{corollary}

\subsection{Refined Output Coupling}\label{app:refined_coupling}

Define the \emph{per-coordinate gauge-fixed head mismatch}
\begin{equation}\label{eq:head_mismatch_inf}
  M_{\Delta \tilde V}^\infty
  \;\coloneqq\;
  \max_{j=1,\ldots,p}
  \Bigl(\frac{1}{N}\sum_{i=1}^N
  \bigl(\tilde v_{A,i} - P_\tau \tilde v_{B,\nperm(i)}\bigr)_j^2\Bigr)^{\!1/2}.
\end{equation}

\begin{proposition}[Refined output coupling]\label{prop:transfer_refined}
Under Assumptions~\ref{ass:arch} and~\ref{ass:activ}\,(i)--(iii),
and~\ref{ass:input_mom} (including the fourth-moment condition
$\E_P[\|x\|^4] \le M_{x,4}^4$):
\begin{equation}\label{eq:transfer_refined}
  \sup_{t\in[0,1]}
  \bigl(\E_P[\|z_{M_t}(x) - \hat z_t(x)\|_\infty^2]\bigr)^{1/2}
  \;\le\;
  \tfrac{1}{4}\,L_\sigma\,M_{\Delta \tilde V}^\infty\,\sqrt{\tilde B_N}
  \;+\;
  \tfrac{1}{8}\,L_{\sigma'}\,M_V^\infty\,\sqrt{Q_N},
\end{equation}
where $Q_N = \frac{1}{N}\sum_{i} \E_P[(\Delta w_i^\top x)^4]$.
Under $W$-only alignment, $M_{\Delta \tilde V}^\infty$ can be as large as
$2M_V^\infty$.
\end{proposition}

\subsection{Data-Dependent Barrier}\label{app:data_dep_barrier}

\begin{proposition}[Data-dependent barrier]\label{prop:dir_barrier}
Under the same hypotheses as Theorem~\ref{thm:general},
define the \emph{directional semi-convexity cost}
\begin{equation}\label{eq:rho_dir}
  \rho_{\mathrm{dir}}(x,y)
  \;\coloneqq\;
  \sup_{t\in[0,1]}\;
  \Delta z(x)^\top
  \bigl(-\nabla^2_z \ell(\hat z_t(x),\,y)\bigr)_{\!+}\;
  \Delta z(x),
\end{equation}
and the \emph{data-dependent Lipschitz constant}
\begin{equation}\label{eq:CL_dd}
  \bar C_L
  \;\coloneqq\;
  \sup_{t\in[0,1]}
  \sqrt{\E_P\!\bigl[\|\nabla_z \ell(\tilde z_t(x),\,y)\|_1^2\bigr]}.
\end{equation}
Then:
\begin{equation}\label{eq:dir_bound}
  \cB \;\le\;
  \tfrac{1}{2}\,\bar C_L\, L_\sigma M_V^\infty \sqrt{\tilde B_N}
  \;+\; \frac{1}{8}\,\E_P\!\bigl[\rho_{\mathrm{dir}}(x,y)\bigr],
\end{equation}
with $\bar C_L \le C_L^\infty$, so~\eqref{eq:dir_bound} is at least
as sharp as~\eqref{eq:barrier_form}.
The transfer term replaces the worst-case $C_L^\infty$ by
the data-averaged gradient norm;
for well-fitted models this yields a $\sim\!5000\times$ reduction
(Appendix~\ref{app:nonvacuous_detail}).
\end{proposition}

%======================================================================
\section{Mean-Field Justification of Assumption~\ref{ass:mf_law}}\label{app:A4_mf}
%======================================================================

Assumption~\ref{ass:mf_law} requires that the post-training neuron
parameters of both networks are i.i.d.\ draws from a common law
$\mathfrak{m}_*$.
This appendix justifies the assumption via propagation of chaos,
first under gradient flow (Section~\ref{app:mf_gf}) and then under
SGD (Section~\ref{app:mf_sgd}), and summarises the resulting
validity conditions (Section~\ref{app:mf_summary}).

Throughout, the training-time variable is $s \in [0,T]$
(distinct from the interpolation parameter $t \in [0,1]$) and
$f = O_T(g)$ means $|f| \le C_T g$ for a constant $C_T$ depending
on $T$ but not on $N$ or $K$.

%----------------------------------------------------------------------
\subsection{Gradient Flow}\label{app:mf_gf}
%----------------------------------------------------------------------

The argument proceeds in three steps: (i)~verify the Lipschitz
condition on the mean-field drift, (ii)~deduce McKean--Vlasov uniqueness,
(iii)~apply propagation of chaos.

\paragraph{Step 1: Lipschitz drift for the mixture-model loss.}
Define the mean-field drift
$b(\theta,\mathfrak{m})
\coloneqq -\E_P[\nabla_\theta\ell(z_{\mathfrak{m}}(x),y)]$.
The mixture-model loss is non-convex, but propagation of chaos requires only
that $b$ be Lipschitz in
$(\theta,\mathfrak{m})$~\cite{mei2018meanfield,rotskoff2022},
not convexity.
Under Assumptions~\ref{ass:bdd_resp}, \ref{ass:activ},
and~\ref{ass:input_mom}, Proposition~\ref{prop:loss_regularity}
controls the output-space gradient on any compact domain $\cD_R$.
Hence, along finite-time trajectories whose mean-field outputs remain
in some $\cD_R$, the head drift is bounded and the input-weight drift
has at most linear growth in $\|v\|$
(cf.~Proposition~\ref{prop:a6_auto}).
This is the finite-time regularity regime in which the standard
propagation-of-chaos heuristics apply.

\paragraph{Step 2: McKean--Vlasov uniqueness.}
Under the $1/N$ mean-field parameterisation, per-neuron gradient flow
converges to the McKean--Vlasov equation
$\partial_s\mathfrak{m}_s
= -\nabla\cdot(\mathfrak{m}_s\,b(\cdot,\mathfrak{m}_s))$.
Because $b$ is Lipschitz, Cauchy--Lipschitz uniqueness on
$(\mathcal{P}_2,\cW_2)$~\cite{carmona2018probabilistic} gives a
unique trajectory $(\mathfrak{m}_s)_{s\in[0,T]}$ from any initial law
$\mathfrak{m}_0$.
Given $\mathfrak{m}_0$, both networks follow the same deterministic
flow and reach the \emph{same} limiting measure $\mathfrak{m}_T$,
regardless of whether $\mathfrak{m}_T$ is a global or local stationary
point.
The individual parameters $\theta^i_A, \theta^j_B$ are different
realisations, but their empirical distributions both converge to
$\mathfrak{m}_T$.

\paragraph{Step 3: Propagation of chaos.}
A standard coupling argument (i.i.d.\ initialisation from
$\mathfrak{m}_0$, Lipschitz drift, Gr\"{o}nwall) gives
\begin{equation}\label{eq:coupling_poc}
  E_N(T)
  \;\coloneqq\;
  \tfrac{1}{N}\sum_i\E[\|\theta_{i,T}-\bar\theta_{i,T}\|^2]
  \;=\; O_T(N^{-1}),
\end{equation}
where $\bar\theta_{i,T}$ is the mean-field limit trajectory for
neuron $i$ and the $O_T$ constant is
$C_T = O(e^{C_L T})$~\cite{mei2018meanfield,rotskoff2022}.
As $N$ grows, the neurons decouple and
Assumption~\ref{ass:mf_law} holds with $\mathfrak{m}_* = \mathfrak{m}_T$.

%----------------------------------------------------------------------
\subsection{Extension to SGD}\label{app:mf_sgd}
%----------------------------------------------------------------------

Practical training uses SGD with finite learning rate $\alpha$ and
mini-batch size $b$, introducing two additional error sources beyond
the finite-width error of~\eqref{eq:coupling_poc}.
The total approximation error decomposes as
\begin{equation}\label{eq:sgd_error}
  \cW_2(\hat\mu_{\mathrm{SGD}}^N,\,\mathfrak{m}_T)
  \;\le\;
  \underbrace{O_T(N^{-1/2})}_{\text{finite width}}
  \;+\; \underbrace{O_T(\alpha)}_{\text{time discret.}}
  \;+\; \underbrace{O_T(b^{-1/2})}_{\text{mini-batch}}.
\end{equation}
All three $O_T$ constants inherit the Lipschitz
constant $C_L$ (Proposition~\ref{prop:loss_regularity}).
We now discuss the theoretical basis for each term.

\paragraph{Learning-rate regimes.}
De~Bortoli et al.~\cite{debortoli2020} establish quantitative
propagation of chaos for a continuous-time SGD counterpart applied to
two-layer overparameterised networks, with the learning rate scaling
as $\alpha = \Theta(N^\beta)$ in the $1/N$ parameterisation.
They identify two regimes:
\begin{itemize}[itemsep=2pt]
  \item \textbf{Regime~1} ($\beta < 1$):
    the SGD noise averages out at rate $O(N^{\beta-1})$, and the
    limiting dynamics is the \emph{same deterministic}
    McKean--Vlasov ODE as gradient flow.
    Propagation of chaos holds with the standard $O_T(N^{-1})$ rate.
  \item \textbf{Regime~2} ($\beta = 1$):
    the SGD noise survives in the limit, yielding a
    McKean--Vlasov SDE with data-dependent diffusion---an implicitly
    regularised optimisation.
\end{itemize}
Standard practice uses $\alpha = O(1)$ (constant, independent of $N$),
corresponding to $\beta = 0 < 1$, squarely in Regime~1.

\paragraph{Law of large numbers and CLT.}
Sirignano and Spiliopoulos~\cite{sirignano2020lln} prove a law of
large numbers for networks trained by SGD: the empirical parameter
distribution converges to the solution of a nonlinear PDE in the
joint limit of large width and many iterations.
The accompanying CLT~\cite{sirignano2020clt} characterises
the $O(N^{-1/2})$ fluctuations around this deterministic limit.
Both results require Lipschitz regularity of the loss gradient,
supplied for each application by the loss regularity analysis
(Proposition~\ref{prop:loss_regularity} for MDN,
Proposition~\ref{prop:vmf_regularity} for vMF).

\paragraph{Non-convexity and finer-grained objects.}
Convexity is not required: Nguyen and
Pham~\cite{nguyen2023multilayer} treat non-convex losses explicitly;
any mixture-model loss with bounded gradient qualifies since its drift is Lipschitz
(Propositions~\ref{prop:loss_regularity} and~\ref{prop:vmf_regularity}).
For finer-grained analysis, stochastic modified
flows~\cite{gess2024smf} capture residual SGD fluctuations, and
mean-field Langevin bounds~\cite{suzuki2023mfld} avoid the
$e^{C_L T}$ blow-up under log-Sobolev conditions.

%----------------------------------------------------------------------
\subsection{Validity Conditions and Limitations}\label{app:mf_summary}
%----------------------------------------------------------------------

Table~\ref{tab:a5_validity} collects the conditions under which
Assumption~\ref{ass:mf_law} is justified.

\begin{table}[H]
\centering\small
\renewcommand{\arraystretch}{1.15}
\caption{Validity conditions for Assumption~\ref{ass:mf_law}.}
\label{tab:a5_validity}
\begin{tabular}{@{}p{3.8cm}p{5.2cm}p{3.6cm}@{}}
\toprule
\textbf{Condition} & \textbf{Role} & \textbf{Practical status} \\
\midrule
Same initialisation distribution $\mathfrak{m}_0$
  & McKean--Vlasov uniqueness: same $\mathfrak{m}_0 \Rightarrow$ same trajectory $\mathfrak{m}_T$
  & Satisfied by standard init.\ schemes \\
\addlinespace
Same training recipe (lr, batch, schedule)
  & Different recipes $\Rightarrow$ different drifts $\Rightarrow$ different $\mathfrak{m}_T$
  & Satisfied in model merging / FL \\
\addlinespace
Learning rate $\alpha = o(N)$
  & Regime~1~\cite{debortoli2020}: SGD noise vanishes, gradient-flow limit applies
  & $\alpha = O(1)$: trivially satisfied \\
\addlinespace
Width $N \gg e^{C_L T}$
  & Propagation of chaos: $E_N(T) = O(e^{C_L T}/N)$
  & Bottleneck for long training \\
\addlinespace
Lipschitz drift
  & Required for propagation of chaos; \emph{not} convexity
  & Props.~\ref{prop:loss_regularity},~\ref{prop:vmf_regularity}: $C_L$ indep.\ of $K$ \\
\bottomrule
\end{tabular}
\end{table}

\paragraph{When the approximation is weakest.}
Two regimes deserve attention:
\begin{enumerate}[label=(\roman*),leftmargin=*,itemsep=2pt]
  \item \emph{Diffuse responsibilities} ($\gamma_k \approx 1/K$):
    the shared softmax denominator creates inter-neuron correlations;
    propagation of chaos still holds (the Lipschitz condition is
    unaffected by curvature), but the $O_T$ constant may be larger.
    This regime is typically transient: gradient descent drives
    components to specialise, decreasing
    $\bar\chi = \sum_k \gamma_k(1-\gamma_k)$ along the training
    trajectory.
  \item \emph{Long training horizon} ($T \gg 1$):
    the Gr\"{o}nwall constant $C_T = O(e^{C_L T})$ grows
    exponentially.
    Uniform-in-time bounds require additional structure (log-Sobolev
    inequalities~\cite{suzuki2023mfld} or dissipativity).
    In practice, early-stopped or converged networks have moderate
    effective $T$.
\end{enumerate}

\paragraph{Implications for the barrier bound.}
Under the conditions of Table~\ref{tab:a5_validity},
substituting~\eqref{eq:sgd_error} into the $\cW_2$ rate of
Corollary~\ref{cor:asymptotic} shows that the barrier
$\cB = O(\kappa_W\sqrt{d/N} + K/N)$ remains valid for SGD-trained networks,
provided the time-discretisation and mini-batch errors are dominated
by the finite-width term---a condition easily met by taking
$\alpha = O(1)$ and $b = \Omega(\sqrt{N})$.

%======================================================================
\section{Empirical Validation of Assumptions~\ref{ass:mf_law}--\ref{ass:subgauss}}\label{app:assumption_validation}
%======================================================================

This appendix provides empirical evidence supporting the
mean-field assumptions (\ref{ass:mf_law}--\ref{ass:subgauss}) using a mean-field MDN
training setup.
Networks are trained with $1/N$ mean-field scaling,
sigmoid activation ($S_\sigma = 1$, $L_\sigma = 1/4$),
Gaussian initialisation, and full-batch gradient descent
(lr $= 10^{-2}$, $10{,}000$ epochs)
on the
\texttt{synthetic} ($y \mid x \sim 0.5\,\cN(x, 0.3^2)
+ 0.5\,\cN(-x, 0.3^2)$) and
\texttt{synthetic\_transpose} datasets
with $K \in \{1, 5\}$ and
$N \in \{1024, 2048, 4096, 8192, 16384\}$ (10 seeds each).

%----------------------------------------------------------------------
\subsection{\ref{ass:mf_law}: Mean-field limiting law}\label{app:a5_validation}
%----------------------------------------------------------------------

Assumption~\ref{ass:mf_law} posits that post-training neuron parameters
are i.i.d.\ draws from a common law $\mathfrak{m}_*$.
Figure~\ref{fig:a5_overlay} overlays the neuron density estimates
$(w_i, v_i)$ from two independently seeded networks at the same width.
At $N = 4096$ the density contours from different seeds are visibly
closer than at $N=1024$, qualitatively consistent with propagation of
chaos driving both sets of neurons toward the same $\mathfrak{m}_T$.
The quantitative mismatch diagnostics in
Appendix~\ref{app:nonvacuous_detail} show, however, that for
$K=5$ the range $N \le 16384$ is still pre-asymptotic.
Figure~\ref{fig:a5_predictive} compares the predictive densities
$p(y \mid x)$ across seeds; again the outputs converge as $N$ grows.

\begin{figure}[H]
\centering
\includegraphics[width=0.48\textwidth]{synthetic/K5/full_batch/overlay/w1024_s0_s1.png}
\hfill
\includegraphics[width=0.48\textwidth]{synthetic/K5/full_batch/overlay/w4096_s0_s1.png}
\caption{Neuron density overlays (seed 0 vs.\ seed 1) at
  $N=1024$ (left) and $N=4096$ (right) for
  \texttt{synthetic}, $K=5$.
  As width increases, the two distributions move closer.
  This provides qualitative support for
  Assumption~\ref{ass:mf_law}.}
\label{fig:a5_overlay}
\end{figure}

\begin{figure}[H]
\centering
\includegraphics[width=0.48\textwidth]{synthetic/K5/full_batch/predictive/w1024_predictive.png}
\hfill
\includegraphics[width=0.48\textwidth]{synthetic/K5/full_batch/predictive/w4096_predictive.png}
\caption{Predictive density $p(y \mid x)$ for two seeds at
  $N=1024$ (left) and $N=4096$ (right).
  Functional convergence accompanies the qualitative parameter-space
  convergence shown in Figure~\ref{fig:a5_overlay}.}
\label{fig:a5_predictive}
\end{figure}

%----------------------------------------------------------------------
\subsection{Finite moments from finite training time under bounded outputs}\label{app:a6_proposition}
%----------------------------------------------------------------------

The moment quantities $m_{2,*}$, $m_{4,*}$ appearing in the barrier
bounds are finite under suitable a priori control of the mean-field
output trajectory.
The following proposition gives a sufficient criterion.

\begin{proposition}[Finite moments under gradient flow with bounded outputs]\label{prop:a6_auto}
Suppose the network is trained by gradient flow
$\dot\theta_i(s) = b(\theta_i(s), \mathfrak{m}_s)$ for
$s \in [0,T]$, with $\mathfrak{m}_0$ satisfying
$\E_{\mathfrak{m}_0}[\|\theta\|^4] < \infty$
(e.g.\ Gaussian, He, or Xavier initialisation).
Assume additionally that there exists $R < \infty$ such that
$z_{\mathfrak{m}_s}(x) \in \cD_R$ for all $s \in [0,T]$ and
$P$-almost every $x$.
Under Assumptions~\ref{ass:bdd_resp}, \ref{ass:activ},
and~\ref{ass:input_mom}:
\begin{enumerate}[label=(\roman*),leftmargin=*,itemsep=2pt]
  \item The head-weight drift is uniformly bounded:
    $\|b_v(\theta, \mathfrak{m})\| \le C_b$
    for a constant $C_b$ depending on $R$, $S_\sigma$, $Y_{\max}$,
    and $s_{\min}$, but not on $\|v\|$ or $\|w\|$.
  \item $\E_{\mathfrak{m}_T}[\|V\|^4] < \infty$;
    in particular, $m_{2,*} = \int\|V\|^2\,d\mathfrak{m}_T < \infty$
    and $m_{4,*} = \max_j \E_{\mathfrak{m}_T}[(V)_j^4] < \infty$.
  \item For the input weights, $\|w_i(s)\|$ grows at most linearly
    in $\|v_i(s)\|$:
    $\|w_i(T)\| \le \|w_i(0)\| + C'(\|v_i(0)\|T + C_b T^2)$,
    so $\E_{\mathfrak{m}_T}[\|W\|^2] < \infty$.
\end{enumerate}
\end{proposition}

\begin{proof}
\textbf{Part (i).}
The mean-field drift for the head weights is
$b_v(\theta, \mathfrak{m})
= -\E_P[\sigma(w^\top x)\,\nabla_z \ell(z_\mathfrak{m}(x),\, y)]$,
where $z_\mathfrak{m}(x) = \int \sigma(W^\top x)\,V\,d\mathfrak{m}$
is determined by the distribution $\mathfrak{m}$, not by the
individual particle $(v,w)$.
By hypothesis, $z_{\mathfrak{m}_s}(x) \in \cD_R$.
Since $|\sigma| \le S_\sigma$ (Assumption~\ref{ass:activ}(i)) and
$\|\nabla_z \ell\|$ is bounded on $\cD_R$
(by the loss-specific $C_L^\infty$),
$\|b_v\|$ is bounded by a constant $C_b$ that depends on
$R$, $S_\sigma$, and the loss regularity constants, but crucially not on
$\|v\|$ or $\|w\|$.

\textbf{Part (ii).}
Since $\dot v(s) = b_v(\theta(s), \mathfrak{m}_s)$ with
$\|b_v\| \le C_b$, we have
$\|v(T)\| \le \|v(0)\| + C_b T$.
Hence
$\|v(T)\|^4 \le 8(\|v(0)\|^4 + (C_b T)^4)$, so
$\E[\|V_T\|^4] \le 8(\E[\|V_0\|^4] + (C_b T)^4) < \infty$
under the hypothesis $\E[\|\theta_0\|^4] < \infty$.
Since $m_{2,*} \le (\E[\|V_T\|^4])^{1/2}$ and
$m_{4,*} \le \E[\|V_T\|^4]$ by Jensen, both are finite.

\textbf{Part (iii).}
The input-weight drift is
$b_w(\theta, \mathfrak{m})
= -\E_P[\sigma'(w^\top x)\,x\cdot
  v^\top\nabla_z \ell(z_\mathfrak{m}(x),\,y)]$.
Since $|\sigma'| \le L_\sigma$, $\E_P[\|x\|] \le M_x$, and
$\|\nabla_z \ell\|$ is bounded on the output domain,
$\|b_w(\theta, \mathfrak{m})\| \le C_w(1 + \|v\|)$
for a constant $C_w$.
Integrating:
$\|w(T)\| \le \|w(0)\| + \int_0^T C_w(1 + \|v(s)\|)\,ds
\le \|w(0)\| + C_w T + C_w T(\|v(0)\| + C_b T)$.
Squaring and taking expectations,
$\E[\|W_T\|^2]$ is controlled by $\E[\|\theta_0\|^2]$
(and hence by $\E[\|\theta_0\|^4]^{1/2}$), which is finite.
\end{proof}

%----------------------------------------------------------------------
\subsection{\ref{ass:subgauss}: Sub-Gaussian input weights}\label{app:a7_validation}
%----------------------------------------------------------------------

Assumption~\ref{ass:subgauss} requires that the $W$-marginal of
$\mathfrak{m}_*$ be sub-Gaussian.
Unlike the moment conditions (for which
Proposition~\ref{prop:a6_auto} gives a sufficient finite-time
criterion),
sub-Gaussianity cannot be shown to hold automatically:
preserving sub-Gaussianity under the
gradient flow requires the flow map
$\theta_0 \mapsto \theta_T$ to be \emph{globally} Lipschitz
(so that the Gaussian concentration of the initialisation is inherited
by the push-forward), but the mixture-model flow map is only \emph{locally}
Lipschitz---the drift $b_w$ depends linearly on $\|v\|$
(Proposition~\ref{prop:a6_auto}, Part~(iii)),
making the Jacobian of the flow grow with $\|v_0\|$.

Empirically, however, the trained hidden-weight marginals remain
compatible with sub-Gaussian tails across the mean-field experiments.
Figure~\ref{fig:a7_kappa2} plots the empirical sub-Gaussian parameter
$\hat\kappa^2 \coloneqq \max_\lambda 2\log M(\lambda)/\lambda^2$
(estimated via the moment-generating function on random projections)
across widths.
The values remain moderate ($\hat\kappa^2 \le 1.5$) and do not grow
with width.
Figure~\ref{fig:a7_hist} compares the standardised weight histograms
against the $\cN(0,1)$ reference for an easy case
(\texttt{synthetic}, $K=1$) and the hardest case
(\texttt{synthetic\_transpose}, $K=5$).
In the latter, the post-training law deviates substantially from a
Gaussian shape (reflecting component specialisation in the
multimodal target), yet the tails remain no heavier than those of a
sub-Gaussian distribution with a moderately inflated constant.

\begin{figure}[H]
\centering
\includegraphics[width=0.65\textwidth]{a7/a7_kappa2.png}
\caption{Empirical sub-Gaussian parameter $\hat\kappa^2$ across widths.
  The dashed line marks the Gaussian baseline ($\hat\kappa^2 = 1$).
  Values remain bounded across all settings, supporting
  Assumption~\ref{ass:subgauss} in practice.}
\label{fig:a7_kappa2}
\end{figure}

\begin{figure}[H]
\centering
\includegraphics[width=0.48\textwidth]{a7/a7_hist_synthetic_K1.png}
\hfill
\includegraphics[width=0.48\textwidth]{a7/a7_hist_synthetic_transpose_K5.png}
\caption{Standardised weight histograms vs.\ $\cN(0,1)$ reference.
  Left: \texttt{synthetic}, $K=1$ (easy case; close to Gaussian).
  Right: \texttt{synthetic\_transpose}, $K=5$ (hardest case;
  non-Gaussian shape but compatible sub-Gaussian tails).}
\label{fig:a7_hist}
\end{figure}

%======================================================================
\section{Detailed Numerical Evaluation}\label{app:nonvacuous_detail}
%======================================================================

We evaluate the worst-case bound~\eqref{eq:barrier_form} and
the data-dependent bound~\eqref{eq:dir_bound}
on mean-field MDN models from the assumption-validation experiments
(Appendix~\ref{app:assumption_validation}).

\paragraph{Setup.}
Two-layer MDNs with sigmoid activation
($S_\sigma = 1$, $L_\sigma = 1/4$, $d = 1$)
and $1/N$ mean-field scaling
are trained on the \texttt{synthetic} and \texttt{synthetic\_transpose}
datasets for $10{,}000$ epochs of full-batch gradient descent
(lr $= 10^{-2}$),
with $K \in \{1, 5\}$ and $N \in \{1024, 2048, 4096, 8192, 16384\}$.
For each $(K, N)$ combination, 10 seeds are trained and
5 seed pairs $(0\text{--}1, 2\text{--}3, \ldots, 8\text{--}9)$
are aligned via $W$-OT (hidden) $+$ Hellinger (component).
All quantities below are averaged over seed pairs.

%----------------------------------------------------------------------
\paragraph{Measured parameters ($K=1$, \texttt{synthetic}, $N=1024$).}
$s_{\min} = 0.062$,
$M_V^\infty = 13.4$,
$m_{2,*} = 260$,
$Y_{\max} = 1.43$,
$\kappa_W = 19.4$,
$\lambda_{\max}(\Sigma_x) = 0.36$.
These give $R = S_\sigma M_V^\infty = 13.4$,
$\Lambda = Y_{\max} + R = 14.8$.

\emph{Worst-case bound~\eqref{eq:barrier_form}.}
Substituting into Proposition~\ref{prop:loss_regularity}
(with $K=1$, so $\rho_{K=1} \le 1/s_{\min}^2$):
\begin{align*}
  C_L^\infty
  &\le 3 + \Lambda/s_{\min}^2 + \Lambda^2/s_{\min}^2
  \approx 6.1 \times 10^4, \\
  \rho
  &\le 1/s_{\min}^2 \approx 261.
\end{align*}
The empirical mismatches are
$\tilde B_N = 5.21$ and $\Delta_N = 6.4 \times 10^{-4}$.
The worst-case bound is
$\cB_{\mathrm{wc}} \approx 2.3 \times 10^5$,
entirely dominated by the transfer term
($\rho\,\Delta_N/8 \approx 0.02$ is negligible),
versus the empirical barrier $\cB_{\mathrm{emp}} \approx 0.11$.

\emph{Data-dependent bound~\eqref{eq:dir_bound}.}
The data-dependent constants measured on the trained model pair are:
\[
  \bar C_L \approx 12
  \quad(\text{vs.}\ C_L^\infty \approx 6.1 \times 10^4,
  \;\text{a}\;\sim\!5000\times\;\text{reduction}),
  \qquad
  \E_P[\rho_{\mathrm{dir}}] \approx 3 \times 10^{-5}.
\]
The transfer contribution is
$\tfrac{1}{2}\bar C_L L_\sigma M_V^\infty\sqrt{\tilde B_N}
\approx \tfrac{1}{2} \times 12 \times 0.25 \times 13.4 \times 2.3
\approx 45$.
The semi-convexity contribution is
$\tfrac{1}{8}\E_P[\rho_{\mathrm{dir}}] < 10^{-5}$---negligible.
The combined data-dependent bound is
$\cB_{\mathrm{dd}} \approx 45$,
a $5000\times$ improvement over the worst-case bound,
though still $\sim 400\times$ the empirical barrier.

%----------------------------------------------------------------------
\paragraph{Width-scaling behaviour.}
Table~\ref{tab:bound_scaling} reports the bound components
and empirical barrier across widths.

\begin{table}[H]
\centering
\small
\renewcommand{\arraystretch}{1.15}
\caption{Width scaling of barrier bounds (5-pair average,
  \texttt{synthetic} dataset).
  $\cB_{\mathrm{dd}}^{\mathrm{tr}}$ and
  $\cB_{\mathrm{dd}}^{\mathrm{sc}}$ denote the transfer
  and semi-convexity terms of~\eqref{eq:dir_bound}.}
\label{tab:bound_scaling}
\begin{tabular}{c|r|rrr|rr}
$K$ & $N$ & $\cB_{\mathrm{emp}}$ & $\cB_{\mathrm{dd}}$
  & $\cB_{\mathrm{wc}}$
  & $\tilde B_N$ & $\Delta_N$ \\
\hline
\multirow{5}{*}{1}
  & 1024  & 0.105 & $4.5 \times 10^{1}$ & $2.3 \times 10^{5}$
    & 5.21 & $6.4 \times 10^{-4}$ \\
  & 2048  & 0.060 & $3.6 \times 10^{1}$ & $1.8 \times 10^{5}$
    & 3.37 & $4.2 \times 10^{-4}$ \\
  & 4096  & 0.026 & $2.8 \times 10^{1}$ & $1.4 \times 10^{5}$
    & 2.11 & $1.3 \times 10^{-4}$ \\
  & 8192  & 0.008 & $2.2 \times 10^{1}$ & $1.1 \times 10^{5}$
    & 1.31 & $1.7 \times 10^{-4}$ \\
  & 16384 & 0.004 & $1.7 \times 10^{1}$ & $8.7 \times 10^{4}$
    & 0.80 & $3.7 \times 10^{-5}$ \\
\hline
\multirow{5}{*}{5}
  & 1024  & 0.73  & $4.4 \times 10^{3}$ & $7.4 \times 10^{21}$
    & 18.2 & 154 \\
  & 2048  & 0.60  & $2.2 \times 10^{3}$ & $2.9 \times 10^{24}$
    & 24.4 & 180 \\
  & 4096  & 0.97  & $3.5 \times 10^{3}$ & $3.5 \times 10^{23}$
    & 11.8 & 152 \\
  & 8192  & 0.48  & $2.3 \times 10^{3}$ & $3.3 \times 10^{22}$
    & 23.4 & 76 \\
  & 16384 & 0.25  & $2.2 \times 10^{3}$ & $7.2 \times 10^{22}$
    & 14.5 & 123 \\
\end{tabular}
\end{table}

\begin{figure}[H]
\centering
\includegraphics[width=\textwidth]{barrier_scaling_full_batch.pdf}
\caption{Barrier vs.\ width (log--log) for both datasets and $K \in \{1,5\}$.
  Blue: empirical barrier after alignment;
  orange: data-dependent bound~\eqref{eq:dir_bound};
  green: empirical barrier without alignment.
  Error bars show standard deviation over 5 seed pairs.
  Dashed/dotted grey lines indicate $O(1/N)$ and $O(1/\sqrt{N})$ reference slopes.
  For $K=1$, $\cB_{\mathrm{emp}}$ decays as $\sim N^{-1.3}$
  (faster than $O(1/N)$), while $\cB_{\mathrm{dd}}$ follows
  $\sim N^{-0.34}$ ($\approx O(1/\sqrt{N})$).
  For $K=5$, neither quantity shows clear decay over this width range.}
\label{fig:barrier_scaling}
\end{figure}

\emph{$K=1$ scaling.}
The empirical barrier decays as $\cB_{\mathrm{emp}} \propto N^{-1.3}$,
faster than the theoretical $O(1/N)$.
The data-dependent bound decays as
$\cB_{\mathrm{dd}} \propto N^{-0.34}$,
consistent with the $O(\sqrt{\tilde B_N})$ rate since
$\tilde B_N \propto N^{-0.68}$.
The bound is entirely transfer-dominated
($\cB_{\mathrm{dd}}^{\mathrm{sc}} < 10^{-5}$ across all widths)
because $\Delta_N < 10^{-3}$.
The gap $\cB_{\mathrm{dd}}/\cB_{\mathrm{emp}}$ grows from
$\sim\!400\times$ to $\sim\!5000\times$ as $N$ increases,
reflecting the different scaling exponents.

\emph{$K=5$ scaling.}
For $K=5$, neither $\tilde B_N$ nor $\Delta_N$ exhibits
systematic $O(1/N)$ decay over the tested width range:
$\tilde B_N$ varies between 12--24 and $\Delta_N$ between 76--180
without a clear downward trend.
This indicates that propagation of chaos has not converged
at these widths for $K=5$,
consistent with the neuron-density overlays in
Figure~\ref{fig:a5_overlay}.
Accordingly, $\cB_{\mathrm{dd}} \approx 2 \times 10^3$
remains roughly flat across widths.
The worst-case bound is astronomically vacuous
($10^{21}$--$10^{24}$), driven by
$s_{\min} \approx 10^{-4}$ pushing
$\rho_{\mathrm{wc}}$ to $10^{20}$--$10^{22}$.
In contrast, $\rho_{\mathrm{emp}} \approx 10^{4}$--$10^{6}$,
a gap of $\sim\!10^{16}$.

%----------------------------------------------------------------------
\paragraph{Transfer-dominated regime.}
A key empirical finding is that for all settings the
data-dependent bound is dominated by the transfer term,
not the semi-convexity term.
Three factors explain why the semi-convexity contribution is small:
\begin{enumerate}[leftmargin=*,itemsep=2pt]
\item For $K=1$, $\Delta_N < 10^{-3}$ because the aligned outputs
  are nearly identical, making $\rho\,\Delta_N/8$ negligible
  regardless of $\rho$.
\item For $K \ge 2$,
  the $z_\pi$-block of $\nabla^2\cL$ is positive semidefinite
  ($C_\pi \succeq 0$),
  so the directional semi-convexity $\rho_{\mathrm{dir}}$
  automatically discards the (potentially large) $\pi$-component
  of $\|\Delta z\|^2$.
\item At well-fitted data points, $|y-\mu_k|/s_k = O(1)$,
  so the local Hessian eigenvalues are far from the $s_{\min}$-dependent
  worst-case bounds.
\end{enumerate}
The bottleneck is thus the transfer term
$\tfrac{1}{2}\bar C_L L_\sigma M_V^\infty \sqrt{\tilde B_N}$,
whose magnitude is driven by $M_V^\infty$---the per-coordinate head
moment, which grows substantially during training
(from $\sim\!1$ at initialisation to $\sim\!13$--$33$ post-training).

\paragraph{Results for \texttt{synthetic\_transpose}.}
The same qualitative pattern holds:
for $K=1$, $\cB_{\mathrm{dd}}$ decays from $181$ to $55$
($\propto N^{-0.45}$)
while $\cB_{\mathrm{emp}}$ decays from $0.056$ to $0.002$
($\propto N^{-1.1}$).
For $K=5$, $\cB_{\mathrm{dd}} \approx 5 \times 10^3$--$10^4$
remains flat, with $\cB_{\mathrm{emp}} \approx 2.4$--$3.6$
also showing no clear decay, again indicating
the mean-field regime has not been reached for $K=5$
at these widths.

%======================================================================
\subsection{vMF Numerical Evaluation}\label{app:vmf_numerical}
%======================================================================

We evaluate the vMF barrier bounds
(Corollary~\ref{cor:vmf_barrier}) on synthetic directional data on
$S^1$ ($d_y = 2$), with $K \in \{1, 3\}$ and
$N \in \{512, 1024, 2048, 4096\}$
using the \texttt{rotating} dataset
(Section~\ref{sec:vmf_experiments}).

\paragraph{Measured parameters ($K=1$, $N=1024$).}
$M_V^\infty = 65.6$,
$\bar C_L = 0.40$
(vs.\ $C_{L,\vmf}^\infty \le 4.83$, a $12\times$ margin),
$\rho = 6$ (universal).
The empirical mismatches are
$\tilde B_N = 5.4$ and $\Delta_N = 0.003$.
The worst-case bound is
$\cB_{\mathrm{wc}} \approx 257$
and the data-dependent bound is
$\cB_{\mathrm{dd}} \approx 21$,
versus the empirical barrier $\cB_{\mathrm{emp}} \approx 0.013$.
These ratios ($\sim\!20{,}000\times$ and $\sim\!1{,}600\times$
respectively) are dramatically better than the MDN's $K=1$
worst-case ratio of $\sim\!2 \times 10^6$.

\paragraph{Width-scaling behaviour.}
Table~\ref{tab:vmf_bound_scaling} reports the vMF bounds across widths.

\begin{table}[H]
\centering
\small
\renewcommand{\arraystretch}{1.15}
\caption{Width scaling of vMF barrier bounds (5-pair average,
  \texttt{rotating} dataset, $d_y = 2$).}
\label{tab:vmf_bound_scaling}
\begin{tabular}{c|r|rrrr}
$K$ & $N$ & $\cB_{\mathrm{emp}}$ & $\cB_{\mathrm{dd}}$
  & $\cB_{\mathrm{wc}}$
  & $\cB_{\mathrm{dd}}/\cB_{\mathrm{emp}}$ \\
\hline
\multirow{4}{*}{1}
  & 512   & 0.019 & 24.4 & 294
    & 1300 \\
  & 1024  & 0.013 & 21.4 & 257
    & 1600 \\
  & 2048  & 0.006 & 15.8 & 191
    & 2500 \\
  & 4096  & 0.004 & 12.0 & 144
    & 3100 \\
\hline
\multirow{4}{*}{3}
  & 512   & 1.02  & $1.8 \times 10^{3}$ & $3.8 \times 10^{3}$
    & 1800 \\
  & 1024  & 1.16  & $1.5 \times 10^{3}$ & $3.6 \times 10^{3}$
    & 1300 \\
  & 2048  & 0.83  & $1.9 \times 10^{3}$ & $3.7 \times 10^{3}$
    & 2300 \\
  & 4096  & 1.23  & $1.9 \times 10^{3}$ & $3.9 \times 10^{3}$
    & 1600 \\
\end{tabular}
\end{table}

\emph{$K=1$ scaling.}
The empirical barrier decays from $0.019$ to $0.004$
as $N$ grows from $512$ to $4096$, consistent with the predicted
$O(1/N)$ rate.
Both $\cB_{\mathrm{wc}}$ and $\cB_{\mathrm{dd}}$ also decay
(by roughly $2\times$), confirming that the bounds track the empirical
trend.
The bound-to-empirical gap increases from $\sim\!1300\times$ to
$\sim\!3100\times$, mirroring the MDN observation that the empirical
barrier decays faster than the bounds.

\emph{$K=3$ scaling.}
The empirical barrier is approximately flat around $1.0$,
while the bounds are flat at $\sim\!2{,}000$--$4{,}000$.
The output-space mismatch $\Delta_N$ does not converge at these
widths, indicating that the component alignment quality is the
bottleneck.
Unlike the MDN $K=5$ case, however, the vMF worst-case bound
($\sim\!4 \times 10^3$) is only one order of magnitude above the
data-dependent bound, whereas the MDN gap is $10^{18}$.

%======================================================================

\begin{thebibliography}{46}

\bibitem{ferbach2024}
D.~Ferbach, B.~Goujaud, G.~Gidel, and A.~Dieuleveut,
``Proving linear mode connectivity of neural networks via optimal
transport,'' in \emph{Proceedings of AISTATS}, 2024.
arXiv:2310.19103.

\bibitem{bishop1994}
C.~M.~Bishop,
``Mixture density networks,''
Tech.\ Rep.\ NCRG/94/004, Aston University, 1994.

\bibitem{frankle2020}
J.~Frankle, G.~K.~Dziugaite, D.~Roy, and M.~Carlin,
``Linear mode connectivity and the lottery ticket hypothesis,''
in \emph{Proceedings of ICML}, 2020.

\bibitem{entezari2022}
R.~Entezari, S.~A.~Sedghi, O.~Saukh, and B.~Neyshabur,
``The role of permutation invariance in linear mode connectivity of
neural networks,'' in \emph{Proceedings of ICLR}, 2022.

\bibitem{ainsworth2023}
S.~K.~Ainsworth, J.~Hayase, and S.~Srinivasa,
``Git re-basin: Merging models modulo permutation symmetries,''
in \emph{Proceedings of ICLR}, 2023.

\bibitem{tran2025moe}
V.-H.~Tran, V.~H.~Trinh, K.~V.~Bui, and T.~M.~Nguyen,
``On linear mode connectivity of mixture-of-experts architectures,''
in \emph{Advances in Neural Information Processing Systems}, 2025.
arXiv:2509.11348.

\bibitem{zhan2025permutation}
K.~Zhan, P.~Li, and L.~Wu,
``Analyzing the role of permutation invariance in linear mode
connectivity,'' in \emph{Proceedings of AISTATS}, 2025.
arXiv:2503.06001.

\bibitem{villani2009}
C.~Villani,
\emph{Optimal Transport: Old and New}.
Springer, 2009.

\bibitem{bobkov1999}
S.~G.~Bobkov and F.~G\"otze,
``Exponential integrability and transportation cost related to
logarithmic Sobolev inequalities,''
\emph{Journal of Functional Analysis}, 163(1):1--28, 1999.

\bibitem{ledoux2001}
M.~Ledoux,
\emph{The Concentration of Measure Phenomenon},
Mathematical Surveys and Monographs, vol.~89,
American Mathematical Society, 2001.

\bibitem{weed2019}
J.~Weed and F.~Bach,
``Sharp asymptotic and finite-sample rates of convergence of empirical
measures in Wasserstein distance,''
\emph{Bernoulli}, 25(4A):2620--2648, 2019.
arXiv:1707.00087.

\bibitem{facco2017}
E.~Facco, M.~d'Errico, A.~Rodriguez, and A.~Laio,
``Estimating the intrinsic dimension of datasets by a minimal neighborhood
information,'' \emph{Scientific Reports}, 7:12140, 2017.

\bibitem{rotskoff2022}
G.~M.~Rotskoff and E.~Vanden-Eijnden,
``Trainability and accuracy of artificial neural networks: An interacting
particle system approach,'' \emph{Communications on Pure and Applied
Mathematics}, 75(9):1889--1935, 2022.

\bibitem{garipov2018}
T.~Garipov, P.~Izmailov, D.~Podoprikhin, D.~Vetrov, and A.~G.~Wilson,
``Loss surfaces, mode connectivity, and fast ensembling of {DNN}s,''
in \emph{Advances in Neural Information Processing Systems}, 2018.
arXiv:1802.10026.

\bibitem{draxler2018}
F.~Draxler, K.~Veschgini, M.~Salmhofer, and F.~A.~Hamprecht,
``Essentially no barriers in neural network energy landscape,''
in \emph{Proceedings of ICML}, 2018.
arXiv:1803.00885.

\bibitem{theus2025}
M.~Theus, R.~M\"{u}ller, J.~Vogt, T.~Hofmann, and K.~Roth,
``Generalized linear mode connectivity for transformers,''
in \emph{Advances in Neural Information Processing Systems}, 2025.
arXiv:2506.22712.

\bibitem{wortsman2022soups}
M.~Wortsman, G.~Ilharco, S.~Y.~Gadre, R.~Roelofs, R.~Gontijo-Lopes,
A.~S.~Morcos, H.~Namkoong, H.~Fang, M.~Luca, L.~Schmidt, and
A.~Farhadi,
``Model soups: Averaging weights of multiple fine-tuned models improves
accuracy and robustness,''
in \emph{Proceedings of ICML}, 2022.
arXiv:2203.05482.

\bibitem{ilharco2023task}
G.~Ilharco, M.~T.~Ribeiro, M.~Wortsman, S.~Gururangan, L.~Schmidt,
H.~Hajishirzi, and A.~Farhadi,
``Editing models with task arithmetic,''
in \emph{Proceedings of ICLR}, 2023.
arXiv:2212.04089.

\bibitem{yadav2023ties}
P.~Yadav, D.~Tam, L.~Choshen, C.~Raffel, and M.~Bansal,
``{TIES}-Merging: Resolving interference when merging models,''
in \emph{Advances in Neural Information Processing Systems}, 2023.
arXiv:2306.01708.

\bibitem{jacobs1991}
R.~A.~Jacobs, M.~I.~Jordan, S.~J.~Nowlan, and G.~E.~Hinton,
``Adaptive mixtures of local experts,''
\emph{Neural Computation}, 3(1):79--87, 1991.

\bibitem{shazeer2017}
N.~Shazeer, A.~Mirhoseini, K.~Maziarz, A.~Davis, Q.~Le, G.~Hinton,
and J.~Dean,
``Outrageously large neural networks: The sparsely-gated
mixture-of-experts layer,''
in \emph{Proceedings of ICLR}, 2017.
arXiv:1701.06538.

\bibitem{kanoh2025dte}
R.~Kanoh and M.~Sugiyama,
``Linear mode connectivity in differentiable tree ensembles,''
in \emph{Proceedings of ICLR}, 2025.
arXiv:2405.14596.

\bibitem{graves2013}
A.~Graves,
``Generating sequences with recurrent neural networks,''
arXiv:1308.0850, 2013.

\bibitem{ha2018worldmodels}
D.~Ha and J.~Schmidhuber,
``World models,''
arXiv:1803.10122, 2018.

\bibitem{makansi2019}
O.~Makansi, E.~Ilg, O.~Cicek, and T.~Brox,
``Overcoming limitations of mixture density networks: A sampling and
fitting framework for multimodal future prediction,''
in \emph{Proceedings of CVPR}, 2019.
arXiv:1906.03631.

\bibitem{papamakarios2016}
G.~Papamakarios and I.~Murray,
``Fast $\varepsilon$-free inference of simulation models with
{B}ayesian conditional density estimation,''
in \emph{Advances in Neural Information Processing Systems}, 2016.
arXiv:1605.06376.

\bibitem{mei2018meanfield}
S.~Mei, A.~Montanari, and P.-M.~Nguyen,
``A mean field view of the landscape of two-layer neural networks,''
\emph{Proceedings of the National Academy of Sciences},
115(33):E7665--E7671, 2018.
arXiv:1804.06561.

\bibitem{carmona2018probabilistic}
R.~Carmona and F.~Delarue,
\emph{Probabilistic Theory of Mean Field Games with Applications~I},
Probability Theory and Stochastic Modelling, vol.~83,
Springer, 2018.

\bibitem{boyd2004}
S.~Boyd and L.~Vandenberghe,
\emph{Convex Optimization}.
Cambridge University Press, 2004.

\bibitem{debortoli2020}
V.~De~Bortoli, A.~Durmus, X.~Fontaine, and U.~Simsekli,
``Quantitative propagation of chaos for {SGD} in wide neural networks,''
in \emph{Advances in Neural Information Processing Systems}, 2020.
arXiv:2007.06352.

\bibitem{sirignano2020lln}
J.~Sirignano and K.~Spiliopoulos,
``Mean field analysis of neural networks: A law of large numbers,''
\emph{SIAM Journal on Applied Mathematics}, 80(2):725--752, 2020.
arXiv:1805.01053.

\bibitem{sirignano2020clt}
J.~Sirignano and K.~Spiliopoulos,
``Mean field analysis of neural networks: A central limit theorem,''
\emph{Stochastic Processes and their Applications}, 130(3):1820--1852,
2020.
arXiv:1808.09372.

\bibitem{gess2024smf}
B.~Gess, S.~Kassing, and V.~Konarovskyi,
``Stochastic modified flows, mean-field limits and dynamics of
stochastic gradient descent,''
\emph{Journal of Machine Learning Research}, 25(30):1--55, 2024.
arXiv:2302.07125.

\bibitem{nguyen2023multilayer}
P.-M.~Nguyen and H.~T.~Pham,
``A rigorous framework for the mean field limit of multilayer neural
networks,''
\emph{Mathematical Statistics and Learning}, 6(3):201--357, 2023.
arXiv:2001.11443.

\bibitem{suzuki2023mfld}
T.~Suzuki, D.~Wu, and A.~Nitanda,
``Convergence of mean-field {L}angevin dynamics: Time and space
discretization, stochastic gradient, and variance reduction,''
in \emph{Advances in Neural Information Processing Systems}, 2023.
arXiv:2306.07221.

\bibitem{yakowitz1968}
S.~J.~Yakowitz and J.~D.~Spragins,
``On the identifiability of finite mixtures,''
\emph{The Annals of Mathematical Statistics},
vol.~39, no.~1, pp.~209--214, 1968.

\bibitem{pinkus1999}
A.~Pinkus,
``Approximation theory of the {MLP} model in neural networks,''
\emph{Acta Numerica}, vol.~8, pp.~143--195, 1999.

\bibitem{luce1959}
R.~D.~Luce,
\emph{Individual Choice Behavior: A Theoretical Analysis}.
John Wiley \& Sons, 1959.

\bibitem{lei2020}
J.~Lei,
``Convergence and concentration of empirical measures under
{W}asserstein distance in unbounded functional spaces,''
\emph{Bernoulli}, 26(1):767--798, 2020.

\bibitem{banerjee2005}
A.~Banerjee, I.~S.~Dhillon, J.~Ghosh, and S.~Sra,
``Clustering on the unit hypersphere using von {M}ises-{F}isher
distributions,''
\emph{Journal of Machine Learning Research}, 6:1345--1382, 2005.

\bibitem{gopal2014}
S.~Gopal and Y.~Yang,
``Von {M}ises-{F}isher clustering models,''
in \emph{Proceedings of ICML}, 2014.

\bibitem{mardia2000}
K.~V.~Mardia and P.~E.~Jupp,
\emph{Directional Statistics}.
John Wiley \& Sons, 2000.

\bibitem{prokudin2018deep}
S.~Prokudin, P.~Gehler, and S.~Nowozin,
``Deep directional statistics: Pose estimation with uncertainty
quantification,''
in \emph{Proceedings of ECCV}, 2018.
arXiv:1805.03430.

\bibitem{liu2024probabilistic}
H.~Liu, Y.~Liu, Y.~Yang, and J.-F.~Yang,
``Probabilistic rotation modeling based on directional mixture density
networks,''
\emph{Information Sciences}, 662:120225, 2024.

\bibitem{han2023sphericalmdn}
S.~Han, H.~Li, H.~Shi, and S.~Huang,
``Uncertainty quantification in molecular property prediction through
spherical mixture density networks,''
\emph{Engineering Applications of Artificial Intelligence},
126:106180, 2023.

\bibitem{ghrear2024direction}
M.~Ghrear, P.~Sadowski, and S.~E.~Vahsen,
``Deep probabilistic direction prediction in {3D} with applications to
directional dark matter detectors,''
\emph{Machine Learning: Science and Technology}, 5(2):025048, 2024.
arXiv:2403.15949.

\bibitem{jena_climate}
Max Planck Institute for Biogeochemistry,
``Jena Climate dataset,'' 2016.
Available at \url{https://www.bgc-jena.mpg.de/wetter/}.

\end{thebibliography}

\end{document}
